% concatenated from main.tex
\documentclass[12pt, reqno]{amsart}
%


\author[M.~Caprio and S.~Mukherjee]{Michele Caprio and Sayan Mukherjee}
\address{Department of Computer  Science, University of Manchester, Oxford Road, Manchester M13 9PL, United Kingdom}
\email{michele.caprio@manchester.ac.uk}
\urladdr{\url{https://michelecaprio.wixsite.com/caprio}}   
\address{Max Planck Institute for Mathematics in the Sciences, Inselstraße 22, Leipzig 04103, Germany}
\email{sayan.mukherjee@mis.mpg.de}
\urladdr{\url{https://sayanmuk.github.io/}}    

 
%%
\keywords{Convex geometry; Choquet theory; Finite mixture models; admixture models.}
\subjclass[2010]{Primary: 52B11; Secondary: 60D05, 62G20.}

%40A35  	Ideal and statistical convergence
%40A05  	Convergence and divergence of series and sequences
%37C70  	Attractors and repellers of smooth dynamical systems and their topological structure
%49j99 Calculus of variations and optimal control; optimization
%37C99  	None of the above, but in this section

\title{Finite Admixture Models: a Bridge with Stochastic Geometry and Choquet Theory}


\usepackage[T1]{fontenc}
\usepackage{amsmath}
\usepackage{amssymb}
\usepackage{amsthm}
\usepackage[left=2.5cm, right=2.5cm, top=3cm]{geometry}
\usepackage{hyperref}
\usepackage{algorithm}
\usepackage{algorithmicx}
\usepackage{algpseudocode}
\usepackage{bbm}
\usepackage{tikz}
\usepackage{caption}
\usepackage{subcaption}
\usepackage{fancyhdr}
\addtolength{\topmargin}{-2.0pt}
\usepackage[inline]{enumitem}
\usepackage{comment}
\usepackage{nicefrac}
\usepackage{bm}
\usepackage{mathrsfs}
\usepackage{graphicx}
\usepackage[utf8]{inputenc}
\usepackage{cancel}
\usepackage{mathtools}
%\usepackage{mathabx}

\newcommand{\vertiii}[1]{{\left\vert\kern-0.25ex\left\vert\kern-0.25ex\left\vert #1 
    \right\vert\kern-0.25ex\right\vert\kern-0.25ex\right\vert}}
		
\AtBeginDocument{%
   \def\MR#1{}
}


\algdef{SE}[DOWHILE]{Do}{doWhile}{\algorithmicdo}[1]{\algorithmicwhile\ #1}%

\newtheorem{thm}{Theorem}[section]
\newtheorem{cor}[thm]{Corollary}%[section]
\newtheorem{lem}[thm]{Lemma}
\newtheorem{prop}[thm]{Proposition}
\newtheorem{question}{Question}[section]
\theoremstyle{definition} 
\newtheorem{defi}[thm]{Definition}%[section]
\let\olddefi\defi
\renewcommand{\defi}{\olddefi\normalfont}
\newtheorem{rmk}[thm]{Remark}
\let\oldrmk\rmk
\renewcommand{\rmk}{\oldrmk\normalfont}

\DeclareMathOperator*{\argmax}{arg\,max}
\DeclareMathOperator*{\argmin}{arg\,min}

\newtheorem{theorem}{Theorem}%[section]
\newtheorem{lemma}[theorem]{Lemma}
\newtheorem{proposition}[theorem]{Proposition}
\newtheorem{corollary}{Corollary}[theorem]
\newtheorem{claim}[theorem]{Claim}
\newtheorem{definition}[theorem]{Definition}
\newtheorem{remark}{Remark}
\newtheorem{example}[theorem]{Example}

\pagestyle{fancy}
\fancyhf{}
\fancyhead[CO]
{\textsc{Admixture Models, Stochastic Geometry, and Choquet Theory}}
\fancyhead[CE]
{\textsc{Michele Caprio and Sayan Mukherjee}}
\fancyhead[RO,LE]{\thepage}

\setlength{\headheight}{12pt}

\hypersetup{
    pdftitle={prove},
    pdfauthor={autori},
    pdfmenubar=false,
    pdffitwindow=true,
    pdfstartview=FitH,
    colorlinks=true,
    linkcolor=blue,
    citecolor=green,
    urlcolor=cyan
}

\uchyph=0
\def\lcm{\operatorname{lcm}}

\providecommand{\MR}[1]{}
\providecommand{\bysame}{\leavevmode\hbox to3em{\hrulefill}\thinspace}
\providecommand{\MR}{\relax\ifhmode\unskip\space\fi MR }
\providecommand{\MRhref}[2]{%
  \href{http://www.ams.org/mathscinet-getitem?mr=#1}{#2}
}
\providecommand{\href}[2]{#2}

\newcommand{\textoperatorname}[1]{%
  \operatorname{\textnormal{#1}}%
}

\newcommand\xqed[1]{%
  \leavevmode\unskip\penalty9999 \hbox{}\nobreak\hfill
  \quad\hbox{#1}}
\newcommand\demo{\xqed{$\triangle$}}


\newcommand{\robin}[1]{{\color{red!70!black}{ \bf \sf \scriptsize Robin:} \sf \scriptsize #1}}
                                 
\begin{document}





\begin{abstract}
Given a finite admixture model whose components and weights are unknown, let the number of identifiable components be a function of the amount of data sampled from a known distribution on the unit simplex. We use techniques from stochastic convex geometry to find the growth rate of its expected value. In addition, when the components are known but the weights are not, we provide an application of the classic Glivenko-Cantelli's theorem that allows us to retrieve the Choquet measure supported on the identifiable admixture components. In turn, this gives us the identifiable admixture weights. Finally, we propose a novel algorithm that estimates the model capturing the complexity of the data using only the strictly necessary number of components.
\end{abstract}
%We also show that by placing a Dirichlet process prior on the densities supported on the unit simplex, we are able to retrieve the Dirac measure at the Choquet measure supported on the identifiable components. 
%In Bayesian density estimation, a question of interest is how the number of components in a finite mixture model grows with the number of observations. We provide a novel perspective on this question by using results from stochastic geometry. We find that the growth rate of the expected number of components of a finite mixture model whose components belong to the unit simplex $\Delta^{J-1}$ of Euclidean space $\mathbb{R}^J$ is $(\log n)^{J-1}$. We also provide a central limit theorem for the number of components. In addition, we relate our model to a classical non-parametric density estimator based on a P\'olya tree. Combining the latter with techniques from Choquet theory, we are able to retrieve mixture weights. We also give the rate of convergence of the P\'olya tree posterior to the Dirac measure on the weights. We further present an algorithm to correctly specify the number of components in a latent Dirichlet allocation (LDA) analysis.

\maketitle
\thispagestyle{empty}

\section{Introduction}\label{intro}
Finite mixture models go back at least to \cite{pearson2,pearson1} and have served as a workhorse in stochastic modeling \cite{everitt,lindsay,mengersen}. Applications include clustering \cite{mclachlan1}, hierarchical or latent space models \cite{masyn}, and semiparametric models \cite{mcnicholas} where a mixture of simple distributions is used to model data that is putatively generated from a complex distribution. In finite mixture models, the mixing distribution is over a finite number of components; there are also many examples of infinite mixture models in the Bayesian nonparametrics literature \cite{antoniak1974,west94}. 


We consider a finite mixture of multinomials. We start with the basic multinomial model where an observation $X$ takes on $J$ possible values $\{1,\ldots,J\}$, and
$X \sim \mbox{Mult}(\pi)$, with $\pi \equiv (\pi_1,\ldots,\pi_J)^\top$ where $\pi_j=\mathbb{P}(X= j)$, with $\pi_j \geq 0$ for all $j$ and $\sum_{j=1}^J \pi_j=1$. Notice that here and in the remainder of the paper, we adopt the notational convention that, when not specified, the number of trials for a Multinomial distribution is equal to $1$. A mixture of $L$ multinomials can be specified as follows
\begin{equation}\label{fam_first}
    X \sim \mbox{Mult}(\pi), \quad \pi = \sum_{\ell=1}^L \phi_{\ell} f_\ell,
\end{equation}
where probability vector $\phi \equiv (\phi_{1},\ldots,\phi_{L})^\top$ assigns the probability of the  observation $X$ coming from the $\ell$-th mixture component with multinomial parameter 
$$f_\ell=(f_{\ell,1},\ldots,f_{\ell,J})^\top.$$
We have that $\sum_{j=1}^J f_{\ell,j} =1$ with $f_{\ell,j} \geq 0$, and $\sum_{\ell=1}^L \phi_{\ell} =1$ with $\phi_{\ell} \geq 0$, for all $\ell \in \{1,\ldots,L\}$. An important point throughout the paper is that $\pi$ belongs to the convex hull of probability vectors  $\{f_1,\ldots,f_L\}$. The convex hull of $\{f_1,\ldots,f_L\}$ is a function of the {\em identifiable elements} of $\{f_1,\ldots,f_L\}$, that is, {\em those elements that cannot be written as a convex combination of the other $f_\ell$'s}. In other words, we call identifiable the {\em extremal elements} of the convex hull. Hence, understanding the identifiable elements provides information about the  key model parameter $\pi$.
 
%In a Bayesian model we are interested in the posterior $\mathbb{P}(\theta \mid x_1,\ldots,x_n)$ where $\theta$ consists of the set $\{\phi_1,...,\phi_n\}$ and $\{f_1,...,f_L\}$, all of which are probability vectors. One can obtain point estimates of the parameters using an EM algorithm or the posterior using MCMC procedures; there are also variational approaches to compute the posterior.

The finite mixture model we stated is an example of a finite admixture model; the most popular finite admixture model is the latent Dirichlet allocation (LDA) model \cite{blei, Pritchard945}. A classic application of an admixture model is a generative process for documents. Consider a document as a collection of words; LDA posits  that each document is a mixture of a small number of topics, and that these latter can be modeled by a multinomial distribution on the presence of a word in the topic. The hierarchical Dirichlet process \cite{teh}, and generalizations thereof, may be considered as the natural nonparametric counterpart of the LDA model.




%The data we collect are distributed according to the likelihood $\mathscr{L}(\mathbf{x},\theta) \equiv p(x_1,\ldots,x_n \mid \theta)$. We can then compute the posterior $p \equiv p(\theta \mid x_1,\ldots,x_n)=({p}_1,\ldots,{p}_J)^T$, where ${p}_j \geq 0$ for all $j$, $\sum_{j=1}^J {p}_j=1$, and ${p}=\sum_{i=1}^L \tilde{\phi}_i \tilde{f}_i$. We have that $\tilde{\phi}_i \geq 0$ for all $i$, and $\sum_{i=1}^L \tilde{\phi}_i=1$, $\tilde{f}_i=(\tilde{f}_{i,1},\ldots,\tilde{f}_{i,J})^T$ for all $i \in \{1,\ldots,L\}$ such that $\sum_{j=1}^J \tilde{f}_{i,j}=1$, for all $i$, and $\tilde{f}_{i,j} \geq 0$, for all $i$, for all $j$. The $\tilde{\phi}_i$'s and the $\tilde{f}_i$'s are updated from the $\phi_i$'s and the $f_i$'s, respectively, using e.g. the EM algorithm. In the remainder of the paper, we will avoid using tildes when no confusion arises. This particular finite mixture model is an admixture model; the most popular admixture model is the latent Dirichlet allocation model \cite{blei, Pritchard945}. It allows sets of observations to be explained by unobserved groups that explain why some parts of the data are similar. For example, if observations are words collected into documents, it posits that each document is a mixture of a small number of topics and that each word's presence is attributable to one of the document's topics. The hierarchical Dirichlet process \cite{teh}, and generalizations thereof, may be
%considered as the natural nonparametric counterpart of the latent Dirichlet allocation model.


%the following mixture model. Let $\theta$, the parameter of interest, be modeled as a discrete random variable having $J$ possible values, $\theta_1,\ldots,\theta_J$. Let $\pi \equiv \pi(\theta)=(\pi_1,\ldots,\pi_J)^T$ be the vector whose entries are such that $\pi_j=\mathbb{P}(\theta=\theta_j)$, for all $j \in \{1,\ldots,J\}$. Clearly, $\pi_j \geq 0$ for all $j$, and $\sum_{j=1}^J \pi_j=1$. Suppose then that $\pi=\sum_{i=1}^L \phi_i f_i$, where $\sum_{i=1}^L \phi_i=1$, $\phi_i \geq 0$ for all $i$, and $f_i=(f_{i,1},\ldots,f_{i,J})^T$ for all $i \in \{1,\ldots,L\}$ such that $\sum_{j=1}^J f_{i,j}=1$, for all $i$, and $f_{i,j} \geq 0$, for all $i$, for all $j$. That is, the $f_i$'s are frequency vectors.

The $f_\ell$'s and $\pi$ in \eqref{fam_first} are all elements of $\Delta^{J-1}$, the unit simplex on $\mathbb{R}^J$. Again, $\pi$ belongs to the convex hull of  $\{f_\ell\}_{\ell=1}^L$, or $\pi \in \text{Conv}(f_1,\ldots,f_L)$. Hence, an element of a convex hull in the Euclidean unit simplex represents (the distribution of) a finite admixture model. 

Notice that the number of extremal elements of $\text{Conv}(f_1,\ldots,f_L)$, which we denote as $M$, will probably be less than $L$ because some of the components $f_\ell$ are likely to be a convex combination of the others. A key concept in this paper is what we call the \textit{richest cheap model} representing $\pi$, that is, the finite admixture model representing $\pi$ whose components are $\{f_k\}_{k \in \mathcal{I}}$ such that $f_k \not\in \text{Conv}(f_{\mathcal{I}\setminus \{k\}})$, for all $k \in \mathcal{I}$, $\mathcal{I} \subset \{1,\ldots,L\}$, and $\#\mathcal{I}=M$, where $\#$ denotes the cardinality operator. These conditions tell us that the $M$ components of the richest cheap model are a subset of $\{f_1,\ldots,f_L\}$ and cannot be written as a convex combination of one another. By assuming -- without loss of generality -- that the identifiable elements in $\{f_1,\ldots,f_L\}$ are the first $M$ ones, for the richest cheap model we can rewrite $\pi$ in \eqref{fam_first} as
%we can write the richest cheap model as
\begin{equation}\label{rcm}
    \pi=\sum_{\ell=1}^M \varphi_{\ell} f_\ell,
\end{equation}
where we denote by $\varphi \equiv (\varphi_{1},\ldots,\varphi_{M})^\top$ the probability vector that assigns the probability of the observation $X$ coming from the $\ell$-th identifiable component with multinomial parameter $f_\ell=(f_{\ell,1},\ldots,f_{\ell,J})^\top$, $\ell\in\{1,\ldots,M\}$. That is, the $\varphi_{\ell}$'s are the identifiable admixture weights. Of course the $\varphi_{\ell}$'s are such that $\sum_{\ell=1}^M \varphi_{\ell}=1$, and $\varphi_{\ell} \geq 0$, for all $\ell \in \{1,\ldots,L\}$. As we can see, the richest cheap model captures the underlying complexity associated with the data at hand, using only the strictly necessary number of components.




\iffalse
In general, a finite mixture distribution of $m$ components for a random vector $Y$ is given by
$$Y \sim \sum _{k=1}^m p_k f(y; \theta_k), \quad \sum_{k=1}^m p_k =1, \quad p_k \geq 0,$$
where the elements of the probability vector $p=(p_1,...,p_m)^\top$ are mixture weights and $\theta_k$ denotes the parameter values for the $k$-th component. 

\iffalse
This model can be generalized to 
an admixture or graded membership model \cite{Pritchard945,blei} where each element $y_j$ in the $J$-dimensional random vector $Y$ is given by
$$y_{ij} \sim \sum_{k=1}^m p_{kj} f(y_j; \theta_k), \quad \sum_{k=1}^m p_{kj} =1, \quad p_{kj} \geq 0,$$
where the mixture weights $p_k$ are replaced by a mixture probability vector $p_k = \{p_{k1},...,p_{kJ}\}$ that encodes the probability that the $j$-th element of $Y$ is drawn from 
the $k$-th mixture component. This model is called an admixture model because the vector $y$ can contain elements drawn from different components.  
\fi
%where the weights $p_j$'s are nonnegative and sum to $1$, and the $f_j$'s are the component densities. Often times the component densities are known up to a vector $\theta_j$ of parameters; in this case, we write $f_j(y,\theta_j)$. Typically, in applied works the component densities $f_j(y,\theta_j)$ are assumed to belong to the same parametric family, e.g. the multivariate normal. 

Inference on the number of mixture components for finite mixture models can be difficult. In the Bayesian setting one can place a prior on the number of mixture components \cite{miller}. In \cite{guha}, the authors study the consistency of the posterior distribution of the number of
clusters; they also propose a merge-truncate-merge procedure to consistently estimate the number of
clusters from Dirichlet process mixture models. 
%This procedure can also be adapted to the admixture model settings.
In \cite{fuquene}, the authors introduce the use of non-local priors for choosing the number of components in finite mixture models.


Another approach to inference on the number of components is to test whether
the number of components is a given $k$ or $k^\prime>k$. Such literature is quite rich: classical results are summarized in \cite{titterington}, while more modern works include \cite{henna} and \cite{chen,li2}. In the former, an estimator for the number of components is provided based on transformations of the observed data. The latter two propose an EM test for testing whether the number of true components in the mixture is some $k_0>0$, or is larger than $k_0$.

A further recent work of interest is \cite{ohn}, where the authors use a data dependent prior and achieve optimal estimation of mixing measures, as well as posterior consistency for the number of clusters. They also consider a Dirichlet Process mixture to estimate a finite mixture model and show that the number of clusters can be used for consistent estimation on the number of components.
\fi


Rather than developing new tools for working with or applying finite admixture models, the main goal of this paper is to establish connections between finite admixture models, Choquet theory, and stochastic convex geometry. This paper is an extension of \cite[Chapter 2]{thesis}.
%We do so to show how studying the geometry of FMMs allows to find novel and interesting results.


%in the hope that they will shed light on the workings and properties of finite mixture models. 

\subsection{Choquet theory}
Choquet theory, named after mathematician Gustave Choquet, is an area of functional and convex analyses concerned with measures which have support on the extreme points of a convex set \cite{phelps1}. Its fundamental tenet is that we can represent every element in a convex set $C$ via a weighted average of the extremal elements of the set. Here, weighted average is to be understood as a generalization of the usual notion of convex combination to an integral taken over the set $E$ of extreme points of $C$. The formal, central result to Choquet theory is the following.

\begin{theorem}{\textbf{(Choquet, cf. \cite[Theorem, page 14]{phelps1}).}}\label{choq1}
Let $C$ be a metrizable compact convex subset of a locally convex space $V$. Pick any $c \in C$. Then, there exists a probability measure $\nu$ on $C$ which represents $c$ and is supported on $E$, that is,
$$f(c)=\int_E f(e) \nu(\text{d}e),$$
for any affine function $f$ on $C$.
\end{theorem}

Choquet also characterized those compact convex sets $C$ with the property that for every $c \in C$ there is a \textit{unique} probability measure $\nu_c$ supported on $E$ that represents $c$. 
%The necessary and sufficient condition is based on the concept of Choquet simplex. Its general definition is given in Appendix \ref{app_choquet}. Here, it is enough to point out that a nonempty convex set $C$ (not necessarily compact) of a finite-dimensional locally convex space $V$ is a \textit{Choquet simplex} if it is an ordinary simplex with number of vertices equal to $\text{dim}(V)+1$.\footnote{Here dim$(V)$ denotes the dimension of space $V$.} The characterization of $C$, then, is the following.



%In the case when $V$ is finite-dimensional, a Choquet simplex is 

\begin{theorem}{\textbf{(Choquet, cf. \cite[Theorem, page 60]{phelps1}).}}\label{choq2}
Let $C$ be a metrizable closed convex subset of a locally convex space $V$. Then, $C$ is a simplex\footnote{Choquet's result refers to a more general notion than the classical simplex, called {\em Choquet simplex}. Since a classical simplex is a Choquet simplex, and since in this paper we only work with the former, we avoid introducing the notion of Choquet simplex to refrain from unnecessary complications. The interested reader can find it in \cite[Chapter 10, page 52]{phelps}.} if and only if for every $c$ in $C$, there exists a unique measure $\nu_c$ which represents $c$ and is supported on $E$, that is, 
$$f(c)=\int_E f(e) \nu_c(\text{d}e),$$
for any affine function $f$ on $C$.
\end{theorem}

We call $\nu_c$ the \textit{Choquet measure} for $c$. Notice that we differentiate between $\nu$ in Theorem \ref{choq1} and $\nu_c$ in Theorem \ref{choq2} to highlight the fact that the latter uniquely represents $c$. These results entail that studying the extremal elements of a convex set gives us important results concerning the (elements of the) whole set. Choquet theory in the context of finite mixture models has been inspected in \cite{hoff}. There, the author develops an approach that uses Choquet's theorems for inference with the goal of estimating probability measures constrained to lie in a convex set, for example mixture models. The key observation in \cite{hoff} is that  inference over a convex set of measures can be made via unconstrained inference over the set of extreme measures. The main difference between this work and the approach developed in \cite{hoff} is that we consider a convex hull of points in a unit simplex rather than the convex hull of probability measures. Furthermore, our goal is different: we use a result from Choquet theory to retrieve the identifiable weights in the finite admixture model at hand. Notice also that de Finetti's theorem \cite{definetti3,definetti1,definetti2} can be given a geometric interpretation -- inspected in Appendix \ref{infiniteex} -- that is heuristically similar to that of Choquet theory.
%In this paper we bridge the study of finite mixture models with stochastic geometry and Choquet theory. The most important result is Theorem \ref{thm1}. We show that if the geometric representation of a finite mixture model is a simplex, and if we place a Dirichlet process prior on the densities supported on the extrema of such simplex, then the Dirichlet process posterior converges weakly to the Dirac at the Choquet measure $\delta_{\mu_c}$. The Choquet measure retrieves the mixture weights of our model; this result is based on Theorem \ref{choq2}. We also give the rate of convergence of the Dirichlet process posterior to $\delta_{\mu_c}$.


\subsection{Stochastic convex geometry}
This paper also establishes a bridge between finite admixture models and stochastic convex geometry that allows to view finite admixture models as well-studied geometric objects. This insight allows to closely relate the number of identifiable admixture components to the number of extremal elements of a convex body. Thereby, it facilitates studying the asymptotic growth rate and the asymptotic distribution of the number of components. 
%This chapter bridges finite mixture models and (geometric) Choquet theory as well: we give a result to retrieve the weights in a finite mixture model using a uniqueness result by Gustave Choquet (Theorem \ref{choq2}) coupled with a Dirichlet process distribution.
%in a way that makes easy to compute the posterior weights in the finite mixture model.
%(stochastic)

The geometry of finite mixture models has primarily been studied in two contexts: differential geometry \cite{amari,kass_vos} and convex geometry \cite{lindsay,marriott}. The approach in this paper is based on (stochastic) convex geometry. The first to study the geometry of mixture models was Lindsay \cite{lindsay_classic1,lindsay_classic2}. In the first paper, the author  established the geometric properties of the likelihood set and used these properties to study the uniqueness of the maximum likelihood estimator (MLE) as well as other fundamental properties of the MLE. In the second paper, the author established the results for the nonparametric MLE. Lindsay also wrote a book \cite{lindsay}  
%in which he observed how a mixture model can be seen as an element of the unit simplex in some Euclidean space $\mathbb{R}^J$. The 
whose focus was the identifiability of the mixture weights, a Carath\'eodory representation theorem for multinomial mixtures, and the asymptotic mixture geometry. In a recent paper \cite{nguyen}, the author studies the asymptotic behavior of the convex polytope representing a finite admixture model. Their results well complement the ones in section \ref{mainresults}, where the focus is the identifiable admixture components, whose geometric representation is given by the extremal elements of the convex polytope representing the finite admixture model. In \cite{marriott}, the author bridges the differential and convex geometric approaches to identify restrictions for which a mixture model can be written as a tractable geometric quantity that can simplify inference problems. This paper is similar in spirit to Lindsay's work, but uses more modern techniques.






\iffalse
An important concept developed in \cite{hoff} is the relation between infinite exchangeability and Choquet theory. There is a natural relationship between Choquet theory and Bayesian inference grounded by infinite exchangeability and de Finetti's theorem, very nicely outlined in \cite{aldous_2010}. While this relationship is not the focus of our paper, we will discuss the implications of our results with respect to infinite exchangeability in Appendix \ref{infiniteex}.
\fi




%the number of components required by our mixture model may be 
%,  where $M$ is the number of extrema of $\text{Conv}(f_1,\ldots,f_L)$. The number of extremal points $M$ may be smaller than the number of mixture components in a general model as  In section \ref{mainresults}, we derive results 

\subsection{Main results and structure of the paper}
We provide three main results. The first two, Theorem \ref{growth} and Theorem \ref{extrema2}, state the following. Suppose we do not know what the components and the weights in our admixture model are, and we also do not know the number of components. Then, if we assume that the number of identifiable components $M$ is a function $M(n)$ of the amount $n$ of data that we sample from a known distribution on $\Delta^{J-1}$, we are able to tell the speed at which its expected value grows. The other main result is Theorem \ref{thm1}. It states that if we know the number of identifiable components of the model, but not the components themselves nor the weights, we can apply the classic Glivenko-Cantelli's theorem to retrieve the Choquet measure on the components. In turn, the latter gives us the identifiable admixture weights. 
%we can place a Dirichlet process on the densities supported on $\Delta^{J-1}$, which eventually retrieves the identifiable admixture weights. 
We also show how looking for the richest cheap model can be seen as an optimization problem, and we propose an algorithm to solve it. 

The paper is organized as follows. In section \ref{mainresults}, we let the number of identifiable admixture components $M$ depend on the sample size $n$, that is, we let $M=M(n)$. 
%We use ideas from stochastic convex geometry to state properties of a finite admixture model. This constitutes a novelty with respect to the standard approach in Bayesian analysis. These techniques have the potential to uncover other properties of finite mixture models that might otherwise be inaccessible. 
%This can be interpreted as proposing a prior on the number of identifiable mixture components that is a function of the sample size. 
We state and prove Theorems \ref{growth} and \ref{extrema2}. In addition, in Theorems \ref{clt_2} and \ref{clt_gen} we state a central limit theorem (CLT) for the distribution of the number of identifiable admixture components, and in Theorem \ref{conc_unif_thm} we prove that the number of identifiable admixture components concentrates around its expected value. 


%MAYBE RE-ENTERED IN FUTURE: We study the behavior of $M(n)$ as the number of observations increases. In Corollary \ref{growth} we show that if $M(n)$ is given by the cardinality of the extremal set of the convex hull of $n$ elements sampled iid from the uniform over the simplex $\Delta^{J-1}$, then the asymptotic growth rate of $\mathbb{E}[M(n)]$ is $(\log n)^{J-1}$. Finally, in  Theorem \ref{extrema2} we relate -- under a relatively mild condition -- the $(\log n)^{J-1}$ asymptotic growth rate of the uniform case to the more general case where the $n$ elements are sampled iid from a generic distribution.

%In Theorem \ref{hullconv} we show that, as number of extrema of the convex hull grows to infinity, the convex hull tends to an apeirogon, a polytope with infinitely many sides. 




%we relate the $(\log n)^{J-1}$ asymptotic growth rate of the expectation of the number of identifiable components holds also when the $n$ elements are drawn from a generic distribution, under a very mild assumption. We relax this latter in Theorem . 
%For example, in \cite[section 7.3.3]{miller} a Dirichlet process prior is used to find that the number of components of a generic mixture model grows at a $\log n$ rate. If we consider a generic mixture model instead of the one outlined in section \ref{setup}, then we retieve a $\log n$ rate too via Theorem \ref{extrematopost}. This shows that the assumptions in Theorem \ref{extrematopost} are reasonable: they allow us to find a known growth rate by using novel techniques. 

%This implies a $(\log n)^{J-1}$ asymptotic growth rate of the expectation of the number of components of our admixture model in this more general setting. 
%Theorem \ref{growthdist} holds also in the more general settings of Theorems \ref{extrematopost} and \ref{extrema2}.

In section \ref{choq} we consider inference when the number of identifiable admixture components is equal to $J$, the dimension of the Euclidean space $\mathbb{R}^J$ we are working with, but the admixture components and weights are unknown. In Theorem \ref{thm1}, we use Theorem \ref{choq2} and Glivenko-Cantelli to retrieve the Choquet measure and thus also the identifiable admixture weights.


%a Dirichlet process retrieves the identifiable admixture weights. We also give the rate of convergence of the Dirichlet process posterior to the (Dirac at the) weights. 

In section \ref{algorithm}, we use the idea of mixture models based on the extremal set to formulate a novel algorithm that outputs an admixture model composed of only extremal elements, that is, an estimate of the richest cheap model. We state the objective function the algorithm optimizes, and provide a two-stage procedure. We apply this latter to the Associated Press data from the First Text Retrieval Conference (TREC-1), a large collection of terms used in 2246 documents.

Section \ref{concl} concludes our work. In Appendix \ref{app1} we discuss the similarities between de Finetti's theorem and Theorem \ref{choq2}, and we give an approximation of the joint distribution of the admixture components. We also provide the number of extremal elements of the convex hull within a unit simplex having the least amount of vertices.
%, we clarify the meaning of ``almost surely'' in equation \eqref{gk_cvg}, and we provide the general definition of Choquet simplex. 
We prove our results in Appendix \ref{proofs}.

%s \eqref{polya_eq} and
%We apply this procedure to the musiXmatch dataset, a large collection of song lyrics in bag-of-words format. For every word, we create $3$ multinomial bins representing high, medium, or low usage of the word in the tracks, and then we run our algorithm. It turns out that the richest cheap model has $9$ components.
%to solve this latter. In the first step, we fix an arbitrarily high number $L$ of components, and we run an EM procedure to estimate components and weights. Then we compute the convex hull of the estimated components; if the number $M$ of extrema of the convex hull is less than $L$, we run a second EM step using $M$ components. The components and weights estimated by this second step give us the richest cheap model. We apply this procedure to the musiXmatch dataset, a large collection of song lyrics in bag-of-words format. For every word, we create $3$ multinomial bins representing high, medium, or low usage of the word in the tracks, and then we run our algorithm. It turns out that the richest cheap model has $9$ components.

\iffalse
The chapter is organized as follows. In section \ref{mainresults} we state the results on the asymptotic growth rate of the expectation of the number of identifiable components of our model, and the CLT for the distribution of the number of identifiable components. We also show that as the number of extrema of a convex hull associated with our model goes to infinity, the convex hull tends to an apeirogon. In section \ref{choq} we bridge finite mixture model theory and Choquet theory by showing that, using a uniqueness result by Gustave Choquet, we can always retrieve weights in an admixture model. section \ref{algorithm} presents our algorithm to find the richest cheap model. We close with a discussion. In Appendix \ref{infiniteex}, we assume that the number $L$ of admixture components is known, but the parameters $\{f_1,\ldots,f_L\}$ of the multinomial components are not. We model them as random variables, and using de Finetti's theorem we find an approximation for their joint distribution. 
%when we drop the assumption that they are sampled iid from the uniform on the unit simplex. 
In Appendix \ref{app_2}, we give the number of extrema of the convex hull in $\Delta^{J-1}$ having the least amount of vertices. 
%In Appendix \ref{groups} we present the topics of the richest cheap model found in section \ref{algorithm}. We also give the frequency for the first ten words in each topic.
\fi





%A similar analysis holds if the scholar knows the sampling distribution of the estimated mixture components $\hat{f}_1,...,\hat{f}_{T(n)}$, where $\hat{f}_k$ is the empirical estimate of the $k$-th mixture element. 
%As a consequence to Theorem \ref{extrematopost}, we find that the rate of growth of the expected number of components in a generic finite mixture model is $\log n$.






\iffalse
In this paper, we are interested in understanding the asymptotic growth rate of the number of admixture components as one obtains more observations. We use well known results from the stochastic geometry literature on the asymptotic properties of random polytopes  \cite{calka,reitzner1,barany} to obtain a growth rate. A well studied object in stochastic geometry is
the convex hull, $K_n$, of $n$ points drawn iid uniformly from a smooth convex set in $\mathbb{R}^d$. Specifically, the growth rate of the number of extremal points of the convex hull $K_n$ with respect to $n$ is known \cite{calka,reitzner1,barany}. The central idea in this paper is one can couple Choquet theory \cite{phelps,phelps1} with the growth function of extremal points to provide the asymptotic growth rate of the number of admixture components. 

We now state a simple probability model that will provide insight on the growth function for the number of components. Sample $k$ admixture weights uniformly iid from the  unit simplex $\Delta^{(J-1)}$ 
\begin{align}
p_1,\ldots,p_k \sim \mbox{Uniform}(\Delta^{J-1}); \label{eq1_fmm}
\end{align}
the admixture weights and densities contained in the convex hull of these $k$ points determine the admixture densities that can be realized. 
Given $n$ observations, it does not make sense to set the number of mixture components greater than $n$, so $k \leq n$. If we consider the upper bound of $k=n$, then we sample $p_1,\ldots, p_n$  from the uniform distribution on the unit simplex of $\mathbb{R}^J$. By Choquet theory we know that any element in the convex hull of  $p_1,\ldots, p_k$  can also be represented by a convex combination of the  extrema of the convex hull.\footnote{In this setting, we may run into an identifiability problems when the number of extrema of $\mathscr{C}:=\text{Conv}(p_1,\ldots,p_k)$ is greater than $J$, the dimension of the Euclidean space we are working with. If that happens, we are not able to identify the weights of our mixture starting from a point inside the convex hull (as mentioned in \cite{lindsay}). However, this is not an issue in this paper. We simply point out how any point inside a convex hull in the unit simplex can be seen to represent a finite admixture distribution on $\mathbb{R}^J$. A similar approach to the geometry of mixture models is given in \cite{lindsay}.} So if we can obtain reasonable estimates for the number of extrema of the convex hull we obtain an estimate of the asymptotic growth rate of the number of mixture components. We adapt results on the number of extrema of the convex hull of random samples drawn from a convex body in $\mathbb{R}^d$  \cite{reitzner1} to prove the expectation of the growth function is  $(\log n)^{J-1}$ as well as a central limit theorem for the asymptotic growth function, Theorems \ref{growth} and \ref{growthdist}.

Besides convex geometry, the other tool that is central to our analysis is Choquet's theorem which states that, given a compact convex subset $C$ of a normed vector space, for any
$c \in C$ there is a probability measure $w$ supported on the set $E$ of extreme points of $C$ such that for any affine function $f$ on $C$
$$f(c) = \int_{e \in E} f(e) w(\mathrm{ d}e),$$ 
and $w$ is called the Choquet measure. There are two novel results in our paper with respect to convex geometry and Choquet theory. In Theorem \ref{hullconv}, we prove that  a convex hull generated by $n$ iid uniform random variables within a convex polytope tends to a smooth convex body as the number of  its extrema goes to infinity. We abuse
terminology: we say that a convex body is smooth if it has infinitely many sides
(that is, if it is an apeirogon). In Proposition \ref{empchoq}, we prove 
that the sequence of empirical probability measures on the extrema of our convex hull converges strongly to the Choquet measure.
\fi

%Our approach is somehow similar to that in \cite{hoff}. There, Hoff presents a method for estimating probability measures constrained to lie in a convex set. In particular, he uses the Choquet Theorem to point out that inference over a convex set of measures can be made via unconstrained inference over the set of extreme measures. The main difference with respect to our paper is that the convex hull he works with is a convex hull of probability measures, while the one we work with is a convex hull of points in a Euclidean unit simplex.
%Our approach is somehow similar to that in \cite{hoff}. There, Hoff presents a method for estimating probability measures constrained to lie in a convex set. In particular, he uses the Choquet Theorem to point out that inference over a convex set of measures can be made via unconstrained inference over the set of extreme measures. The main difference with respect to our paper is that the convex hull he works with is a convex hull of probability measures, while the one we work with is a convex hull of points in a Euclidean unit simplex.
% We first investigate the behavior of the extrema; in particular, we show that as the number of extrema grows to infinity, our convex hull converges to a smooth convex body. Here, we abuse terminology: we say that a convex body is smooth if it has infinitely many sides (that is, if it is an apeirogon). We also give a Central Limit Theorem (CLT) for the number of extrema. The  other route is more probabilistic in nature: we use Aldous' version of de Finetti's theorem (stated in \cite{aldous_2010}) to approximate the distribution of our random points $X_1,\ldots,X_n$. We also show how the sequence of empirical probability measures generated by our model converges to the Choquet measure, and how a P\'olya tree can approximate the latter.
%Our approach is somehow similar to that in \cite{hoff}. There, Hoff presents a method for estimating probability measures constrained to lie in a convex set. In particular, he uses the Choquet Theorem to point out that inference over a convex set of measures can be made via unconstrained inference over the set of extreme measures. The main difference with respect to our paper is that the convex hull he works with is a convex hull of probability measures, while the one we work with is a convex hull of points in a Euclidean unit simplex.
%In this paper we use many concepts developed in rather distant literatures, and we make them relevant in the study of finite admixture models. 
%The present work fits in the convex geometry approach, and stems from realizing that an element in the convex hull of $n$ points $X_1,\ldots,X_n$ drawn according to a uniform distribution on the unit simplex of $\mathbb{R}^J$ can be seen to represent an $m$-dimensional admixture model of component densities defined on $\mathbb{R}^J$, with $J \leq m \leq n$, as explained in section 2.1 (we denote by $m$ the number of extrema of the convex hull). This representation is useful since it allows to describe the growth rate of the expected number of components in a finite admixture model by studying the growth rate of the expected number of extrema of the convex hull. In particular, the main result of the paper, Theorem 3.3, states that the rate of growth of the expected number of components in a finite admixture model is given by $(\log n)^{J-1}$ (here $n$ is given by the cardinality of $\{X_1,\ldots,X_n\}$, where $X_1,\ldots,X_n \sim \mathcal{U}(\Delta^{J-1})$ iid, and $J$ is the dimension of the Euclidean space we are working on).
%The convex geometry approach is mainly due to Lindsay. In \cite{lindsay}, he focuses especially on multinomial and exponential family mixtures, with many examples on the family of binomial mixtures. He is the first to point out that a mixture model can be seen as an element of the unit simplex in some Euclidean space $\mathbb{R}^J$, a very deep insight which will prove fundamental for the present work. His main concerns are the identifiability of the weights of the mixture of multinomials and of distributions in the exponential family; he also gives a Carath\'eodory representation theorem for multinomial mixtures, which in a sense resembles the Choquet Theorem we use in this paper. He addresses also the problem of reducing the dimensionality of a mixture, the asymptotic multinomial geometry linked to the likelihood ratio of the first type, and the asymptotic mixture geometry.
%The former was pioneered by Amari and further inspected by Vos and Kass; in \cite{amari} and \cite{kass_vos}, they define a fiite mixture model as an element belonging to the mixture family, a subset of the space $\mathfrak{S}$ of statistical models. Then, after introducing statistical $\alpha$-connections, a concept that represents the intrinsic properties of the family of probability distributions, they point out that the $(-1)$-connection manifests the criterion that mixture families should be understood as straight models; that is, a mixture family can be regarded as a straight line connecting two distributions. They also introduce more concepts, such as $\alpha$-families and $\alpha$-curvature that, once applied to the case of mixture families (as before, when $\alpha=-1$), convey the information that finite mixture models can be represented as a flat smooth manifold.
%The present work fits in the convex geometry approach, and stems from realizing that an element in the convex hull of $n$ points $X_1,\ldots,X_n$ drawn according to a uniform distribution on the unit simplex of $\mathbb{R}^J$ can be seen to represent an $m$-dimensional admixture model of component densities defined on $\mathbb{R}^J$, with $J \leq m \leq n$, as explained in section 2.1 (we denote by $m$ the number of extrema of the convex hull). This representation is useful since it allows to describe the growth rate of the expected number of components in a finite admixture model by studying the growth rate of the expected number of extrema of the convex hull. In particular, the main result of the paper, Theorem 3.3, states that the rate of growth of the expected number of components in a finite admixture model is given by $(\log n)^{J-1}$ (here $n$ is given by the cardinality of $\{X_1,\ldots,X_n\}$, where $X_1,\ldots,X_n \sim \mathcal{U}(\Delta^{J-1})$ iid, and $J$ is the dimension of the Euclidean space we are working on).
%The rest of the paper is organized as follows. section 2 presents some of the properties of our model: we find that there exists a unique measure on the extreme points that allows to represent any point inside the convex hull, which we call the Choquet measure. We also give an approximation of the distribution of the sequence of points $X_1,\ldots,X_n$ we initially drew using Aldous' functional analysis version of de Finetti's theorem. section 3 is the main section of the paper: it deals with the behavior of the extreme points; we show that as the number of extrema grows to infinity, our convex hull approaches a smooth convex body; we also discover that the growth rate of the expected number of extrema is $(\log n)^{J-1}$, which corresponds to the growth rate of the expected number of components of our finite admixture model. We provide a Central Limit Theorem for the number of extreme points too, and we prove that the sequence of empirical probability measures with support on the extrema converges strongly to the Choquet measure. section 4 inspects how to approximate the Choquet measure using a P\'olya tree prior; we show that the P\'olya tree posterior is weakly consistent for the Choquet measure, and that the rate of convergence of the P\'olya tree posterior is given by $( \frac{\log n}{n} )^{\frac{\alpha}{2\alpha + 1}}$, for a properly defined $\alpha$. section 5 provides an application of our model to a classical problem in population and statistical genetics: we generalize the hierarchical model in \cite{Pritchard945} by using mixtures to approximate unknown distributions; we then use our main result to establish the growth rate of the expected number of components of these mixtures. section 6 concludes our work. There is an Appendix where the proofs to the results are given.






\section{Growth rates for extremal elements and admixture components}\label{mainresults}
\iffalse
The three main results of this paper are: Theorem \ref{growth} which states the growth rate of the extremal elements and number of components in the admixture model with uniform draws from the unit simplex, the model specified by \eqref{eq1_fmm}; 
Theorem  \ref{extrematopost} which states conditions under which the previous growth rate also holds for Bayesian non-parametric models; lastly, Theorem \ref{polya} that examines when a Bayesian inferential procedure with a prior on extremal points
recovers the Choquet measure. 


In this section, we build a bridge between finite mixture models and stochastic convex geometry. We give the growth rate of the expected number of identifiable components in our finite mixture model by studying the growth rate of the expected number of extrema of a convex body associated to our model. We also give a CLT for the number of identifiable components as a result of a CLT for the number of extrema of a convex body associated to our model. 
\fi
%Recall that $L$ denotes the number of admixture components, while $M$ denotes the number of identifiable admixture components, that is, those components that cannot be written as a convex combination of the other ones. We make the mild assumption that $M$ and $L$ are linearly related, that is, that we can find some $a \in \mathbb{Q}$ such that $M=aL$. This is indeed a mild assumption in that, given two natural numbers, we can always retrieve one as a linear combination of the other, using a rational coefficient.

Suppose the number $M$ of identifiable admixture components is a function $M(n)$ of the amount $n$ of data $x_1,\ldots,x_n$ that we sample. Such function is defined as follows. Let
\begin{equation}\label{eq1_fmm}
S_1,\ldots,S_n \stackrel{iid}{\sim} \mbox{Uniform}(\Delta^{J-1}), \quad J\geq 2,
\end{equation}
and call $K_n:=\text{Conv}(s_1,\ldots,s_n)$, where $s_j$ denotes the realization of $S_j$. In the stochastic convex geometry literature \cite{vu}, $K_n$ is called a \textit{random polytope}. Function $M(\cdot)$ is defined as
$$M:\mathbb{N} \rightarrow \mathbb{N}, \quad n \mapsto M(n):=\#\text{ex}(K_n),$$
that is, $M(n)$ is given by the cardinality of the extremal set of $K_n$. This construction allows us to explicitly relate the field of stochastic convex geometry with applied probability. A simple representation of the procedure to elicit $M(n)$ when $n=10$ and $J=3$ is given in Figure \ref{fig:procedure}.


\begin{figure}[h!]
\centering
\begin{subfigure}{.5\textwidth}
  \centering
  \includegraphics[width=.7\linewidth]{proc1}
  \caption{Figure \ref{fig:sub1}}
  \label{fig:sub1}
\end{subfigure}%
\begin{subfigure}{.5\textwidth}
  \centering
  \includegraphics[width=.7\linewidth]{proc2}
  \caption{Figure \ref{fig:sub2}}
  \label{fig:sub2}
\end{subfigure}
\caption{Suppose that we are working in $\mathbb{R}^3$. We sample $S_1,\ldots,S_{10}\sim \text{Uniform}(\Delta^2)$ iid; the realizations $s_1,\ldots,s_{10}$ are the ten gray points in the purple unit $2$-simplex in Figure \ref{fig:sub1}. Their convex hull $K_{10}$ is the orange polygon in Figure \ref{fig:sub2}. As we can see, it has five vertices, so $M(10)=5$.}
\label{fig:procedure}
\end{figure}







Before presenting the results in this section, we need to introduce the following geometric concepts.
\begin{itemize}
    \item As pointed out in \cite[Definition 2.1]{ziegler}, in higher-dimensional geometry, the faces of a polytope are features of all dimensions. A face of dimension $i$ is called an \textit{$i$-face}. For example, the polygonal faces of an ordinary polyhedron are $2$-faces. For any $h$-dimensional polytope, we have that $-1 \leq i \leq h$, where $-1$ is the dimension of the empty set. Let us  give a clarifying example. The faces of a cube comprise the cube itself ($3$-face), its facets ($2$-faces), the edges ($1$-faces), its vertices ($0$-faces), and the empty set (having dimension $-1$). Given a generic $h$-dimensional polytope $P$, we denote by $\mathcal{F}_i(P)$ one of its $i$-faces, $i \in \{-1,0,\ldots,h\}$.
    \item We call $\mathscr{F}_i(P)$ the \textit{collection of its $i$-faces}, and ${F}_i(P)$ the \textit{number of its $i$-faces}, that is, ${F}_i(P)=\#\mathscr{F}_i(P)$, for all $i$.
\end{itemize}

In view of the above definition of $0$-faces, we denote by $F_0(K_n)$ the number of extremal elements of $K_n$, so $M(n)=F_0(K_n)$. In the remainder of this section, we keep both notations to highlight the relationship between (stochastic convex) geometry and finite admixture models.


First we find the expected number of identifiable components and we show that it grows at rate $(\log n)^{J-1}$. Equations \eqref{first_theorem} and \eqref{equation_imp_first} are a consequence of \cite[Theorem 6]{reitzner} and \cite[Theorem 5]{barany2}, respectively.

\begin{theorem}{\textbf{(Growth rate of $\mathbb{E}[M(n)]$).}}\label{growth}
Let $K_n :=\text{Conv}({s}_1,\ldots,{s}_{n})$, where ${S}_1,\ldots,{S}_{n}$ are sampled as in \eqref{eq1_fmm}. Then,
\begin{align}\label{first_theorem}
\mathbb{E}[M(n)] = \mathbb{E} \left[F_0(K_n)\right] = \frac{J}{(J+1)^{J-1}}  (\log n)^{J-1} + \mathcal{O}\left((\log n)^{J-2}\log\log n\right),
\end{align}
where $\mathcal{O}$ denotes Bachmann–Landau big-O notation.
%and $\tau(\Delta^{J-1})$ is the number of towers of $\Delta^{J-1}$. 
In turn, this implies that
\begin{equation}\label{equation_imp_first}
    \lim_{n \rightarrow \infty} (\log n)^{-(J-1)}\mathbb{E} \left[M(n)\right] = \lim_{n \rightarrow \infty} (\log n)^{-(J-1)}\mathbb{E} \left[F_0(K_n)\right] = \frac{J}{(J+1)^{J-1}}  =: c(J).
\end{equation}
%where $T(\Delta^{J-1})$ is the number of towers of $\Delta^{J-1}$.
\end{theorem}

%\tau(\Delta^{J-1})  (J-1)!


Then, we see how, for $n$ large enough, the variance $\mathbb{V}[M(n)]=\mathbb{V}[F_0(K_n)]$ of the number of identifiable components can be approximated by $(\log n)^{J-1}$. Equation \eqref{var_eq} is a consequence of \cite[Theorem 1.3]{reitzner_var}.


\begin{theorem}{\textbf{(Approximation of $\mathbb{V}[M(n)]$).}}\label{var_thm}
Let $K_n :=\text{Conv}({s}_1,\ldots,{s}_{n})$, where ${S}_1,\ldots,{S}_{n}$ are sampled as in \eqref{eq1_fmm}. Then, 
\begin{equation}\label{var_eq}
    \mathbb{V}[M(n)]=\mathbb{V}[F_0(K_n)]=\mathcal{O}\left( (\log n)^{J-1} \right).
\end{equation}
\end{theorem}



If $J=3$, we have the following central limit theorem for the number of identifiable admixture components. 


\begin{theorem}{\textbf{(CLT for $M(n)$ when $J=3$).}}\label{clt_2}
Let $J=3$ and $K_n :=\text{Conv}({s}_1,\ldots,{s}_{n})$, where ${S}_1,\ldots,{S}_{n}$ are sampled as in \eqref{eq1_fmm}. Then, 
\begin{align}\label{clt_2_eq}
\begin{split}
    \lim_{n\rightarrow\infty} &\sup_{x\in\mathbb{R}} \left| \mathbb{P}\left( \frac{M(n)-\mathbb{E}[M(n)]}{\sqrt{\mathbb{V}[M(n)]}} \leq x \right) - \Phi(x) \right|\\
    &= \lim_{n\rightarrow\infty} \sup_{x\in\mathbb{R}} \left| \mathbb{P}\left( \frac{F_0(K_n)-\mathbb{E}[F_0(K_n)]}{\sqrt{\mathbb{V}[F_0(K_n)]}} \leq x \right) - \Phi(x) \right|=0,
\end{split}
\end{align}
where $\Phi$ denotes the cdf of a Standard Normal distribution.
\end{theorem}

Equation \eqref{clt_2_eq} is a consequence of \cite[Corollary 1.2]{pardon}. There, the author conjectures that it holds also for $J> 3$, but to the best of our knowledge such a conjecture has not been proven yet. To overcome this shortcoming, consider the following modification to our setup. Call $\mathbb{K}^2_+(J)$ the space of compact convex sets in $\mathbb{R}^J$, $J\geq 2$, with nonempty interior, boundary of differentiability class $\mathcal{C}^2$, and positive Gaussian curvature. That is, $\mathbb{K}^2_+(J)$ is the space of smooth compact convex sets in $\mathbb{R}^J$; the unit simplex ${\Delta}^{J-1}$ \textit{does not} belong to $\mathbb{K}^2_+(J)$. Pick then any $\epsilon>0$, and call $\hat{\Delta}^{J-1}_\epsilon$ a set in $\mathbb{K}^2_+(J)$ such that ${\Delta}^{J-1} \subset \hat{\Delta}^{J-1}_\epsilon$ and
$$d_H\left( \Delta^{J-1},\hat{\Delta}^{J-1}_\epsilon \right):=\max\left\lbrace{\sup_{x\in \Delta^{J-1}} d_2(x,\hat{\Delta}^{J-1}_\epsilon), \sup_{y\in \hat{\Delta}^{J-1}_\epsilon} d_2({\Delta}^{J-1},y)}\right\rbrace = \epsilon,$$
where $d_H$ denotes the Hausdorff distance and $d_2$ the Euclidean metric. Examples of  $\hat{\Delta}^{2}_\epsilon$ and $\hat{\Delta}^{3}_\epsilon$ are given in Figure \ref{fig_smooth_simpl}. Notice that set $\hat{\Delta}^{J-1}_\epsilon$ always exists, and that it is an $\epsilon$-approximation of ${\Delta}^{J-1}$ belonging to $\mathbb{K}^2_+(J)$. 

\iffalse
\begin{figure}[h!]
\centering
\includegraphics[width=.45\textwidth]{smooth_simpl}
\caption{For some $\epsilon>0$, $\hat{\Delta}^{2}_\epsilon$ is given by the triangle with round edges. It contains $\Delta^2$, the unit $2$-simplex in $\mathbb{R}^3$, colored in purple.}
\label{fig_smooth_simpl}
\centering
\end{figure}
\fi


\begin{figure}[h!]
\centering
\begin{subfigure}{.5\textwidth}
  \centering
  \includegraphics[width=.73\linewidth]{smooth_simpl.pdf}
  \caption{Figure \ref{fig:sub3}}
  \label{fig:sub3}
\end{subfigure}%
\begin{subfigure}{.5\textwidth}
  \centering
  \includegraphics[width=.5\linewidth]{michl1}
  \caption{Figure \ref{fig:sub4}}
  \label{fig:sub4}
\end{subfigure}
\caption{For some $\epsilon>0$, $\hat{\Delta}^{2}_\epsilon$ is given by the triangle with round edges in Figure \ref{fig:sub3} containing $\Delta^2$, the unit $2$-simplex in $\mathbb{R}^3$. $\hat{\Delta}^{3}_\epsilon$, instead, is given by the smooth tetrahedron in Figure \ref{fig:sub4} containing $\Delta^3$, the unit $3$-simplex in $\mathbb{R}^4$.}
\label{fig_smooth_simpl}
\end{figure}













Now, sample 
\begin{equation}\label{eq1_hat}
    \hat{S}_1,\ldots,\hat{S}_n \stackrel{iid}{\sim} \text{Uniform}(\hat{\Delta}^{J-1}_\epsilon), \quad J\geq 2,
\end{equation}
and call $\hat{K}_n:=\text{Conv}(\hat{s}_1,\ldots\hat{s}_n)$, where $\hat{s}_j$ is the realization of $\hat{S}_j$, $j\in\{1,\ldots,n\}$. Notice that in \eqref{eq1_hat} we could sample elements that are in $\hat{\Delta}^{J-1}_\epsilon\setminus {\Delta}^{J-1}$, but this happens with a probability that shrinks with $\epsilon$. 
%is more negligible the smaller $\epsilon$ is. 
Define then 
$$\hat{M}:\mathbb{N}\rightarrow\mathbb{N}, \quad n \mapsto \hat{M}(n):=\#\text{ex}(\hat{K}_n)=F_0(\hat{K}_n).$$
We can now give a version of Theorem \ref{clt_2} that holds for any $J\geq 2$. Equation \eqref{clt_gen_eq} is a consequence of \cite[Theorems 2, 6]{reitzner1}.

\begin{theorem}{\textbf{(CLT for $\hat{M}(n)$).}}\label{clt_gen}
Let $\hat{K}_n:=\text{Conv}(\hat{s}_1,\ldots\hat{s}_n)$, where $\hat{S}_1,\ldots,\hat{S}_n$ are sampled as in \eqref{eq1_hat}. Then, 
\begin{align}\label{clt_gen_eq}
\begin{split}
    &\left| \mathbb{P}\left( \frac{\hat{M}(n)-\mathbb{E}[\hat{M}(n)])}{\sqrt{\mathbb{V}[\hat{M}(n)]}} \leq x \right) - \Phi(x)\right| \\
    &=\left| \mathbb{P}\left( \frac{F_0(\hat{K}_n)-\mathbb{E}[F_0(\hat{K}_n)])}{\sqrt{\mathbb{V}[F_0(\hat{K}_n)]}} \leq x \right) - \Phi(x)\right| = \mathcal{O}\left( n^{-\frac{1}{2(J+1)}}\left(\log n\right)^{2+\frac{2}{J+1}} \right).
\end{split}
\end{align}
\end{theorem}

Then we prove that the number of identifiable admixture components concentrates around its expected value. Equation \eqref{conc_unif_eq} is a consequence of \cite[Theorem 2.11, Section 7]{vu}.


\begin{theorem}{\textbf{(Concentration inequality for $\hat{M}(n)$).}}\label{conc_unif_thm}
    Let $\hat{K}_n:=\text{Conv}(\hat{s}_1,\ldots\hat{s}_n)$, where $\hat{S}_1,\ldots,\hat{S}_n$ are sampled as in \eqref{eq1_hat}. Then, there are fixed positive constants $c,\Xi,\epsilon_0$ such that for any $\epsilon\in(0,\epsilon_0]$, $V\geq \Xi n^{\frac{J-1}{J+1}}$, $C\geq n\epsilon$, and $\lambda\in(0,\frac{V}{4C^2})$, the following holds
    \begin{align}\label{conc_unif_eq}
        \begin{split}
            \mathbb{P}&\left( \left| \hat{M}(n) -\mathbb{E}[\hat{M}(n)] \right| \geq \sqrt{\lambda V} \right)\\
            &=\mathbb{P}\left( \left| F_0(\hat{K}_n) -\mathbb{E}[F_0(\hat{K}_n)] \right| \geq \sqrt{\lambda V} \right) \leq 2 \exp(-\lambda/4)+\exp(-c\epsilon n) + \exp\left(-cn^\frac{J-1}{3J+5}\right).
        \end{split}  
    \end{align}
\end{theorem}
























\iffalse
$n^{\frac{J-1}{J+1}}$. Following \cite[page 2137]{reitzner_aop}, define the affine
surface area of $\Delta^{J-1}$ as $\Omega(\Delta^{J-1}):=\int_{\partial \Delta^{J-1}} \kappa(x)^{\frac{1}{J+1}}\text{d}x$, where $\partial \Delta^{J-1}$ is the boundary (in the Euclidean topology) of $\Delta^{J-1}$ and $\kappa(x)$ is the Gaussian curvature of $\partial \Delta^{J-1}$ at $x$ \cite{kuhnel}, and the constant (which depends only on the dimension $J$ of our Euclidean space)
$$\tau_J:=\frac{1}{2}\left( \frac{\kappa_{J-1}}{J+1}^{-\frac{2}{J+1}} \frac{(J^2+J+2)(J^2+1)}{(J+3)(J+1)!} \Gamma\left( \frac{J^2+1}{J+1} \right) \right),$$
where $\kappa_{J}$ is the volume of the $J$-dimensional unit ball and $\Gamma(\cdot)$ is the Gamma function.

\begin{theorem}\label{exp_pois}
If $J\geq 2$, then the following is true
\begin{equation}\label{exp_pois_eq}
    \mathbb{E}[F_0(K_n)] = \tau_J \Omega(\Delta^{J-1})n^{\frac{J-1}{J+1}}\left(1+o(1)\right)
\end{equation}
as $n$ grows to infinity.
\end{theorem}

\begin{corollary}\label{exp_pois_growth}
If $J\geq 2$, then the following is true
$$\lim_{n\rightarrow\infty} n^{-\frac{J-1}{J+1}}\mathbb{E}[F_0(K_n)] = \tau_J \Omega(\Delta^{J-1})\left(1+o(1)\right).$$
\end{corollary}

Then we prove that the number of identifiable admixture components concentrates around its expected value. 


\begin{theorem}\label{conc_unif_thm}
    If $J\geq 2$, then there exist fixed positive constants $c,\Xi,\epsilon_0$ such that for any $\epsilon\in(0,\epsilon_0]$, $V\geq \Xi n^{\frac{J-1}{J+1}}$, $C\geq n\epsilon$, and $\lambda\in(0,\frac{V}{4C^2})$, the following holds
    \begin{equation}\label{conc_unif_eq}
        P\left( \left| F_0(K_n) -\mathbb{E}[F_0(K_n)] \right| \geq \sqrt{\lambda V} \right) \leq 2 \exp(-\lambda/4)+\exp(-c\epsilon n) + \exp\left(-cn^\frac{J-1}{3J+5}\right)
    \end{equation}
\end{theorem}



\begin{proof}[Proof of Theorem \ref{conc_unif_thm}]
    Consider again the set $\mathbb{K}^2_+(d)$ of compact convex sets in $\mathbb{R}^d$, $d \geq 2$, having nonempty interior, boundary of differentiability class $\mathcal{C}^2$, and positive Gaussian curvature. Pick any $K \in \mathbb{K}^2_+$, and let $K_n$ denote the convex hull of $n\in\mathbb{N}$ random points sampled independently according to the uniform distribution in $K$. 
%$\Pi_n$ denote the convex hull of the intersection of $K$ with a Poisson point process with intensity $n$.
Then, in \cite[Theorem 2.11]{vu}, the author shows that
there exist positive constants $c,\Xi,\epsilon_0$ such that for any $\epsilon\in(0,\epsilon_0]$, $V\geq \Xi n^{\frac{J-1}{J+1}}$, $C\geq n\epsilon$, and $\lambda\in(0,\frac{V}{4C^2})$, the following holds
        $$P\left( \left| F_0(K_n) -\mathbb{E}[F_0(K_n)] \right| \geq \sqrt{\lambda V} \right) \leq 2 \exp(-\lambda/4)+\exp(-c\epsilon n) + \exp\left(-cn^\frac{J-1}{3J+5}\right).$$
        Since $\Delta^{J-1}\in \mathbb{K}_+^2(J)$, and because we assumed $J\geq 2$, then \eqref{conc_unif_eq} follows.
\end{proof}



Finally, we show that when the number of extreme points $F_0(\Pi_n)$ goes to infinity, $\Pi_n$ converges to an apeirogon, a polytope with infinitely many sides.

\begin{theorem}\label{hullconv2}
If $F_0(K_n)$ grows to infinity, then $K_n$ tends to an apeirogon.
\end{theorem}


***COPY THE PROOF HERE***













\iffalse
Call $PP(n)$ a stationary Poisson point process in $\mathbb{R}^J$ with intensity $n$, and $\Pi_n$ the convex hull of the intersection of $PP(n)$ with the unit simplex $\Delta^{J-1}$. Then,
$$M:\mathbb{N}\rightarrow\mathbb{N}, \quad n \mapsto M(n):=\#\text{ex}\Pi_n,$$
that is, $M(n)$ is given by the cardinality of the extremal set of $\Pi_n$. This means that $M(n)$ is equal to the number of vertices ($0$-faces) of $\Pi_n$, denoted by $F_0(\Pi_n)$. Then, we have the following important results.

First we find the expected number of identifiable components and we show that it grows at rate $n^{\frac{J-1}{J+1}}$. Following \cite[page 2137]{reitzner_aop}, define the affine
surface area of $\Delta^{J-1}$ as $\Omega(\Delta^{J-1}):=\int_{\partial \Delta^{J-1}} \kappa(x)^{\frac{1}{J+1}}\text{d}x$, where $\partial \Delta^{J-1}$ is the boundary (in the Euclidean topology) of $\Delta^{J-1}$ and $\kappa(x)$ is the Gaussian curvature of $\partial \Delta^{J-1}$ at $x$, and the constant (which depends only on the dimension $J$ of our Euclidean space)
$$\tau_J:=\frac{1}{2}\left( \frac{\kappa_{J-1}}{J+1}^{-\frac{2}{J+1}} \frac{(J^2+J+2)(J^2+1)}{(J+3)(J+1)!} \Gamma\left( \frac{J^2+1}{J+1} \right) \right),$$
where $\kappa_{J}$ is the volume of the $J$-dimensional unit ball and $\Gamma(\cdot)$ is the Gamma function.

\begin{theorem}\label{exp_pois}
If $J\geq 2$, then the following is true
\begin{equation}\label{exp_pois_eq}
    \mathbb{E}[F_0(\Pi_n)] = \tau_J \Omega(\Delta^{J-1})n^{\frac{J-1}{J+1}}\left(1+o(1)\right)
\end{equation}
as $n$ grows to infinity.
\end{theorem}




\begin{corollary}\label{exp_pois_growth}
If $J\geq 2$, then the following is true
$$\lim_{n\rightarrow\infty} n^{-\frac{J-1}{J+1}}\mathbb{E}[F_0(\Pi_n)] = \tau_J \Omega(\Delta^{J-1})\left(1+o(1)\right).$$
\end{corollary}
\fi 


Then we prove that after a finite number of observations, the number of identifiable admixture components concentrates around its expected value. 

\begin{theorem}\label{conc_ineq_pp}
Let $\psi\geq 0$ and $\beta(J):=4+\frac{2}{J+1}$. Then, for $n \geq c_1(J)$,
\begin{align}\label{conc_ineq_pp_eq}
\begin{split}
    P\bigg( \left| F_0(\Pi_n)-\mathbb{E}[F_0(\Pi_n)] \right| \geq &\psi \sqrt{\mathbb{V}[F_0(\Pi_n)]} \bigg)\\ &\leq 2 \exp \left( -\frac{1}{4} \min\left\lbrace{\frac{\psi^2}{2^{\beta(J)}},c_2(J) n^{\frac{J-1}{2\beta(J)(J+1)}}\psi^{\frac{1}{\beta(J)}}}\right\rbrace \right),
\end{split}
\end{align}
where $\mathbb{V}$ denotes the variance operator, and $c_1(J),c_2(J)\in(0,\infty)$ are constants depending (only) on $J$. Notice that, following \cite[Lemma 2]{reitzner1}, $\mathbb{V}[F_0(\Pi_n)]$ can be approximated by $n^{\frac{J-1}{J+1}}$.
\end{theorem}



We go on establishing a central limit theorem for the number of identifiable admixture components.

\begin{theorem}\label{clt_pois}
If $J\geq 2$, then the following is true
\begin{equation}\label{clt_pois_eq}
    \left| P\left( \frac{F_0(\Pi_n)-\mathbb{E}[F_0(\Pi_n)])}{\sqrt{\mathbb{V}[F_0(\Pi_n)]}} \leq x \right) - \Phi(x)\right| = \mathcal{O}\left( n^{-\frac{1}{2}+\frac{1}{J+1}}\left(\log n\right)^{2+\frac{2}{J+1}} \right),
\end{equation}
where $\mathcal{O}$ denotes Bachmann–Landau big-O notation.
\end{theorem}
\fi















\iffalse
%Here we characterize the growth rate of the expected number of extrema of $K_n$ and state geometric properties of asymptotic collection of extrema. 
%Throughout this subsection the number of extrema is equivalent to the number of mixture components via Choquet's theorem.

We first state the growth rate of the expected number of extrema of the convex body built as the convex hull of uniform draws from the unit simplex. 
%As explained in the Introduction, this immediately gives us the growth rate of the expected number of components of our mixture model. 
The growth rate is based on results in  \cite{barany2, reitzner}.
%that states that the expected number of $\ell$-faces, which we denote as $f_\ell$, of a convex hull of $n$ iid uniform random variables on a convex polytope $\mathscr{C}$ grows to infinity at the rate of $(\log n)^{d-1}$, where $d$ is the dimension of the Euclidean ball containing $\mathscr{C}$. 

Let us briefly introduce the concept of an $i$-face.  


%We call a chain $F_0(\Delta^{J-1})\subset F_1(\Delta^{J-1}) \subset \cdots \subset F_{J-1}(\Delta^{J-1})$ of $i$-dimensional faces $F_i$ of $\Delta^{J-1}$ a tower of $\Delta^{J-1}$.
\begin{theorem}\label{growth}
Let $K_n :=\text{Conv}({s}_1,\ldots,{s}_{n})$, where ${S}_1,\ldots,{S}_{n}$ are sampled as in \eqref{eq1_fmm}. Then,
\begin{align}\label{first_theorem}
\lim_{n \rightarrow \infty} (\log n)^{-(J-1)}\mathbb{E} \left[F_0(K_n)\right] = \frac{1}{(J+1)^{J-1}  (J-1)!} \tau(\Delta^{J-1}) =: c(J),
\end{align}
where $\tau(\Delta^{J-1})$ is the number of towers of $\Delta^{J-1}$.
\end{theorem}



Notice that $F_0(K_n)$ corresponds to $M(n)$. Theorem \ref{growth} tells us that the expected number of identifiable mixture components grows at rate $(\log n)^{J-1}$.
\fi




%Since $M(n) \equiv F_0(K_n)$ and $M(n)=bL(n)$, for some $b \in \mathbb{Q}$, then equation \eqref{first_theorem} implies that $\lim_{n \rightarrow \infty} (\log n)^{-(J-1)}\mathbb{E} \left[L(n)\right] = b \cdot c(J)$: the growth rate of the expected number of components of our mixture model is $(\log n)^{-(J-1)}$.

\iffalse
 the growth rate was stated for the number of $\ell$-dimensional faces of the convex body $K_n:=\text{Conv}(\mathscr{P}_n \cap K)$, and $\mathscr{P}_n$ is a Poisson point process with intensity $n$. If we denote 
$f_\ell(K_n)$ as  the number of $\ell$-dimensional faces then 
$$\lim_{n \rightarrow \infty} \mathbb{E} f_\ell(K_n)   = c_{d,\ell} (\log n)^{(d-1)},$$
where $c_{d,\ell}$ is a constant that depends on $K$, $\ell$, and $d$.

Note that in our case $K= \Delta^{J-1}$ is the unit simplex and a Poisson process on $\Delta^{J-1}$ with intensity $n$ is equivalent to sampling $\{p_1,\ldots,p_n\} \stackrel{iid}{\sim} \mathcal{U}(\Delta^{J-1})$.  So the results from
 \cite{barany} apply directly to our setting. Again $f_0$ is the number of $0$-dimensional faces or extremal elements and for the case $K= \Delta^{J-1}$ the constant $c_{J,\ell}$ can be computed and
 $$\lim_{n \rightarrow \infty} \mathbb{E} f_0(K_n)   = c_{J} (\log n)^{(d-1)}.$$



We can also state the limiting distribution of $F_0(K_n)$; specifically, we give the following central limit theorem for $F_0(K_n)$. It immediately implies the same result for $M(n)$. We denote by $\mathbb{V}[F_0(K_n)]$ the variance of the number of extreme points of $K_n$.
\begin{theorem}\label{growthdist}
Let $K_n :=\text{Conv}({s}_1,\ldots,{s}_{n})$, where ${S}_1,\ldots,{S}_{n}$ are sampled as in \eqref{eq1_fmm}. Then,
\begin{equation}\label{second_theorem}
\lim_{n \rightarrow \infty} \mathbb{P} \left( \frac{F_0(K_n) -\mathbb{E}[F_0(K_n)]}{\sqrt{\mathbb{V}[F_0(K_n)]}} \leq t \right) 
= \Phi({t})
\end{equation}
where $\Phi$ is the cumulative distribution function of the standard normal distribution.
\end{theorem}











We now show that the sequence of empirical probability measures on the extremal elements $\mathscr{E}$ converges strongly to the Choquet measure. The sequence of empirical measures we consider are constructed via the following procedure (P1). 
\begin{itemize}\label{P1}
\item[(1)] Sample  $\mathbf{p}_0   = \{p_1,\ldots,p_n\} \stackrel{iid}{\sim} \mathcal{U}(\Delta^{J-1})$;
\item[(2)] Compute the convex hull $\mathscr{C}:=$ Conv$(p_1,\ldots,p_n)$;
\item[(3)] Denote $\mathscr{E}$ as the set of extrema of $\mathscr{C}$ and $|\mathscr{E}| = k_n$;
\item[(4)] Denote  $\mu_0$ as a measure on $\mathscr{E}$ for which $p = \sum_{e \in  \mathscr{E}} e \cdot \mu_0(e)$, for all $p\in \mathbf{p}_0$;
\item[(5)] Denote a draw of $k_n$ points as $\mathbf{p} =  \{p_1,...,p_{k_n}\} \sim \mathcal{U}( \mathscr{C})$ iid;
\item[(6)] Generate $m$ samples $\{\mathbf{p}_1,...,\mathbf{p}_m \}$;
\item[(7)] Denote  $\mu_m$ as a measure for which $p = \sum_{e \in  \mathscr{E}} e \cdot \mu_m(e)$ for all $p\in \{ \cup_{j=0}^m \mathbf{p}_j\}$.
\end{itemize}


%The sequence of empirical probability measures on $\mathscr{E}$ is built as follows. We first sample $p_1,\ldots,p_n \sim \mathcal{U}(\Delta^{J-1})$ iid; then, we compute the convex hull $\mathscr{C}:=$ Conv$(p_1,\ldots,p_n)$, and we let $\mathscr{E}$ be the set of extrema of $\mathscr{C}$. After that, we consider $\mu_0$ on $\mathscr{E}$ such that any point in $\{p_1 \ldots,p_n\}$ is represented by $\mu_0$. Then, we sample $p_{k_{1_1}},\ldots,p_{k_{1_n}}$ uniformly from $\mathscr{C}$, and we let $\mu_1$ be the measure on $\mathscr{E}$ that represents any point in $\{p_1 \ldots,p_n\} \cup \{p_{k_{1_1}},\ldots,p_{k_{1_n}}\}$. Similarly, we sample $p_{k_{2_1}},\ldots,p_{k_{2_n}}$ uniformly from $\mathscr{C}$, and we let $\mu_2$ be the measure on $\mathscr{E}$ that represents any point in $\{p_1 \ldots,p_n\} \cup \{p_{k_{1_1}},\ldots,p_{k_{1_n}} \} \cup \{p_{k_{2_1}},\ldots,p_{k_{2_n}} \}$. By iterating this process, we build our sequence $\{\mu_n\}$.

\begin{proposition}\label{empchoq}
Let $\{\mu_m \}$ be the sequence of probability measures on $\mathscr{E}$ generated by procedure (P1).Then, $\mu_m \rightarrow \nu$ strongly as $m$ goes to infinity.
\end{proposition}
\begin{proof}
Consider the sequence generated by procedure (P1).  The sequence exists by the axiom of dependent choice. By construction, $\{\mu_m\}$ converges strongly to some $\mu \in \mathscr{P}(\mathscr{E})$, because $\mu_m$ has to agree with $\{\mu_1,\ldots,\mu_{m-1}\}$ in representing any element in $\bigcup_{i=0}^{m-1}\mathbf{p}_i$. The proof is by contradiction.

Assume  $\{\mu_m\}$  converges strongly to some $\mu \neq \nu$. Then there exists $\widetilde{p} \in  \bigcup\limits_{j \in \mathbb{N}} \mathbf{p}_j$ such that
$$\widetilde{p} \neq \sum\limits_{e \in \mathscr{E}} e \cdot \mu(e) \text{.}$$
Notice that $\bigcup_{j \in \mathbb{N}} \mathbf{p}_j=\mathscr{C}$, so we reach a contradiction.
 \end{proof}
\begin{remark}
The above result requires that the points we observe to update the sequence of measures belong to $\mathscr{C}$. If  we observe a point  in $\Delta^{J-1} \backslash \mathscr{C}$ then such a point would be a new extreme point and
change the support of our probability measures. We can overcome this problem by letting $\text{supp}(\mu_m)=\text{ex}(\Delta^{J-1})$, the extrema of the unit simplex $\Delta^{J-1}$ for all $m$. In this case, the sequence
$\{\mu_m\}$ will converge strongly to $\tilde{\nu}$ where $\tilde{\nu}$ is the Choquet measure on $\text{ex}(\Delta^{J-1})$.
\end{remark}


A further interesting result concerns the shape that $\hat{K}_n$ converges to asymptotically. The next theorem states that, as its number of vertices goes to infinity,  $\hat{K}_n$ converges to an apeirogon, a polytope with infinitely many sides.





%Suppose now we do not know the number of extreme points of $\mathscr{C}$. This situation corresponds to having a finite admixture model whose number of mixture components is unknown. We are interested in the rate at which the expectation of the number of extrema grows to infinity as the number $n$ of points drawn uniformly from the unit simplex grows to infinity. 
%Before finding such rate, we claim that the number of $\ell$-faces of $\mathscr{C}$, $f_\ell (\mathscr{C})$, is related to $f_0(\mathscr{C})$.
%\begin{lemma}\label{lemma1}
%In a generic Euclidean space $\mathbb{R}^d$, $f_\ell (\mathscr{C}) = k f_0(\mathscr{C})$, for some $k \geq 0$, $k=k(\ell)$, and $\ell \in \{1,\ldots,d-1\}$.
%\end{lemma}
%Notice that a flag is a maximal chain of faces, each a sub-face of the next in the chain. This last result allows us to say that the expected number of extrema of our set $\mathscr{C}$ grows at rate $(\log n)^{J-1}$, where $J-1$ is given by the dimension of our simplex, $\Delta^{J-1}$, and $n$ is given by the cardinality of $\{p_1,\ldots,p_n\}$, where $p_1,\ldots,p_n \sim \mathcal{U}(\Delta^{J-1})$ iid, as in (\ref{eq1_fmm}) (we assumed without loss of generality that $n=k$). This is also the growth rate of the expected number of components of the finite mixture our model represents.
%Denote $f_0(\mathscr{C})$ as the number of extreme points of the convex hull of $\mathscr{C}$  given the number of observations $n$ goes to infinity. In this paper we consider a convex body as smooth if it has infinitely many sides (that is, if it is an apeirogon). The following theorem proves that the extrema converge to an apeirogon.

\begin{theorem}\label{hullconv}
Let $\hat{K}_n :=\text{Conv}(\hat{s}_1,\ldots,\hat{s}_{n})$, where $\hat{S}_1,\ldots,\hat{S}_{n}$ are sampled as in \eqref{eq1_hat}. If $F_0(\hat{K}_n)$ grows to infinity, then $\hat{K}_n$ tends to an apeirogon.
\end{theorem}
\fi


%are going to explore some properties of our model. We first argue that a finite admixture model can be represented by a point in the convex hull of a finite sequence of points drawn uniformly from the unit simplex of a Euclidean space. Then, we provide what we call the Choquet properties of the geometric representation of a finite admixture model. Finally, we give a way of approximating the distribution of the sequence we drew in the first place to build our convex hull. 
%\subsection{Geometric representation of finite admixture models} 
%Here we explain how an element of a convex hull in a unit simplex can be seen to represent a finite admixture model. This representation is useful since it allows us to incorporate results developed in other literatures to the finite admixture models literature. In particular, it allows us to give the growth rate of the expected number of elements of the mixture by studying the growth rate of the expected number of extrema of the convex hull (Theorem 3.3). 

%Let us briefly review what a finite admixture model is. Given a random variable $Y$ that takes values $y \in \mathbb{R}^J$, we say that its distribution is given by (or is approximated by) an $m$-dimensional admixture model if:
%\begin{itemize}
%\item $f(y)=(f(y_1),\ldots,f(y_J))$, with $f(y_j) \geq 0$ for all $j \in \{1,\ldots,J\}$ and $\sum_{j=1}^J f(y_j)=1$; %Essentially, this corresponds to saying that $Y$ is a discrete random variable with $J$ possible values;
%\item $f(y)=\sum\limits_{k=1}^m p_k f_k(y),$ where the $p_k$'s are nonnegative and sum to $1$, and the $f_k$'s are such that, for all $k \in \{1,\ldots,m\}$,
%\end{itemize}
%\begin{align}
%\begin{split}
%f_k(y)&=(f_k(y_1),\ldots,f_k(y_J)) \text{, with } f_k(y_j) \geq 0 \text{ for all } j \in \{1,\ldots,J\}\\ &\text{and }\sum_{j=1}^J f_k(y_j)=1. \label{eq5}
%\end{split}
%\end{align}
%Such a finite mixture of densities on $\mathbb{R}^J$ can be represented by an element of a convex hull in the unit simplex of $\mathbb{R}^J$. To see this, draw $n$ points uniformly from the unit simplex $\Delta^{J-1}:= \{ \mathbf{x} \in \mathbb{R}^J : \sum_{j=1}^J x_j=1 \text{ and } x_j \geq 0 \text{, for all } j \}$
%\begin{align}
%X_1,\ldots,X_n \distas{iid} \mathcal{U}(\Delta^{J-1}) \text{,} \label{eq1_fmm}
%\end{align}
%and construct the convex hull $\mathscr{C} := \text{Conv}(X_1,\ldots,X_n)$ of these random points. It is going to be a polytope with $m$ extreme points, $J \leq m \leq n$. Then, any point inside the convex hull represents a finite admixture model, that is, it represents $f=\sum_{k=1}^m p_k f_k \text{,}$ where $f_k$ is the density of the $k^{\text{th}}$ mixture component, $p_k \in (0,1)$ for all $k$, $\sum_{k=1}^m p_k=1$, and $m$ is the number of components in the finite admixture model. Let us give an example.

%Suppose we are in $\mathbb{R}^J$, for some $J \in \mathbb{N}$, and suppose we have:
%\begin{align*}
%f=a f_1 + b f_2 + (1-a-b) f_3, \qquad a,b \in (0,1), \qquad a+b<1,
%\end{align*}
%where $f_1$, $f_2$, and $f_3$ are as in (\ref{eq5}). Then, suppose we want to compute $f(x)$, for some $x \in \mathbb{R}^J$. If we compute $f_1(x)$, $f_2(x)$, and $f_3(x)$ and calculate their convex hull, we have that ${f}(x)$ belongs to the convex hull.
%f_1 &\sim \mathcal{N}(\mu_1,\Sigma_1), \qquad f_2 \sim \mathcal{N}(\mu_2,\Sigma_2), \qquad f_3 \sim\mathcal{N}(\mu_3,\Sigma_3),

\subsection{From the Uniform to the general case}\label{from-to}
In this section we generalize Theorem \ref{growth} to the case where $S_1,\ldots,S_n$ are iid samples from a generic distribution $G$.\footnote{That is, $G$ is any distribution that is absolutely continuous with respect to the Lebesgue volume measure of the simplex.} We defer the study of the non-iid case to future work. We ask ourselves whether the asymptotic growth function of the expected number of identifiable components based on draws from the uniform distribution can inform us about the growth rate based on draws from a generic distribution. 
%Let us drop the uniform assumption in \eqref{eq1_fmm}; s

Suppose $S_1,\ldots,S_n$ are now sampled iid from a generic distribution $G$ on $\Delta^{J-1}$. Call then $\breve{K}_n:=\text{Conv}(s_1,\ldots,s_n)$ and define
$$T:\mathbb{N} \rightarrow \mathbb{N}, \quad n \mapsto T(n):=\#\text{ex}(\breve{K}_n)=F_0(\breve{K}_n),$$
so we assume that the number of identifiable admixture component corresponds to the number of vertices of $\breve{K}_n$. In a way, we can see $T(n)$ as a ``generalization'' of $M(n)$; of course, $T(n)$ and $M(n)$ may be different.

%We only focus on the uniform distribution over $\Delta^{J-1}$ as the $\hat{\Delta}_\epsilon^{J-1}$ case is analogous.



%At this point, an obvious question is what the asymptotic growth function of the expected number of identifiable components based on draws from the uniform distribution on the unit simplex tells us about the asymptotic growth rate based on draws from a generic distribution.
%in a real statistical problem. 
%That is, we want to drop the assumption that ${f}_1,\ldots,{f}_{M(n)}$ are sampled iid from $\mathcal{U}(\Delta^{J-1})$. We generalize the result in Theorem \ref{growth}. 




%When we drop the uniform assumption in \eqref{eq1_fmm}, the number of extrema of the convex hull may be different from $M(n)$. We denote this quantity by $T$, and it is too going to be a function of the amount $n$ of data we collect, that is, $T=T(n)$. In particular, suppose that $S_1,\ldots,S_n$ are now sampled iid from a generic distribution $G$ on $\Delta^{J-1}$. Call then $\breve{K}_n:=\text{Conv}(s_1,\ldots,s_n)$. Function $T$ is defined as 
%$$T:\mathbb{N} \rightarrow \mathbb{N}, \quad n \mapsto T(n):=\#\text{ex}(\breve{K}_n)=F_0(\breve{K}_n).$$
%that is, $T(n)$ is given by the cardinality of the extremal set of $\breve{K}_n$.

%In Theorem \ref{extrematopost} we state a condition on $T(n)$ that allow us to directly apply the result we have for the uniform case to obtain the $(\log n)^{J-1}$ growth rate of the number of identifiable mixture components. related to our model

\begin{remark}
In this work, we only consider the case of $M(n),T(n)\geq J$ (of course, $J \geq 2$). This because, if that were not the case, our convex hull would be a (proper) subset of a smaller-dimensional Euclidean space than $\mathbb{R}^J$. 
%We defer to future work the study of the case where $M(n)$ or $T(n)$ can be smaller than $J$. We conjecture that such a possibility may be linked to model misspecification, since potentially that could imply that the number of identifiable components is smaller than the $J$ many values that we deemed possible for $X$ to take values on.
It is worth to notice, though, that {\em if $n \geq J$ and is sufficiently large, then $M(n),T(n)\geq J$ holds almost surely}. Let us be more precise. In Theorem \ref{growth}, we show that $\mathbb{E}[M(n)]$ grows as $(\log n)^{J-1}$, as $n\rightarrow \infty$. By concentration of measure arguments, then, it follows that the probability of having only \( J \) extreme points is vanishingly small. 
\iffalse 
Specifically, let \( A_n \) be the event that the number of extreme points is exactly \( J \), that is $A_n=\{M(n)=J\}$. Then,
\[
\mathbb{P}(A_n) \to 0 \quad \text{as } n \to \infty.
\]
Let us show this. 
\fi 
More formally, if we assume \( n \geq J \), it follows that for sufficiently large \( n \), the expected number of extreme points satisfies
\[
\mathbb{E}[M(n)] \approx c(J) (\log n)^{J-1} \geq J.
\]
Thus, by Markov's inequality and standard deviation bounds, we have that

\[
\mathbb{P}(M(n) < J) \rightarrow 0 \quad \text{as } n \rightarrow \infty.
\]
That is,

\[
\mathbb{P}(M(n) \geq J) \rightarrow 1 \quad \text{as } n \rightarrow \infty.
\]
Therefore, almost surely, as we sample more and more points uniformly from $\Delta^{J-1}$, the number $M(n)$ of extreme elements of $K_n$ is greater than \( J \). A similar argument holds for $T(n)$ as a consequence of Theorem \ref{extrema2}. 

This argument shows that, while it is necessary to assume that $M(n),T(n) \geq J$ if we consider a fixed generic sample size $n$, this condition is satisfied almost surely when $n \rightarrow \infty$.
\end{remark}


%Notice that $M(n),T(n)\geq J$ (of course, $J \geq 2$). If that is not the case, we can still have a convex hull, but it will be a proper subset of a smaller dimensional Euclidean space, and we are not interested in this eventuality.




%We now relax the assumption in \eqref{eq1_fmm}. 
%A fundamental quantity of interest in Bayesian density estimation is the posterior distribution $p(\theta \mid  x_1,\ldots,x_n)$ of the admixture components as the number of observations  $x_1,\ldots,x_n$ goes to infinity, when the number of components itself depends on the number of observations.
%In particular we will study the increase in the cardinality of $\{p_1,\ldots,p_k\}$ as $n$ goes to infinity. 
%The main idea is that we can use the result in Theorem \ref{growth} to  prove results in a more general setting, that is, a setting in which we do not require $S_1,\ldots,S_n$ to be uniformly distributed on the unit simplex $\Delta^{J-1}$. 
\iffalse
In Theorem \ref{extrematopost}, we require that the expected number of identifiable admixture components $\mathbb{E}[T(n)]$ in the general setting is in a fixed linear relation with the expected number of identifiable components $\mathbb{E}[M(n)]$ in the simple uniform model. This can be interpreted as the stochasticity around the number of components entering the  general model through the simple uniform one, and then being linearly ``passed on''. 
% to the more general one. 
In Theorem \ref{extrema2}, we weaken this rather strong regularity condition. 
%We only require that the sequence of rational numbers relating the expected number of identifiable components in the uniform and in the general models does not have $0$ as an accumulation point. 

%It is not standard for a result on the extremal set of convex bodies to translate into result on the posterior distribution of a Bayesian procedure. 
%As we pointed out before, the number $T(n)$ of extrema of the convex hull in this more general setting may be different from the one we had for the uniform case, that is, $T(n)$ may be different from $M(n)$. In the following theorem we state a regularity conditions on $T(n)$ 
%and one on $\mathbb{P}(f_1,\ldots,f_{L(n)} \mid  x_1,\ldots,x_n)$ 
%such that we can apply Theorem \ref{growth}. 





\begin{theorem}\label{extrematopost}
Suppose that, for all $n$, 
%and all observations $x_1,\ldots,x_n$, 
we can always find $\gamma \in \mathbb{R}\setminus\{0\}$ such that $\mathbb{E}[T(n)]=\gamma \cdot \mathbb{E}[M(n)]$.
%\item[(i)] There exists a uniform lower bound on the density function of the posterior $\mathbb{P}(f_1,\ldots,$ $f_{M(n)} \mid  x_1,\ldots,x_n)$ over the support $\mathfrak{S}_n$, for all $n$ and all observations $x_1,\ldots,x_n$.
%\item[(ii)] For all $n$ and all observations $x_1,\ldots,x_n$, support $\mathfrak{S}_n$ is a subset of the convex hull of the finite union of $k_0$ many convex bodies $K_{j,n}$ whose number of extrema depend on $n$, that is, $\mathfrak{S}_n \subset \text{Conv} (\bigcup_{j=1}^{k_0} K_{j,n})=:\mathscr{K}_n$.
%\end{itemize}
Then,
\begin{equation}\label{gen_post}
\lim_{n \rightarrow \infty} (\log n)^{-(J-1)} \mathbb{E}[T(n)] = c^{\prime}(J,\gamma):=\gamma \cdot c(J).
\end{equation}
\end{theorem}

%Here the constant $c^{\prime\prime}$ will be larger than the one in Theorem \ref{growth}.




there is a uniform lower bound $\rho$ on the posterior density, that is, $\mathfrak{g}(\tilde{f}_1,\ldots,\tilde{f}_{M(n)}) \geq \rho$, for all $\tilde{f}_1,\ldots,\tilde{f}_{M(n)} \in \mathfrak{S}_n$, where $\mathfrak{g}$ is the density function of  $\mathbb{P}(f_1,\ldots,f_{M(n)} \mid  x_1,\ldots,x_n)$. Then, we can adapt Theorem \ref{growth} to obtain
$$ \lim_{n \rightarrow \infty} (\log n)^{-(J-1)} \mathbb{E}[F_0(K_n)] = c^{\prime}(J,b),$$
where $c^{\prime}(J,\rho)$ is now also be a function of the uniform lower bound $\rho$. Notice that $K_n$ is the convex hull of $\tilde{f}_1,\ldots,\tilde{f}_{M(n)}$, distributed according to $\mathbb{P}(f_1,\ldots,f_{M(n)} \mid  x_1,\ldots,x_n)$.


We now relax the assumption that support $\mathfrak{S}_n$ is convex to condition (ii), which states that the support belongs to the convex hull of the union of $k_0$ convex bodies whose number of extremal elements depend on $n$. By Theorem \ref{growth}, we know that for each body $K_{j,n}$ with $j \in \{1, \dots,k_0\}$, we have that $\lim_{n \rightarrow \infty} (\log n)^{-(J-1)} \mathbb{E}[F_0(K_{j,n})] = c(J)$. This implies that 
$$\lim_{n \rightarrow \infty} (\log n)^{-(J-1)} \mathbb{E} \left[ F_0 \left( \mathscr{K}_n \right) \right] = k_0 c(J).$$
This, together with condition (i), immediately retrieves equation \eqref{gen_post}, with $c^{\prime\prime}=k_0 c^{\prime}(J,\alpha)$.
\fi

\iffalse
The assumption that for all $n$ we can always find $\gamma \in \mathbb{Q}$ -- not depending on $n$ -- such that $\mathbb{E}[T(n)]=\gamma \cdot \mathbb{E}[M(n)]$
% implies that $\mathbb{E}[T(n)]=\gamma \cdot b \cdot \mathbb{E}[L(n)]$, since in section \ref{mainresults} we pointed out that $M(n)=bL(n)$, for some $b \in \mathbb{Q}$.
%the Introduction we assumed that we can alwys find $a \in \mathbb{Q}$ such that $M=aL$.
%In addition, such an assumption 
is mild. 
%Indeed, we can always obtain a natural number as a linear function of another one, using a rational coefficient. 
It means that there is a fixed (linear) relation between the expected number of identifiable components of the uniform and the general models. This entails that all the stochasticity around the number of identifiable components enters the model through the number of extrema of the simple uniform model. Then, it is ``passed'' to the number of extrema of the more general model through coefficient $\gamma$. 
%In turn, this gives us the (expected) number of components of our finite mixture model.
\fi

%Although the assumption in Theorem \ref{extrematopost} is mild, it can further be weakened, as we show in the next theorem. 

\iffalse
is still necessary, as we show in this next example. 
\begin{example}\label{ex_gam}
Let $\mathbb{E}[M(n)]$ be some function $\xi(n)$ of $n\in\mathbb{N}$, and let 
$$\mathbb{E}[T(n)]=\begin{cases}
J+\lceil{\log \log n}\rceil  & \text{, for all } n \leq \tilde{n}\\
J+\lceil{\log n}\rceil  & \text{, for all } n > \tilde{n}\\
\end{cases},$$
for some $\tilde{n} \in \mathbb{N}$. Then, we can find $\gamma \in \mathbb{Q}$ such that $J+\lceil{\log \log n}\rceil=\gamma \cdot \xi(n)$, for all $n \leq \tilde{n}$, and also $\gamma^\prime \in \mathbb{Q}$ such that $J+\lceil{\log n}\rceil=\gamma^\prime \cdot \xi(n)$, for all $n > \tilde{n}$. But we may have that $\gamma \neq \gamma^\prime$. In this case, we cannot use the previous theorem to determine the growth rate of $\mathbb{E}[T(n)]$. 
\demo
\end{example}
%We give now a more general version of Theorem \ref{extrematopost}.
In light of Example \ref{ex_gam}, we present a generalization of Theorem \ref{extrematopost}.
\fi

The next result, Theorem \ref{extrema2}, states the following. Up until the $(N-1)$-th data point, the expected number of identifiable components $\mathbb{E}[T(n)]$ in the more general case can take on any possible real value. From the $N$-th observation onward, though, it must be in a fixed (possibly highly nonlinear) relationship $\gamma_n$ with $\mathbb{E}[M(n)]$. If this happens, we are able to relate their growth rates. Such an assumption is made primarily for mathematical convenience: with it, deriving the result in equation \eqref{gen_gr_rate} becomes relatively easy. In the future, we plan to relax this assumption in order to derive a more general result.

\begin{theorem}{\textbf{(Growth rate of $\mathbb{E}[T(n)]$).}}\label{extrema2}
Call $(\gamma_n)$ a sequence in $\mathbb{R}^\mathbb{N}$ for which $0$ is not an accumulation point, and let $\mathbb{E}[T(n)]=g_n(\mathbb{E}[M(n)])$, where $g_n$ is a functional on $\mathbb{R}$ that depends on $n$. Then, if there exists $N\in\mathbb{N}$ such that for all $n\geq N$, $g_n(\mathbb{E}[M(n)])=\gamma_n \mathbb{E}[M(n)]$, we have that
%Consider sequence $(\gamma_n) \in \mathbb{Q}^\mathbb{N}$ such that $0$ is not an accumulation point and $\mathbb{E}[T(n)]=\gamma_n \mathbb{E}[M(n)]$, for all $n \in \mathbb{N}$. Then,
%Consider the sequence $(\gamma_n) \in \mathbb{Q}^\mathbb{N}$ such that $\mathbb{E}[T(n)]=\gamma_n \mathbb{E}[M(n)]$, for all $n \in \mathbb{N}$. Suppose that $(\gamma_n)$ is such that $0$ is not an accumulation point, and call $(r_{\gamma_n})$ a sequence such that $0$ is not an accumulation point and $\gamma_n/r_{\gamma_n} \rightarrow \zeta$ as $n \rightarrow \infty$, $\zeta\in\mathbb{R}$. Then,
\begin{equation}\label{gen_gr_rate}
\lim_{n \rightarrow \infty} \frac{1}{{\gamma_n} (\log n)^{J-1}} \mathbb{E}[T(n)] =  c(J).
\end{equation}
\end{theorem}


In \eqref{gen_gr_rate}, $c(J)$ is a quantity that only depends on the dimension $J$ of the Euclidean space $\mathbb{R}^J$ that we consider, and it is the same as in \eqref{equation_imp_first}. The following corollary is a direct consequence of Theorem \ref{extrema2}.

\begin{corollary}\label{extrema2_cor}
Suppose the assumptions of Theorem \ref{extrema2} hold. Then, if there exists a sequence $\varpi_n \in \mathbb{R}^\mathbb{N}$ such that $\gamma_n=\mathcal{O}(\varpi_n)$, then the growth rate of $\mathbb{E}[T(n)]$ is $\varpi_n(\log n)^{J-1}$.
\end{corollary}






%is free to ``act crazy'', but from the $N$th observation onward, it must ``set its mind straight'' and be in a steady relationship $\gamma_n$ with $\mathbb{E}[M(n)]$. If this happens, we are able to relate their growth rates.
%growth rate of $\mathbb{E}[T(n)]$ to that of $\mathbb{E}[M(n)]$.


%, we let the expected number of identifiable components in the more general case be ``free'' for the first $N-1$ 

%5Hence, the expected number of identifiable components in this more general case grows at rate ${\gamma_n}(\log n)^{J-1}$: the growth rate depends on ${\gamma_n}$ as well.
%, the limiting behavior of the sequence of rationals that links the simple uniform model to the general one.

\iffalse
In both Theorems \ref{extrematopost} and \ref{extrema2} we do not need to put restrictions on the support $\mathfrak{S}_n$ of the posterior of the identifiable admixture components, that is, on $\mathfrak{S}_n := \text{supp}(\mathbb{P}(f_1,\ldots,f_{T(n)} \mid  x_1,\ldots,x_n))$. Indeed, suppose the worst case scenario takes place, that is, $\mathfrak{S}_n$ is a fractal. Working with probability measures whose support is a fractal is a particular challenging task, see e.g. \cite{prats}. In the context of the present work however, the fractal is contained in the unit simplex. We give an instance of a fractal contained in the unit $3$-simplex in $\mathbb{R}^4$ in Figure \ref{fig4}. We can then consider the extension $\mathfrak{P}$ of $\mathbb{P}(f_1,\ldots,f_{L(n)} \mid  x_1,\ldots,x_n)$ to the unit simplex $\Delta^{J-1}$, 
%$$\mathbb{P}(f_1,\ldots,f_{L(n)} \mid  x_1,\ldots,x_n)=\frac{\mathfrak{P}(f_1,\ldots,f_{L(n)})}{\int_{\mathfrak{S}_n} \{f_1,\ldots,f_{L(n)}\} \text{ d}\mathfrak{P} },$$
%for all $\{f_1,\ldots,f_{L(n)}\} \subset \mathfrak{S}_n$, 
and then sample from it. 
%Integrals over fractals are inspected e.g. in \cite{edgar}, a complete account of the subject, and in \cite{stric}, where the author uses self-similarity to ease computations. 
This procedure is only viable because $\mathfrak{S}_n$ is a subset of $\Delta^{J-1}$, a convex compact body, for all $n$, by construction. 


\begin{figure}[h!]
\centering
\includegraphics[width=.75\textwidth]{frac_1}
\includegraphics[width=.25\textwidth]{frac_2}
\caption{This figures are reproduced from \cite[Figures 6 and 7]{roshchina}. We consider the unit simplex  $\Delta^3$ in $\mathbb{R}^4$, and consider a sphere such that the edges of $\Delta^3$ touch the sphere. Then, we consider the intersection of the sphere with $\Delta^3$. After that, we continue slicing off spherical caps in such a way that they are tangential to the existing slices. The resulting body is a spherical fractal.}
\label{fig4}
\centering
\end{figure}
\fi



 
\begin{remark}
It is immediate to see that there is a universal upper bound for the Euclidean distance between two points in a unit simplex: for all $x,y \in \Delta^{J-1}$, $d_2(x,y) \equiv \|x-y\| \leq 2$, where $\|\cdot\|$ denotes the Euclidean norm. This gives us an interesting result: the Hausdorff distance between $K_n$ and $\breve{K}_n$ has a universal upper bound as well. Indeed,  
\begin{align*}
    d_H({K}_n,\breve{{K}}_n)&=  \max\left\lbrace{ \sup_{{x} \in {K}_n} d_2 (x,\breve{{K}}_n), \sup_{y \in \breve{{K}}_n} d_2({K}_n,y)}\right\rbrace  \\
    &=\max\left\lbrace{ \sup_{{x} \in {K}_n} \inf_{y \in \breve{{K}}_n}d_2 (x,y), \sup_{y \in \breve{{K}}_n} \inf_{x \in {K}_n} d_2(x,y)}\right\rbrace \leq 2.
\end{align*}
%The same holds for $d_H(\hat{K}_n,\breve{{K}}_n)$.

Notice also that if instead of requiring $\mathbb{E}[T(n)]=\gamma_n \mathbb{E}[M(n)]$, for all $n\geq N$,   we are willing to make the slightly stronger assumption that for all $n\geq N$, $T(n)=\rho_n  M(n)$, $(\rho_n) \in \mathbb{R}^\mathbb{N}$ possibly different from $(\gamma_n)$, then we retrieve Theorem \ref{clt_2} for $T(n)$.\footnote{Of course we need to assume that for $(\rho_n)$, $0$ is not an accumulation point.} This because, since $F_0(K_n) = M(n)$, we have that 
$$\frac{T(n) -\mathbb{E}[T(n)]}{\sqrt{\mathbb{V}[T(n)]}} = \frac{\rho_n F_0(K_n) - \mathbb{E}[\rho_n F_0(K_n)]}{\sqrt{\mathbb{V}[\rho_n F_0(K_n)]}} = \frac{F_0(K_n) -\mathbb{E}[F_0(K_n)]}{\sqrt{\mathbb{V}[F_0(K_n)]}},$$
and so Theorem \ref{clt_2} follows. A similar argument allows us to retrieve Theorem \ref{clt_gen} when we work with $\hat{K}_n$ instead of ${K}_n$.
%In a similar fashion, if in Theorem \ref{extrema2} we require that, for all $n \geq N$, $T(n)=\rho_n \cdot M(n)$, $(\rho_n) \in \mathbb{R}^\mathbb{N}$ possibly different from $(\gamma_n)$, then we retrieve Theorems \ref{clt_2} and \ref{clt_gen}.
\end{remark}

%Our result reduces in the mixture model case to a $\log n$ growth rate.  An interesting question is how restrictive are assumptions (i) and (ii) in Theorem \ref{extrematopost} and how do these assumptions relate to classic assumptions in Bayesian non-parametric analysis and posterior consistency.

%Next, notice how finite mixture models are just a degenarate version of finite admixture models, where the weights are just $0$'s and $1$'s. Therefore, our results carry over to finite mixture models. In particular, since in finite mixture models we are not bound to sample weights from the unit simplex, we retrieve the rate of growth $\log n$ for the expected number of extrema. In \cite[section 7.3.3]{miller}, the authors show that the Dirichlet process mixture model gives a $\log n$ growth rate for the number of components of a finite mixture model. Hence, our result for the expected number of extrema matches a result in the literature concerning the {actual} number of extrema of a finite mixture models.


\section{Choquet measure and identifiable admixture weights}\label{choq}

In this section we build a bridge between finite admixture models and Choquet theory. We show how, thanks to a uniqueness result by Gustave Choquet (i.e. Theorem \ref{choq2}), the classical Glivenko-Cantelli's theorem can be used to retrieve the identifiable admixture weights. 
%ADD? We also give the rate of convergence of the empirical cdf to the cdf of the Choquet measure.

%Dirichlet process posterior to the Dirac measure on the weights.


%Suppose the identifiable admixture components $\{f_\ell\}_{\ell=1}^M$ are known, but their weights $\{\varphi_{i,\ell}\}_{\ell=1}^M$ are not. 
%To deal with this problem, we give an interesting result that relates Choquet theory to finite mixture models. 
By Theorem \ref{choq2}, we have that for every element $p$ in a simplex $C$, there exists a unique measure -- that we call the Choquet measure associated with $p$, and denote by $\nu_p$ -- supported on the extremal elements $E=\text{ex}(C)$ such that $p=\sum_{e \in E} e \cdot \nu_p(e)$.\footnote{We write $\nu_p(e)$ in place of $\nu_p(\{e\})$ for notational convenience. We stick to this abuse of notation for the rest of the paper.} In our analysis, $p$ corresponds to $\pi$, the elements $e$ in $E=\text{ex}(C)$ correspond to the identifiable $f_\ell$'s, and the $\nu_p(e)$'s correspond to the weights of the identifiable $f_\ell$'s, that is, $\nu_{\pi}(f_\ell)=\varphi_{\ell}$, for every identifiable $f_\ell$.

In Theorem \ref{thm1}, we show that if we only assume that the number $M$ of components is known and equal to $J$ -- the dimension of the Euclidean space $\mathbb{R}^J$ we are working with -- we can use Glivenko-Cantelli to retrieve $\nu_{\pi}$. The $\varphi_{\ell}$'s represent the weights of the richest cheap finite admixture model representing $\pi$, 
%That is, the model representing $\pi_i$ having the largest amount of components that cannot be written as a convex combination of one another. 
so $\pi=\sum_{\ell=1}^M f_\ell \nu_{\pi}(f_\ell)$. 
%ADD? By retrieving, we mean that given a Dirichlet process prior $DP(\alpha P_0)$ specified on the distributions supported on $\Delta^{J-1}$ having parameter $\alpha>0$ and $P_0$ as base measure, its posterior converges weakly almost surely to the Dirac at $\nu_{\pi_i}$. 
%We also give the rate of convergence.
% of the Dirichlet process posterior to the Choquet measure.
%This means that if we keep sampling $f_1,\ldots,f_k \mid \mu \sim \mu$ iid, where supp($\mu)=\{f_1,\ldots,f_M\}$, then eventually the P\'olya tree posterior $\Pi_\alpha(\cdot \mid f_1,\ldots f_k)$ retrieves the Dirac on the ``right'' measure $\nu_{\pi_i}$ that gives the $\varphi_{i,\ell}$'s.

%As \cite[section 2.3.3]{lindsay} points out, if the number $M$ of identifiable components $f_1,\ldots,f_M$  of our mixture model is greater than $J$, the dimension of the Euclidean space we are working with, then the weights $\varphi_{i,\ell}$ of the mixture model are not identifiable. Theorem \ref{polya} allows to overcome this problem.




\begin{remark}
Recall that $\pi=\sum_{\ell=1}^M \varphi_{\ell} f_\ell =\sum_{\ell=1}^L \phi_{\ell} f_\ell $, where we labeled the unidentifiable components as $f_{M+1},\ldots,f_L$, $M \leq L$. This is without loss of generality.
\end{remark}


%\subsection{Choquet theory and extrema of convex bodies}\label{back} 

%In this section, we provide the relevant background material on Choquet theory. 
%As we explained in the Introduction, we use Choquet theory to retrieve the weights of the mixture model, so we assume the components $f_1,\ldots,f_L$ or $\tilde{f}_1,\ldots,\tilde{f}_L$ of the mixture model to be known. 
Let us denote by 
$$\mathcal{K}_M:=\text{Conv}(f_1,\ldots,f_M)=\text{Conv}(f_1,\ldots,f_L)$$ 
the convex hull generated by the $M$ identifiable admixture components, and assume $M=J$, so that $\mathcal{K}_M$ is a {\em simplex}. Recall that the latter is defined as $\mathcal{K}_M \coloneqq \{\pi=\sum_{\ell=1}^M \xi_\ell f_\ell : \sum_{\ell=1}^M \xi_\ell =1 \text{, } \xi_\ell \geq 0 \text{ } \forall \ell \in \{1,\ldots,M\} \text{, and } M=J\} \subset \mathbb{R}^J$, while the {\em $(J-1)$-unit simplex} satisfies the extra condition that the $f_\ell$'s are the $J$ standard unit vectors in $\mathbb{R}^J$. An example of a simplex within the unit $2$-simplex in $\mathbb{R}^3$ is given in Figure \ref{fig2}.  Our first goal is to learn about distributions supported on the extremal elements of $\mathcal{K}_M$, $E_M:=\text{ex}(\mathcal{K}_M)=\{f_1,\ldots,f_M\}$.
\begin{figure}[h!]
\centering
\includegraphics[width=.45\textwidth]{simplex}
\caption{A triangular-shaped convex hull within the unit $2$-simplex in $\mathbb{R}^3$. It is a simplex because the number of its vertices coincides with the dimension of the Euclidean space.}
\label{fig2}
\centering
\end{figure}

\iffalse
The first result we present is the famous Choquet theorem \cite[Theorem, page 14]{phelps1}. It states that for every element $x$ of a metrizable compact convex subset $X$ of a locally convex space $S$, there exists a probability measure $\mu$ supported by the extrema $\text{ex}(X)$ of $X$ that represents $x$. That is, $f(x)=\int_{\text{ex}(X)} f(e) \text{ }\mu(\text{d}e)$, for all continuous and affine functions $f$ on $X$.
\fi 

Since $\Delta^{J-1}$ is locally convex, and $\mathcal{K}_M \subset \Delta^{J-1}$ is a metrizable compact convex set, then thanks to Theorem \ref{choq2}, we know that for every $\pi \in \mathcal{K}_M$, there exists a unique probability measure $\nu_{\pi}$ supported on $E_M$ such that $\pi=\sum_{f_\ell \in E_M} f_\ell \cdot \nu_{\pi}(f_\ell)$. We formalize this statement in the next proposition.


%It states that if $V$ is a Banach space, and if we consider a compact convex subset $C \subseteq V$, then, for all $c \in C$ there exists a probability measure $\mu$ such that $\text{supp}(\mu)=exC$, and for all affine functions $f$ on $C$ we have:
%$$f(c)= \int_{exC} f(e) \text{ } \mu(\mathrm{d}e) \text{.}$$

%We will call such a $\mu$ an extreme measure. In our case, $V=\Delta^{J-1} \subset \mathbb{R}^{J}$, which is complete with respect to the Euclidean norm, and $C=\mathscr{C}$, which is compact by Heine-Borel; indeed, it is closed (its complement is open) and bounded (it is contained in $\Delta^{J-1}$, and so also in the $(J-1)$-unit ball).


\begin{proposition}{\textbf{(Choquet uniqueness for admixture models)}.}\label{simplex_conseq}
If $\mathcal{K}_M$ is a simplex, then every element $\pi \in \mathcal{K}_M$ can be represented by a unique measure $\nu_{\pi}$ (the Choquet measure representing $\pi$) supported on $E_M$.
\end{proposition}
For every identifiable admixture component $f_\ell\in E_M$, the Choquet measure $\nu_{\pi}$ gives the corresponding admixture weight, that is, $\nu_{\pi}(f_\ell)=\varphi_{\ell}$.


\iffalse
\begin{remark}
For any convex set built as the convex hull of $M$ points in the simplex $\Delta^{J-1}$, we can always find a subset that is of particular interest. Indeed, consider $M$ points $f_1,\ldots,f_M$ in $\Delta^{J-1}$ that cannot be written as a convex combination of one another; call $\mathcal{K}_M$ their convex hull and $E_M:=ex\mathcal{K}_M$. Let us also denote the set of probability measures on $(E_M,\sigma(E_M))$ by $\Delta(E_M,\sigma(E_M))$, where $\sigma(E_M)=2^{E_M}$. Then by Proposition \ref{simplex_conseq}, if $\mathcal{K}_M$ is a simplex, we have that for all $p \in \mathcal{K}_M$, there exists a unique measure $\nu_p \in \Delta(E_M,\sigma(E_M))$ such that $p=\sum_{f_j \in E_M} f_j \cdot \nu_{p}(f_j)$. Hence, we can always retrieve the set $N_{\mathcal{K}_M} \subset \Delta(E_M,\sigma(E_M))$ of Choquet measures associated with the (elements of the) convex set $\mathcal{K}_M$. It is immediate to see that there is a bijection between $N_{\mathcal{K}_M}$ and $\mathcal{K}_M$. An important consequence of Proposition \ref{simplex_conseq} is that the elements of $\mathcal{K}_M$ inherit the properties of the elements of $E_M$. The Choquet measures are the ``channel'' through which the properties are ``carried'' from the extrema to the points in the convex hull. 
%To this extent, the set $N_{\mathcal{K}_M}$ is an interesting set to study. 
\end{remark}
%All the results given in this section hold for $\tilde{K}_M$ as well.
\fi

%Of course, the same holds if we consider the convex hull of the elements $\tilde{f}_1,\ldots,\tilde{f}_M$.



%for all $q_j \in \{q_1,\ldots,q_n\}$, there exists a unique measure $\nu_{q_j} \in \mathscr{P}(E_n,\sigma(E_n))$ -- the Choquet measure for $q_j$ -- such that $q_j=\sum_{e \in E_n} e \cdot \nu_{q_j}(e)$. Actually, a more general claim is true. Again by Proposition \ref{simplex_conseq}, we have that 





\iffalse
One perspective of de Finetti's theorem \cite{aldous_2010} is that a mixture model can be considered as making inference over a convex set of measures or inference over the set of extreme measures. This idea was operationalized and pushed further in
\cite{hoff} to execute inference over a convex set of measures via unconstrained inference over the set of extreme measures. In this spirit we will consider non-parametric Bayesian analysis with a classic prior, the P\'olya tree prior. 


\subsection{Choquet theory for admixture weights}\label{choq_mix}
%In this section, we show that if we consider a Dirichlet process over the distributions supported on $\Delta^{J-1}$ and if $\mathcal{K}_M$ is a simplex, then for every element $p \in \mathcal{K}_M$ we can select a Dirichlet process that converges weakly to the (Dirac at the) Choquet measure $\nu_p$ associated with $p$. We also provide the rate of convergence.


%It is natural to ask whether placing a prior on the extreme points of $\mathscr{C}$ we do recover the Choquet measure, and -- if that is the case -- at which rate. In this section, we use a non-parametric technique, namely the P\'olya tree, to approximate the Choquet measure.
%We show that the posterior of a P\'olya tree on the distributions on the extrema of $\mathscr{C}$ is weakly consistent for the Choquet measure, and we give the rate of convergence of the P\'olya tree posterior to the Choquet measure.



%P\'olya tree priors were introduced in \cite{ferguson1974} as a special case of tail free processes, and further developed using de Finetti's theorem in \cite{lavine1992},  \cite{lavine1994} and \cite{mauldin1992}. As 
The Dirichlet process (DP) is of fundamental importance in the Bayesian nonparametrics literature \cite{ferguson1974, ghosal_vdv, ghosh}. It is a default prior on spaces of probability measures, and a building block for priors on other structures. Because it ``selects'' almost surely discrete distributions, it is a sensible choice for our case, since the Choquet measure is supported on the finite set $E_M$. The DP possesses the conjugacy property, as shown in the following theorem from \cite{ferguson1974}.



\begin{theorem}
Suppose that $P\sim DP(\alpha P_0)$, for some $\alpha>0$ and base measure $P_0$, and we collect iid observations $X_1,\ldots,X_n\mid P \sim P$. Then, (a version of) the posterior for $P$ is given by $DP(\alpha P_0 + \sum_{j=1}^n \delta_{X_j}) \equiv DP(\alpha P_0 +nP_n)$. 
\end{theorem}
Here $P_n=1/n \sum_{j=1}^n \delta_{X_j}$ denotes the empirical distribution of the observations. An extensive treatment of DP's is given in \cite[Chapter 4]{ghosal_vdv}. The following is the main result of the section; 


\iffalse
P\'olya tree processes are a large class of priors that includes the Dirichlet processes, and provide a flexible framework for Bayesian analysis of non-parametric problems  .
P\'olya tree processes have a tractable expression for the posterior while allowing for the incorporation of a wide range of beliefs.  P\'olya tree priors can be specified so as to give nonzero probability to continuous distributions which is
not possible with the Dirichlet process prior, which . The standard construction of a P\'olya tree prior is as follows. Given an interval $\Omega$, recursively bisect it into subintervals. At each stage,
the probability mass already assigned to an interval is randomly divided and assigned into its subintervals according to the independent draw of a Beta random variable. The parameters in the Beta distribution are set to 
increase rapidly as the partitions become finer.


Recall that in Theorem \ref{hullconv} we stated that as the extremal elements of the convex body $K_n$ (the convex hull of $n$ points drawn from the unit simplex) grow to infinity, $K_n$ converges to a smooth convex body; call this latter $\mathscr{K}$. The following theorem states the (positive) answer to
the natural question of whether Bayesian updating with a  P\'olya tree prior $\mu$ on the extrema of $\mathscr{K}$ recovers the Choquet measures.
\fi



\iffalse
It states that if $\mathcal{K}_M$ is a (Choquet) simplex, 
%, that is, if $M=J$, 
then for all $\pi_i \in \mathcal{K}_M$ we can select a DP that always retrieves the Choquet measure $\nu_{\pi_i}$. This is a useful result because, despite Proposition \ref{simplex_conseq} tells us that every $\pi_i \in \mathcal{K}_M$ can be represented by a unique $\nu_p$ supported on the extrema of $\mathcal{K}_M$, it does not give the analytic form of $\nu_p$. 

%In section \ref{intro}, we wrote our model as $\pi_i=\sum_{\ell=1}^L \phi_{i,\ell} f_\ell$. Equivalently, if we consider the identifiable components, that is, those that are not convex combinations of other components, then we can write $\pi_i=\sum_{\ell=1}^M \varphi_{i,\ell} f_\ell$. 
%Here, $M=aL$, for some $a \in \mathbb{Q}$, and we cannot write any $f_i \in \{f_1,\ldots,f_M\}$ as a convex combination of the other ones. 
%Notice that the $\varphi_{i,\ell}$'s are such that $\sum_{\ell=1}^M \varphi_{i,\ell}=1$, and $\varphi_{i,\ell} \geq 0$, for all $\ell$. 
The Choquet measure $\nu_{\pi_i}$ accruing to $\pi_i$, is important because it gives us the weights $\varphi_{i,\ell}$'s, that is $\nu_{\pi_i}(f_\ell)=\varphi_{i,\ell}$, for all $\ell \in \{1,\ldots,M\}$. This means that, using the DP approach described in the next theorem, we are able to identify the weights of the richest cheap finite admixture model.
%, that is, the mixture model that represents $\pi_i$ using the largest amount of identifiable components. 
%They cannot be written as convex combinations of the other ones.
\fi








%Also, $\nu_p$ allows to compute the weights $\phi_i$ for our mixture model, 












%as described in the Introduction. The same holds if $K_n$ is the convex hull of $\tilde{f}_1,\ldots,\tilde{f}_M$; in that case, $\nu_{\tilde{p}}$ allows to compute the weights $\tilde{\phi}_i$ for the mixture model.
\begin{theorem}\label{thm1}
%Let $\mathcal{K}_M:=\text{Conv}(f_1,\ldots,f_L)$ and call $E_M$ its extrema. Pick any element $\pi_i \in \mathcal{K}_M$. 
Let $\mathcal{K}_M$ be a (Choquet) simplex. If $e_1,\ldots,e_k$ are an iid sample of elements of $E_M$ from $\nu_{\pi_i}$, then 
\begin{equation}\label{polya_eq}
DP \left(\alpha P_0 + \sum_{j=1}^k \delta_{e_j} \right) \xrightarrow[k \rightarrow \infty]{w} \delta_{\nu_{\pi_i}} \quad \text{a.s.}
\end{equation}
where $\alpha$ is a positive real, $P_0$ is a base measure supported on $\Delta^{J-1}$, $\xrightarrow[]{w}$ denotes the weak convergence, and $\delta_{\nu_{\pi_i}}$ is the Dirac measure at $\nu_{\pi_i}$. In addition, the rate of convergence relative to the total variation metric is given by $k^{-1/2}$.
%Notice that $P_0$ is supported on the whole unit simplex $\Delta^{J-1}$. So in Theorem \ref{thm1} we abused notation by letting $\nu_{\pi_i} \equiv \overline{\nu}_{\pi_i}$, where $\overline{\nu}_{\pi_i}$ is a probability measure supported on $\Delta^{J-1}$ such that $\overline{\nu}_{\pi_i}(E_M)=1$, $\overline{\nu}_{\pi_i}(E_M^c)=0$, and $\overline{\nu}_{\pi_i}(A)={\nu}_{\pi_i}(A)$, for all $A\subset E_M$.}
\end{theorem}
The idea is that if we keep observing iid samples $e_j$'s from the Choquet measure $\nu_{\pi_i}$ -- that is, if we are able to observe identifiable admixture components $\{f_1,\ldots,f_M\}$ sampled according to their true weights $\{\nu_{\pi_i}(f_1),\ldots,\nu_{\pi_i}(f_M)\}\equiv\{\varphi_{i,1},\ldots,\varphi_{i,M}\}$ -- then the DP recovers $\nu_{\pi_i}$. Because the DP is a measure over measures, the formal statement is convergence to the Dirac at the Choquet measure $\delta_{\nu_{\pi_i}}$. 
%We need base measure $P_0$ to be supported on $\Delta^{J-1}$ because it has to give positive probability to all the elements of the unit simplex.

\begin{remark}\label{rem_proc}
Let $\text{supp}(P)$ denote the support of a generic measure $P$. Notice that in Theorem \ref{thm1} $\text{supp}(P_0)=\Delta^{J-1}$, while $\text{supp}(\nu_{\pi_i})=E_M\subset\Delta^{J-1}$. This is not a problem since by \cite[Theorem 4.15]{ghosal_vdv} the weak support of a Dirichlet process on measures on $\Delta^{J-1}$ is given by $\{P: \text{supp}(P)\subset \text{supp}(P_0)\}$.

It is also worth to mention that the Choquet measure can be retrieved using a different approach than Theorem \ref{thm1}. 
\end{remark}
\fi



We are now ready for the main result of this section. Suppose that $\mathcal{K}_M$ is a simplex, so that $M=J$, and call $E_M$ the set of extremal elements of $\mathcal{K}_M$. We denote by $e_j$ a generic element of $E_M$, that is, a generic identifiable admixture component. In mathematical terms, we write $e_j\in\{f_1,\ldots,f_M\}=:E_M$, for all $j\in\mathbb{N}$. 

Consider an iid sample from the Choquet measure $\nu_{\pi}$, that is, $e_1,\ldots,e_K \sim \nu_{\pi}$ iid. Because every $e_k$, $k\in \{1,\ldots,K\}$, corresponds to an identifiable admixture element, we can write $e_k=f_{\ell_k}$, where label $\ell_k$ belongs to $\{1,\ldots,M\}$, for all $k\in\{1,\ldots,K\}$. Now, call $\zeta$ the distribution of the labels; we immediately notice that there is a one-to-one correspondence between $\zeta$ and $\nu_{\pi}$ since $\zeta(\ell)=\nu_{\pi}(f_{\ell})$, for all $\ell\in\{1,\ldots,M\}$. Call then $F_\zeta$ the cdf of $\zeta$, and denote by $\ell_1,\ldots,\ell_K\sim \zeta$ the iid sample of labels from $\zeta$ that corresponds to the iid sample $e_1,\ldots,e_K$ from the Choquet measure $\nu_{\pi}$. For all $x\in\mathbb{R}$, denote by
$$F_K(x):=\frac{1}{K}\sum_{j=1}^K \mathbb{I}_{[\ell_j,\infty)}(x)$$
the empirical cdf of $\ell_1,\ldots,\ell_K$, where $\mathbb{I}_A(x)$ stands for the indicator function of $x$ belonging to a generic set $A$. Then, we have the following. Equation \eqref{gk_cvg} is a consequence of the Glivenko-Cantelli's Theorem \cite{vdv_asy}. 
%, and the rate of convergence comes from \cite[Chapter 1, Remarks 1 and 2]{zhou}.

\begin{theorem}{\textbf{($F_K$ retrieves $F_\zeta$)}.}\label{thm1}
    Let $\mathcal{K}_M$ be a simplex in $\Delta^{J-1}$, call $E_M:=\text{ex}(\mathcal{K}_M)$ the set of its extremal elements, so that $M=J$. Suppose that, in the notation introduced above, $\ell_1,\ldots,\ell_K\sim \zeta$ iid. Then, we have that 
\begin{equation}\label{gk_cvg}
    \sup_{x\in\mathbb{R}} \left| F_K(x)-F_\zeta(x) \right| \xrightarrow[K\rightarrow\infty]{a.s.} 0.
\end{equation}
In addition, pick any $0 < \eta < 0.5$. Then, $\sup_{x\in\mathbb{R}} \left| F_K(x)-F_\zeta(x) \right| = o \left( \frac{1}{K^\eta} \right)$ a.s.
%In addition, it goes to $0$ as $\frac{8k+4}{\exp(\frac{\varepsilon^2}{8}k)}$.
\end{theorem}






%by Glivenko-Cantelli's Theorem \cite{vdv_asy}, we have that $F_k$ converges to $F_\zeta$ uniformly almost surely, that is, In addition, \cite[Chapter 1, Remarks 1 and 2]{zhou} shows that the rate of convergence is $ke^{-k}$. This rate is faster than that of Theorem \ref{thm1}, but $(f_1,\ldots,f_M)$ is not an exchangeable sequence since in this alternative procedure the labels ``matter''.\footnote{For a definition of exchangeable sequence, see Appendix \ref{infiniteex}.} As usual, there is no free lunch.
%We refer the reader that is interested in the rate of convergence of \eqref{gk_cvg}  
The idea in Theorem \ref{thm1} is that if we keep observing iid samples $e_k$'s from the Choquet measure $\nu_{\pi}$ -- that is, if we are able to observe identifiable admixture components $\{f_1,\ldots,f_M\}$ sampled according to their true weights $\{\nu_{\pi}(f_1),\ldots,\nu_{\pi}(f_M)\}\equiv\{\varphi_{1},\ldots,\varphi_{M}\}$ -- then their empirical cdf recovers the cdf of $\nu_{\pi}$. In turn, this gives us the identifiable admixture weights, since -- as we saw before -- $\nu_{\pi}(f_\ell)=\varphi_{\ell}$, for all $\ell\in\{1,\ldots,M\}$.





\iffalse
We can also provide the rate of convergence to the Choquet measure $\nu_{\pi_i}$. We require the density $f_{\nu_{\pi_i}}$ corresponding to the Choquet measure to satisfy some regularity conditions. In particular we  consider functions in the class 
$\mathcal{C}^\alpha[0,1]$, $\alpha \in (0,1]$, that denotes the H\"older functions on the interval $[0,1]$, that is, 
$$\mathcal{C}^\alpha[0,1]:= \left\lbrace{ g:[0,1] \rightarrow \mathbb{R} , \sup\limits_{x \neq y \in [0,1]} \frac{|g(x)-g(y)|}{|x-y|^\alpha} < \infty }\right\rbrace \text{.}$$
Let us denote by $\mathcal{E}:=\bigcup_{l \geq 0} \{0,1\}^l \cup \{\emptyset\}$, and by $\#$ the cardinality operator.



\begin{theorem}\label{thm1}
%Let $\mathcal{K}_M:=\text{Conv}(f_1,\ldots,f_L)$ and call $E_M$ its extrema. 
Let $\mathcal{K}_M$ be a simplex. Pick any element $\pi_i \in \mathcal{K}_M$, and sample $e_1,\ldots,e_k \sim \nu_{\pi_i}$ iid, where  $e_1,\ldots,e_k$ are elements of $E_M$, and $\nu_{\pi_i}$ is the Choquet measure associated with $\pi_i$. Clearly, it may be the case that $e_j=e_s$, for some $j,s \in \{1,\ldots,k\}$. We assume that density $f_{\nu_{\pi_i}}$ belongs to $\mathcal{C}^\alpha[0,1]$, $\alpha \in (0,1]$, and is bounded away from zero on $[0,1]$.

Let $\Pi$ be the prior on densities supported on $E_M$ generated  by a P\'olya tree random measure with respect to the canonical dyadic partition of $[0,1]$ with parameters $\mathcal{A}=\{ \alpha_\epsilon : \epsilon \in \mathcal{E} \}$ chosen as $\alpha_\epsilon=a_{\#\epsilon} \vee 8$, for any $\epsilon \in \mathcal{E}$, with
$$a_l=l2^{2l\alpha} , \quad l \geq 0.$$
Then, as $k \rightarrow \infty$, for any $\mathcal{M}_k \rightarrow \infty$, we have that
\begin{equation}\label{polya_speed}
\mathbb{E}^k_{f_{\nu_{\pi_i}}} \left[ \Pi \left( f: \|f - f_{\nu_{\pi_i}}\|_{\infty} \leq \mathcal{M}_k \epsilon^*_{k,\alpha} \mid e_1\ldots,e_k \right) \right] \rightarrow 1 \text{,}
\end{equation}
where $\epsilon^*_{k,\alpha}$ is the minimax rate of convergence
$$\epsilon^*_{k,\alpha} = (\log k/{k} )^{{\alpha}/(2\alpha +1)}.$$
\end{theorem}



Theorem \ref{thm1} states that for a  prior specified by a P\'olya tree random measure the posterior density concentrates in a supremum norm ball around the Choquet density. The rate of convergence of the P\'olya tree posterior to the Choquet measure $\nu_{\pi_i}$ is given by $\epsilon^*_{k,\alpha}$.
%Also note that the rate of convergence of the P\'olya tree posterior to the Choquet measure on the extrema is the same we had for the convergence of the P\'olya tree posterior to the Choquet measure on the extrema of the convex hull.

%The same results hold for $\tilde{K}_M$ as well.


\section{Implications for admixture models} 
In this section, we compare the growth rate of the expected number of components of a finite admixture model we retrieved via Theorem 3.3 to the growth rate of the number of components obtained by using a Dirichlet process mixture model, as described in \cite{li}. We then give an example of how our model can be useful in statistical and population genetics.

In \cite[section 7.3.3]{miller}, Miller and Harrison point out that the Dirichlet process mixture model gives a $\log n$ growth rate for the number of components of a finite mixture model. Theorem 3.3 tells us that the growth rate of the {expected} number of components of a finite admixture model is given by $(\log n)^{J-1}$.

We would expect these rates to be close together, if not the same (because finite admixture models are a subclass of finite mixture models). We realize that they are actually the same in the $2$-dimensional Euclidean space, while the {expected} number of components has a higher growth rate as the dimension of the Euclidean space is higher than $2$, more so the higher the dimension. This is due to the way we built our model: it pays the price of taking into account the dimensionality of the Euclidean space where the unit simplex from which we sample our $X_j$'s lies, and this is reflected in the growth rate of the expected number of extreme points of our convex hull.

We now give an example explaining how the result we gave in Theorem 3.3 may be applied to the first popularly used admixture model \cite{Pritchard945}. In their paper, Pritchard, Stephens, and Donnelly use a Bayesian clustering approach to identify the subpopulations and to probabilistically assign individuals to the populations on the basis of their genotype, while simultaneously estimating population allele frequencies. An essential aspect of the method is that admixed individuals are allowed and the admixture probabilities can be inferred for each individual.

We will use the notation in  \cite{Pritchard945} and write $(x_l^{(i,1)},x_l^{(i,2)} )$ for the genotype of the $i$-th individual at locus $l$, $i \in \{1,\ldots,N \}$, $l \in \{ 1,\ldots,L \}$. Each observed allele copy $x_l^{(i,a)}$ is modeled to be originated in some unknown population $z_l^{(i,a)}$, so that $z_l^{(i,a)}$ represents the population of origin of allele copy $x_l^{(i,a)}$. $q_k^{(i)}$ denotes the proportion of individual $i$'s genome that originated from population $k$, and the full hierarchical model is given as follows:
\begin{align}
&\mathbb{P}\left(x_l^{(i)}=j \mid Z,P,Q \right)=p_{z_l^{(i,a)}lj} \text{;} \nonumber \\
&\mathbb{P}\left( z_l^{(i,a)}=k \mid P,Q\right)=q_k^{(i)} \text{;} \nonumber \\
&\mathbf{q^{(i)}} \equiv \left( q^{(i)}_1,\ldots,q^{(i)}_K \right) \sim \text{Dir}(\alpha,\ldots,\alpha) \text{;} \label{eq2}\\
&\mathbf{p_{kl\cdot}} \equiv \left( p_{kl1},\ldots,p_{klJ_l} \right) \sim \text{Dir}(\lambda_1,\ldots,\lambda_{J_l}) \label{eq3}\text{,}
\end{align}
where $J_l$ is the number of distinct alleles observed at locus $l$, and these alleles are labeled as $1$ through $J_l$. It is also assumed that $\lambda_1= \cdots = \lambda_{J_l}=1$ to have a uniform distribution on the allele frequencies. A Dirichlet distribution with the unit vector as parameter corresponds to the uniform distribution on the simplex, so that $\mathbf{p_{kl\cdot}} \sim \mathcal{U}(\Delta^{J_l-1})$, where $\Delta^{J_l-1}$ is the unit simplex of $\mathbb{R}^{J_l}$ (see \cite{Frigyik2010IntroductionTT}).

We move from this framework and add uncertainty about the distributions of $\mathbf{q^{(i)}}$ and $\mathbf{p_{kl\cdot}}$; this is interesting since it allows the model to be more flexible. We write $\mathbf{q^{(i)}} \sim \phi$ and $\mathbf{p_{kl\cdot}} \sim \varphi$, where $\phi$ and $\varphi$ are unknown distributions. Instead of using a random measure to express our uncertainty, we consider a finite admixture of distributions; we do so because mixtures of distributions are a widely used tool to approximate any given distribution, and both the Dirichlet distribution and finite admixture models take values on the unit simplex. The number of components of the mixture is unknown.

Let then $\phi \approx \sum_{m=1}^M \alpha_m f_m$ and $\varphi \approx \sum_{s=1}^S \alpha_s g_s$, $M,S \in \mathbb{N}$. Now, as pointed out in section 2.1, there always exist two collections of random points, $\{ X_k \}_{k=1}^n$ and $\{ Y_s \}_{s=1}^j$, $X_1,\ldots,X_n \sim \mathcal{U}(\Delta^{K-1})$ iid and $Y_1,\ldots,Y_j \sim \mathcal{U}(\Delta^{J_l-1})$ iid, for $K$ as in (\ref{eq2}) and $J_l$ as in (\ref{eq3}), such that an element of $\text{Conv}(X_1,\ldots,X_n)$ represents the finite admixture distribution that approximates $\phi$, and an element of $\text{Conv}(Y_1,\ldots,Y_j)$ represents the finite admixture distribution that approximates $\varphi$. We have $K \leq M \leq n$, and $J_l \leq S \leq k$.

We can use a Dirichlet process mixture model, as described in \cite{li}, to estimate the values of $M$ and $S$. Our main result, Theorem 3.3, states that $\mathbb{E}M$ grows at rate $(\log n)^{K-1}$ and $\mathbb{E}S$ grows at rate $(\log j)^{J_l-1}$.

We believe this extension to the model in \cite{Pritchard945} can be useful to the field of genomics: not only do we take into account admixed individuals, but also uncertainty in their proportion and in the probability of sets of allele frequencies. The classical tradeoff between model flexibility and computational efficiency applies to this case too: our extension allows the researcher to be more flexible, at the cost of a higher computation time. In particular, the algorithm the researcher has to follow to implement our generalized model is the following:
\begin{itemize}
\item First estimate $M$ and $S$ using a Dirichlet process mixture model, as described in \cite{li};
\item Then implement the model in \cite{Pritchard945}, with the following modification. Equations (\ref{eq2}) and (\ref{eq3}) become, respectively:
\begin{align*}
\mathbf{q^{(i)}} \equiv \left( q^{(i)}_1,\ldots,q^{(i)}_K \right) &\dot\sim \sum_{m=1}^M \alpha_m f_m; \tag{\ref{eq2}'}\\
\mathbf{p_{kl\cdot}} \equiv \left( p_{kl1},\ldots,p_{klJ_l} \right) &\dot\sim \sum_{s=1}^S \alpha_s g_s \tag{\ref{eq3}'} {,}
\end{align*}
where $\dot\sim$ denotes ``approximately distributed as'', and the $f_m$'s and the $g_s$'s are known distributions with the same form as in (\ref{eq5}).
\end{itemize}
\fi

\section{A procedure to estimate the richest cheap model}\label{algorithm}

%In this section we use the idea of a mixture model based on the extremal set to provide a procedure which finds the richest cheap admixture model. We apply the procedure to a document-term matrix \cite{harman} and examine the number of topics inferred and the word frequency distribution of the topics.

In this section, we present a sparse-factor-analysis-inspired algorithm that estimates the parameter of the richest cheap model introduced in section \ref{intro}, and we apply it to the well-know TREC-1 document-term matrix dataset \cite{harman}. 

Consider an admixture model $X_i \sim \text{Mult}(\pi_i)$, $\pi_i = \sum_{\ell=1}^L \phi_{i,\ell} f_\ell$, $i\in\{1,\ldots,n\}$, where the $X_i$'s are independent. There are two sets of parameters: 
\begin{enumerate}
\item the mixing weights for each individual, that can be arranged in an $n\times L$ matrix $\Phi$ whose entries $\phi_{i,\ell}$ represent the probability that the
$i$-th sample is drawn from the $\ell$-th component. In turn, each row of $\Phi$ is the mixture vector of the $i$-th observation $\phi_i = (\phi_{i,1},\ldots,\phi_{i,L})$;
\item the probability vectors parameterizing each admixture component, which we can write as an $L\times J$ matrix $F$ whose $\ell$-th row is $f_\ell$.
\end{enumerate}
The relation between admixture modeling and sparse factor analysis has been  
explored in detail in \cite{engel}. There, conditions are provided when sparse factor analysis and LDA have very similar results, and the implications for population genetics are discussed. The key insight in \cite{engel} is that given an $n\times J$ observation matrix $X$ (whose $i$-th row is $X_i$) from a multinomial admixture model, learning an admixture model amounts to the following minimization procedure
\begin{equation}\label{min_probl}
\min_{F,\Phi}\| \mathbb{E}[X] -  \Phi F\|^2.
\end{equation}
The sparse factor analysis framework can be summarized as minimizing \eqref{min_probl} with the constraint that many of the elements of $\Phi$ will be zero, or that every observation is a sparse combination of each component. The spirit behind the algorithm proposed in this section is to think of sparsity as the extremal set: we want to find a set of components that are extremal yet still accurately solves the above minimization. We first state the likelihood for the admixture model, assuming a maximum of $L$ components,
$$\mathcal{L}(X_1,...,X_n; \{\phi_1,...,\phi_n\}, \{f_1,...,f_L\}) = \prod_{i=1}^n \mbox{Mult}\left(\pi_i =  \sum_{\ell=1}^L \phi_{i,\ell} f_\ell  \right).$$
A maximum likelihood estimator (MLE) for for the above model is
\begin{equation}
\label{mixtureml}
 \{\{\hat{\phi}_1,...,\hat{\phi}_n\}, \{\hat{f}_1,...,\hat{f}_L\}\} \in \argmax_{\{\phi_1,...,\phi_n\}, \{f_1,...,f_L\}}  \mathcal{L}(X_1,...,X_n; \{\phi_1,...,\phi_n\}, \{f_1,...,f_L\}).
 \end{equation}
A notion of sparsity related to the sparse factor analysis framework is to maximize the likelihood subject to the constraint that components are identifiable, that is, no component can be represented as a convex combination of other components. We consider a procedure that maximizes the following objective function
\begin{equation} \label{opt_probl}
\begin{aligned}
& \underset{ \mathcal{I},\{\phi_1,...,\phi_n\}, \{f_k\}_{k \in \mathcal{I}}} {\text{argmax} }
& &   \prod_{i=1}^n \mbox{Mult}\left(\pi_i =  \sum_{k \in \mathcal{I}} \phi_{i,k} f_k  \right). \\
& \text{subject to}
& & f_k \not \in \mbox{Conv}(f_{\mathcal{I} \setminus \{k\}}) , \; \forall k \in \mathcal{I} ,
\end{aligned}
\end{equation}
where $\mathcal{I}$ is a subset of the set $\{1,...,L\}$, and it is the collection of the indices of the extremal set. Constraint  $f_k \not \in \mbox{Conv}(f_{\mathcal{I} \setminus \{k\}})$, for all $k \in \mathcal{I}$, ensures that no admixture component is contained in the convex combination of the others. Notice that the cardinality of $\mathcal{I}$ represents the number of components $M$ of the richest cheap model in \eqref{rcm}.

%where the probability vector $\phi_i \equiv (\phi_{i,1},\ldots,\phi_{i,M})$ assigns the probability of the $i$-th observation coming from the $\ell$-th mixture component with multinomial parameters
%$f_\ell=(f_{\ell,1},\ldots,f_{\ell,J})$. In this setting the maximum likelihood formulation of learning 
%This section draws inspiration from \cite{engel}. There, the authors consider a new method for
%analyzing population structure, sparse factor analysis (sparse factor analysis), which lies in this same model class as admixture-based methods and principal component analysis (PCA). They show that sparse factor analysis produces similar results to admixture-based models when the data conform to discrete and admixed populations, and gives results similar to PCA when allele frequencies vary continuously with geography. Therefore, sparse factor analysis seems to unify admixture-based methods and PCA when the analyst is interested in studying  structure of natural populations. 
%What is especially important in light of the present work is that the authors see an admixture model as an optimization problem \cite[Equations 1 and 2]{engel}. We use the same idea. Given our frequency vector of interest $\pi_i$, we want to find the model such that $\pi_i=\sum_{\ell=1}^M \varphi_{i,\ell} f_\ell$ having the largest amount of identifiable components, that is, such that the cardinality of $\{f_1,\ldots,f_M\}$ is the largest among all the collections of identifiable components that represent $\pi_i$. The other constraints are that the $\varphi_{i,\ell}$'s have to be positive and sum up to $1$, and that the elements of all the $f_\ell$'s have to be positive and sum up to $1$.
%Let us write this optimization problem in mathematical notation. We denote by $\mathcal{L}(\theta; x_i) = \sum_{\ell=1}^M \varphi_{i,\ell} f_\ell(x_i)$ the likelihood function for $X_i$, where the prameters of interest are $\theta=(f_1,\ldots,f_M \text{; } \varphi_{i,1},\ldots,\varphi_{i,M})$. We also denote by $\mathfrak{F}=\{\mathfrak{f}\}$ the collection of collections $\mathfrak{f}=\{f_j\}_{j=1}^k$ of frequency vectors whose mixture gives $\pi_i$. We are particularly interested in the subcollection $\mathfrak{F}_0 \subset \mathfrak{F}$ of collections $\mathfrak{f}_0=\{f_{0,j}\}_{j=1}^k$ of vectors that cannot be written as convex combination of one another. We then denote by $\mathfrak{f}^\star \in \mathfrak{F}_0$ the collection having the highest possible number of components, that is, 
%$$\#\mathfrak{f}^\star=\max_{\mathfrak{f}_0 \in \mathfrak{F}_0} \# \mathfrak{f}_0,$$
%where $\#$ denotes the cardinality operator.
%We want $\theta$ that maximizes $\mathcal{L}(\theta; x_i)$ subject to some constraints:
%\begin{align}\label{opt_probl}
%\begin{split}
%\hat{\theta}&=\argmax_\theta \mathcal{L}(\theta; x_i)\\
%\text{s.t. } & \{f_1,\ldots,f_M\} = \mathfrak{f}^\star;\\
%& \sum_{j=1}^J f_{\ell,j}=1 \text{, } \forall \ell \in \{1,\ldots,M\} \text{; } f_{\ell,j} \geq 0 \text{, } \forall \ell \in \{1,\ldots,M\} \text{, } \forall j \in \{1,\ldots,J\};\\
%& \sum_{\ell=1}^M \varphi_{i,\ell}=1 \text{; } \varphi_{i,\ell} \geq 0 \text{, } \forall \ell \in \{1,\ldots,M\}.
%\end{split}
%\end{align}

The maximization specified by equation \eqref{opt_probl} is non-convex and finding the global optima is difficult; we propose a two-step procedure to solve it. 


\begin{algorithm}
\caption{Approximating the parameters of the richest cheap model}\label{alg:2step}
\begin{algorithmic}
%\item \hspace*{\algorithmicindent} 
\State \textbf{Step 0} Initialize $t=0$ and set the initial number of components $L_0 \in\mathbb{N}$
\Do:
\State \textbf{Step 1} Compute an MLE as in \eqref{mixtureml} so to obtain the estimated parameters $\{\{\hat{\phi}_1,...,\hat{\phi}_n\}, \{\hat{f}_1,...,\hat{f}_{L_{t}}\}\}$

\State \textbf{Step 2} $L_{t+1}:=\#\text{ex}(\text{Conv}(\{\hat{f}_1,...,\hat{f}_{L_{t}}\})$
\doWhile{$L_{t+1} < L_{t}$} \Comment{\% Call $L_T$ the number of components that exits the loop \%}
\State \Return parameters $\{\{\hat{\phi}_1,...,\hat{\phi}_n\}, \{\hat{f}_1,...,\hat{f}_{L_T}\}\}$
\end{algorithmic}
\end{algorithm}

\iffalse
\item \hspace*{\algorithmicindent} \textbf{Step 1} Compute the MLE specified in \eqref{mixtureml} so to obtain the estimated parameters
$\{\{\hat{\phi}_1,...,\hat{\phi}_n\}, \{\hat{f}_1,...,\hat{f}_{L_0}\}\}$
\item \hspace*{\algorithmicindent} \textbf{Step 2} Compute $\text{Conv}(\{\hat{f}_1,...,\hat{f}_{L_0}\})$ and consider the cardinality $M$ of its extremal set, $M:=\#\text{ex}(\mbox{Conv}(\{\hat{f}_1,...,\hat{f}_{L_0}\})$
\If{$M<{L_0}$}
    \State rerun Step 1 with $M$ components
    \Else \text{ } stop and keep the current parameters
\EndIf
\fi



The parameters returned by Algorithm \ref{alg:2step} are estimates of the parameters of the richest cheap model. Notice that computing the convex hull is evocative of the Choquet procedure described in section \ref{choq}. If $L_T=J$ (where $J$ is the dimension of the Euclidean space we work with), we have an estimate $\hat{\nu}_{\pi_i}$ of the Choquet measure for $\pi_i$, since the geometric representation of the richest cheap model is a simplex. 
%,  since $\hat{\nu}_{\pi_i}(\hat{f}_\ell)=\hat{\varphi}_{i,\ell}$, for all $\ell \in \{1,\ldots,L_J\}$, where we denote the estimates of the weights of the richest cheap model as $\hat{\varphi}_{i,\ell}$, for all $\ell$. 
We have the following important result.

\begin{theorem}\label{algo_cons}
    Call $M^\star$ the true number of components of the richest cheap finite admixture model. Then, if $L_0 \geq M^\star$ and the MLE in \eqref{mixtureml} is consistent,\footnote{Conditions for the MLE to be consistent can be found in \cite{wasserman_2}.} we have that $L_T$ is a consistent estimator for $M^\star$.
\end{theorem}

We leave for future work to verify whether the following plausible conjecture holds. Suppose that our model is ``sparse enough'' as in \eqref{opt_probl}. That is, suppose we constrain the admixture elements to be identifiable. Then, if the MLE in \eqref{mixtureml} is consistent, it should be verified that the collection $\{\hat{f}_1,...,\hat{f}_{L_T}\}$ of identifiable components estimated by Algorithm \ref{alg:2step} is a consistent estimator for the collection  $\{{f}^\star_1,...,{f}^\star_{M^\star}\}$ of true identifiable admixture components.

\begin{remark}\label{rem_l0}
    One thing left to discuss before applying our algorithm to a dataset is how to choose $L_0$. The objective function \eqref{mixtureml} is highly non-concave and, in general, optimization algorithms
for finding the optimal solutions to \eqref{mixtureml} can have sub-linear convergence
rates. This is due to the fact that function \eqref{mixtureml} is locally weakly concave around
the optimal solutions (see \cite{dwivedi} for a more detailed discussion). Hence, $L_0$ cannot just be chosen to be arbitrarily large, as the performance of the optimization algorithms can be strongly
affected. At the same time, it cannot be too small, otherwise our algorithm would not be able to capture the underlying complexity associated with the data at hand, and we would also risk not to meet the $L_0 \geq M^\star$ condition of Theorem \ref{algo_cons}. We select $L_0$ via an ``educated guess'' coming from the exploratory data analysis part of our study or from previous results on the same or similar datasets. In the future we will study a more formal way of coming up with a value for $L_0$. 

\end{remark}








%This To solve optimization problem in \eqref{opt_probl}, we propose a $2$-stage expectation maximization (EM) algorithm. In the first stage, we select a deliberately high number $L$ of components. We then run the EM algorithm to estimate admixture weights and admixture components of our model. After that, we compute the convex hull of the estimated admixture components. If the number of extrema of the convex hull is $L$, we stop. If instead the number of extrema is $M<L$, we run another stage of EM algorithm, this time specifying $M$ many components. The results of this latter stage are the estimates we are interested in. Notice that computing the convex hull is evocative of the Choquet procedure we explained in section \ref{choq}. In particular, at the end of this algorithm, we have an estimate $\hat{\nu}_{\pi_i}$ of the Choquet measure for $\pi_i$,  since $\hat{\nu}_{\pi_i}(\hat{f}_\ell)=\hat{\varphi}_{i,\ell}$, for all $\ell$, where we denote the estimates of the components as $\hat{f}_\ell$, and the estimate of the weights as $\hat{\varphi}_{i,\ell}$, for all $\ell$.

We applied our two-step procedure to a well studied dataset \cite{bail} which is a document-term matrix consisting of term frequencies of $10473$ terms in $2246$ documents collected from Associated Press documents \cite{harman}. We used the latent Dirichlet allocation (LDA) function in the R package  \texttt{topicmodels} \cite{topicmodels} to compute the MLE and the convex hull function in the R  package \texttt{geometry} to compute the convex hull. 

The number of topics obtained in previous studies on the same dataset is between $9$ and $12$ \cite{bail,hou}. For this reason, we ran our algorithm three times on the document-term matrix setting $L_0$ equal to $12$, $25$, and $50$. $L_0=12$ seems the proper educated guess, while $L_0=25$ and $L_0=50$ are safety checks; starting with a higher value of $L_0$ can cause our algorithm to incur problems, as discussed in Remark \ref{rem_l0}.


%; it seemed pointless to start with even higher $L_0$'s.

Computing the convex hull over the full topic frequency vectors -- elements belonging to simplex $\Delta^{2245}$ -- is prohibitive and also does not make sense when the number of topics are less than $2245$. We used principal components analysis (PCA) to project the frequency vectors of the topics onto a lower dimensional space and then computed the convex hull of the projections. We used a simple scree plot 
to notice that $3-5$ dimensions are sufficient to capture about $40\%$ of the variation when we carry out our analysis specifying $50$ initial topics. If we choose $L_0<50$, we have that $3-5$ dimension explain more than $40\%$ of the variation. We only need to compute the number of extremal elements of the convex hull and not the extremal elements themselves in our procedure, so it suffices to compute the convex hull in the low dimensional space.

Given the results in the PCA step, we projected down to $5$ dimensions. The number of extremal elements -- i.e. the number of topics -- we obtained were  $8$, $8$, and $9$ having initialized $L_0$ to $12$, $25$, and $50$, respectively. The number of topics of the richest cheap model seems to be $8$, that is, a mixture of $8$ multinomials appears to be the model that captures the complexity in our dataset using the smallest number of components. The estimated topics $\ell\in\{1,\ldots,8\}$, together with the $10$ terms $i$ having the highest estimated probability $\hat{\varphi}_{i,\ell}$ of being generated from topic $\ell$, are reported in Figure \ref{fig_algo}. We do not report the estimated topic frequency vectors $\hat{f}_\ell$, $\ell\in\{1,\ldots,8\}$, because of their high dimension, $J=2246$. Recall that for all $\ell\in\{1,\ldots,8\}$, we can write $\hat{f}_\ell$ as $\hat{f}_\ell=(\hat{f}_{\ell,1},\ldots,\hat{f}_{\ell,J})^\top$, where $\hat{f}_{\ell,j}$ represents the estimated probability of topic $\ell$ being featured in document $j$.  

Let us add a brief discussion. In this section we provide a simple application of Algorithm \ref{alg:2step} to show that the latter is actually implementable. A more in-depth discussion on how it performs in comparison to other methods is postponed to future research. We conjecture that Algorithm \ref{alg:2step} finding a smaller number of topics, compared e.g. to \cite{bail,hou}, is due to the fact that these latter look for the ``most interpretable'' solution, that is, the one whose topics are better interpretable by a human user. On the contrary, we explicitly look for the model involving the smallest possible number of admixture components, which is able to fully capture the complexity of the data at hand. That is, we look for the richest cheap model. Another reason why we find a smaller number of topics may linked to the fact that we adopt PCA. Once again, all of these questions are left for future research.




\begin{figure}[h!]
\centering
\includegraphics[width=\textwidth]{plot}
\caption{Parameters estimated by our algorithm. For every topic, we show the $10$ terms having the highest probability of being generated from that topic. As we can see, Topics 2 and 7 are hard to interpret, a common issue with topic models \cite{fung}.}
\label{fig_algo}
\centering
\end{figure}







%We do not report the estimated parameters as and it seems to be in line with previous studies \cite{bail,hou}. 
%The reason why the number of extremal points seems to decrease as $L$ increases will be studied in future work.

%so it has been thoroughly studied in other works, see e.g. \cite{bail}. The majority of these studies use around 10 topics (9 to 12), as by increasing the number of topics it seems that clarity is sacrificed \cite{hou}. 
%We next computed the convex hull of the 50 topic frequency vectors using function \texttt{convhulln} in the \texttt{geometry} package. This function allows to compute the convex hull of vectors in high dimensional spaces. Using \texttt{convhulln}, though, is not as immediate as it may seem. Indeed, every time it is required to compute a convex hull of frequency vectors, this function builds the simplex in which the vectors lie. This means that \texttt{convhulln} cannot compute the convex hull of a number $d$ of vectors that are more than $d+1$ dimensional. We then projected the topic frequency vectors onto lower dimensional spaces, of dimension $2$, $5$, $10$, and $20$, and computed the convex hull of the projections. We noticed that while the mode of the number of extrema of the convex hulls is roughly the same when we consider the projections onto $2$ and $5$-dimensional spaces ($9$ and $10$, respectively), it is different and higher when we consider the projections onto $10$ and $20$-dimensional spaces ($22$ and $31$, respectively). We used the mode to see the number of extrema in the majority of the projection spaces having the same dimension.


%To illustrate the algorithm, we apply it to the \href{https://rdrr.io/cran/topicmodels/man/associatedpress.html}{Associated Press data}, a document-term matrix which contains the term frequency of 10473 terms in 2246 documents. It is a well-known dataset in topic modeling, see e.g. \cite{bail}. We analyzed the dataset using \texttt{R}. To find the estimates in \eqref{mixtureml}, we ran a latent Dirichlet allocation (LDA) analysis using the \texttt{LDA} function in the \texttt{topicmodels} package. LDA models a document $X_i$ as mixtures of topics, whose number has to be specified ex ante. Each topic is represented by a frequency vector $f_\ell$. Each entry $f_{\ell,j}$ of this latter is the probability that one of the words in the analysis is used within a given topic. LDA gives estimates $\hat{f}_\ell$ of the frequency vectors associated with the topics, and also of the document-specific weights $\hat{\phi}_{i}$. In our analysis, we first specified $L=50$ topics, so we obtained the estimates for 50 10473-dimensional topic frequency vectors, and for 2246 50-dimensional mixture frequency vectors, each giving the weights for the 50 topics in every document. The first 7 rows and 3 columns of both these estimated matrices are shown in Tables \ref{data} and \ref{data2}, respectively.
%\begin{table}[ht]
%\centering
%\begin{tabular}{rrrr}
 % \hline
% & Topic1 & Topic2 & Topic3\\
 % \hline
%aaron & 1.18265879980859e-150 & 1.18221871851056e-115 & 0.000102968156941872\\
%abandon & 5.48461105804886e-48 & 0.000136362896867619 & 5.34318904496137e-152\\
%abandoned & 0.000268240337778773 & 5.73649091875299e-57 & 2.18202776749332e-05\\
%abandoning & 6.27823716013166e-152 & 1.26846233563446e-151 & 1.85914414311264e-07\\
%abbott & 1.48477951219737e-116 & 2.14293175560996e-151 & 2.38989688999521e-152\\
%abboud & 1.35351225683029e-153 & 2.73465190100513e-153 & 1.67501088973385e-189\\
%abc & 1.99351411989207e-71 & 4.78746229300328e-151 & 1.89037703917877e-07\\
  % \hline
%\end{tabular}
%\caption{First 7 rows and 3 columns of the estimated topics matrix $\hat{F}$. Every column $\hat{f}_\ell$ is a topic frequency vector, and every row represents one of the words in the dataset. Entry $\hat{F}_{i,j}$ is the probability that word $i$ is used in topic $j$. We use notation $F$ from equation \eqref{min_probl}.}
%\label{data}
%\end{table}
%\begin{table}
%\centering
%\begin{tabular}{rrrr}
 % \hline
%& Topic1 & Topic2 & Topic3\\ 
 % \hline
%Document1 & 4.60919586013466e-05 & 4.60919586013569e-05 & 0.0816654070448194\\
%Document2 & 4.47343537629079e-05 & 4.47343537629223e-05 & 0.535202552055116\\
%Document3 & 4.75345440316063e-05 & 4.75345440316166e-05 & 4.7534544031604e-05\\
%Document4 & 5.85247925104753e-05 & 5.85247925104352e-05 & 5.85247925105158e-05\\
%Document5 & 0.000182414663610274 & 0.000182414663610255 & 0.000182414663610258\\
%Document6 & 2.88187144157083e-05 & 2.88187144157225e-05 & 2.88187144157281e-05\\
%Document7 & 7.24918976235351e-05 & 7.24918976235246e-05 & 7.24918976235294e-05\\
  % \hline
%\end{tabular}
%\caption{First 7 rows and 3 columns of the estimated weights matrix $\hat{\Phi}$. Every row $\hat{\phi}_{i}$ is a mixture frequency vector, and every column represents one of the topics. Entry $\hat{\Phi}_{i,j}$ is the probability that topic $j$ is dealt with in document $i$. We use notation $\Phi$ from equation \eqref{min_probl}.}
%\label{data2}
%\end{table}


%The difference in the number of extrema is again due to the way \texttt{convhulln} works. Since it has to build the simplex, the higher the dimensionality of the space we are considering, the higher the number of vectors the function ``needs", and in turn the higher the number of extrema of the convex hull. In the limit case, if we project onto a $61$-dimensional space, \texttt{convhulln} tells us that the number of extrema is $60$. Because of this, we conjecture that the true number of extrema of our convex hull is somewhat similar to the number of extrema of its projection onto small dimensional subspaces. Empirical evidence supporting this intuition are the following. First, we have that the highest number of topics used in a single document is 11. To verify this, we looked for documents assigning weight $\hat{\phi}_{i,\ell}$ greater than $0.01$ to at least 12 topics, founding none. Only two assign probability greater than $0.01$ to 11 topics, one of which is depicted in Figure \ref{fig1}. The other empirical evidence we have is that this dataset is a classic one in topic modeling, and so it has been thoroughly studied in other works, see e.g. \cite{bail}. The majority of these studies use around 10 topics (9 to 12), as by increasing the number of topics it seems that clarity is sacrificed \cite{hou}. Finally, we checked for the robustness of this empirical evidence: starting our analysis with 100, 150, and 200 topics yielded basically the same results.

%\begin{figure}[h!]
%\centering
%\includegraphics[width=.8\textwidth]{fig1}
%\caption{This figure represents document 4 in our analysis, which assigns probability $\hat{\phi}_{i,\ell}$ greater than $0.01$ to 11 topics out of 50.}
%\label{fig1}
%\centering
%\end{figure}
%The above (empirical) argument led us to choose 10 as the number of extrema of the convex hull. We then re-run the LDA analysis specifying 10 topics, and produced bar graphs that describe the top terms for each topic. These are shown in Figure \ref{fig3} in Appendix \ref{groups}.

\iffalse
In Table \ref{data}, we plot the first $10$ rows and $10$ columns of our dataset. The rows correspond to words, the columns to track IDs, and the numbers are the number of time a given word is used in a given track.

\begin{table}[h!]
\centering
\begin{tabular}{rrrrrrrrrrr}
  \hline
 & 1 & 2 & 3 & 4 & 5 & 6 & 7 & 8 & 9 & 10 \\ 
  \hline
you & 6 & 10 & 28 & 5 & 4 & 16 & 1 & 4 & 39 & 6 \\ 
  to & 4 & 17 & 15 & 4 & 5 & 4 & 1 & 11 & 30 & 9 \\ 
  and & 2 & 8 & 2 & 3 & 7 & 1 & 1 & 2 & 10 & 3 \\ 
  a & 2 & 2 & 12 & 2 & 2 & 3 & 2 & 7 & 10 & 4 \\ 
  me & 5 & 2 & 22 & 1 & 4 & 5 & 1 & 3 & 28 & 4 \\ 
  it & 3 & 1 & 2 & 11 & 1 & 5 & 1 & 5 & 21 & 3 \\ 
  not & 1 & 3 & 2 & 4 & 1 & 3 & 1 & 1 & 1 & 5 \\ 
  in & 1 & 2 & 4 & 9 & 9 & 4 & 1 & 3 & 20 & 3 \\ 
  my & 1 & 3 & 2 & 3 & 1 & 3 & 1 & 6 & 11 & 2 \\ 
  is & 2 & 4 & 1 & 2 & 2 & 2 & 1 & 6 & 12 & 1 \\ 
   \hline
\end{tabular}
\caption{First $10$ rows and $10$ columns of the musiXmatch dataset.}
\label{data}
\end{table}


In our analysis, we kept all the lyrics, but we only used $100,000$ tracks to save computation time. The results are compelling nonetheless. We analyzed the dataset using \texttt{R}. For every word, we created $3$ multinomial bins using the function \texttt{makemultdata}. Each bin represents the level of usage -- low, medium, and high -- of the word. So, for example, we have that the frequency vector associated with the word ``you'' is $(10445, 8949, 80605)$. This means that the word ``you'' was used in $10445 + 8949 + 80605=99999$ tracks, and out of all them, $10445$ times its level of usage was low, $8949$ times it was medium, and $80605$ times it was high. To visualize the output of function \texttt{makemultdata}, we plot its first $10$ rows in Table \ref{table_mult}.


\begin{table}[ht]
\centering
\begin{tabular}{rrrr}
  \hline
 & \text{low usage} & \text{medium usage} & \text{high usage} \\ 
  \hline
you & 10445 & 8949 & 80605 \\ 
  to & 12363 & 10678 & 76939 \\ 
  and & 14239 & 12140 & 73536 \\ 
  a & 15882 & 14636 & 69294 \\ 
  me & 18753 & 15383 & 65505 \\ 
  it & 23094 & 17528 & 58857 \\ 
  not & 25198 & 18558 & 55550 \\ 
  in & 24980 & 18396 & 55777 \\ 
  my & 26814 & 18975 & 53172 \\ 
  is & 29082 & 19973 & 49740 \\ 
   \hline
\end{tabular}
\caption{First $10$ rows of the output of function \texttt{makemultdata}.}
\label{table_mult}
\end{table}




Then, we ran the EM algorithm using the function \texttt{multmixEM}, having selected $L=25$ admixture components. Given the results of this first EM step, it was clear that the number of elements was some $M<L$, since some of the estimated components were equal to one another. In Figure \ref{fig:algo} we first plot the estimated components as points in the unit simplex in $\mathbb{R}^3$. Then, we compute the convex hull of these points and plot the projection of this convex hull to $2$-dimensional spaces in order to enhance legibility.
\begin{figure}[h!]\label{fig_algo}
  \begin{subfigure}[t]{.4\textwidth}
    \centering
    \includegraphics[width=\linewidth]{algo1}
  \end{subfigure}
  \hfill
  \begin{subfigure}[t]{.4\textwidth}
    \centering
    \includegraphics[width=\linewidth]{algo2}
  \end{subfigure}

  \medskip

  \begin{subfigure}[t]{.4\textwidth}
    \centering
    \includegraphics[width=\linewidth]{algo3}
  \end{subfigure}
  \hfill
  \begin{subfigure}[t]{.4\textwidth}
    \centering
    \includegraphics[width=\linewidth]{algo4}
  \end{subfigure}
\caption{Top left: Graphical representation of the 25 admixture components estimated by the EM algorithm. Remaining figures: Projection of their convex hull to $2$-dimensional spaces.}\label{fig:algo}
\end{figure}

The number of extremal elements of the convex hull is $9$. So we re-ran the EM algorithm, this time with $M=9$ components. The results were the following: the estimated admixture weights $\hat{\varphi}_{i,\ell}$ are the entries of vector $\hat{\varphi}_i=(0.01877830, 0.59164685, 0.18697996, 0.01051828, 0.04569413$, $0.03003144, 0.06341927, 0.03740000, 0.01553177)^\top$, and the estimated admixture components are the rows of Table \ref{comp}.



\begin{table}[ht]
\centering
\begin{tabular}{rrrr}
  \hline
 & \text{low usage} & \text{medium usage} & \text{high usage} \\ 
  \hline
$\hat{f}_1$ & 0.3957634 & 0.2164343 & 0.38780228 \\ 
  $\hat{f}_2$ & 0.8192425 & 0.1050865 & 0.07567096 \\ 
  $\hat{f}_3$ & 0.7581605 & 0.1315361 & 0.11030344 \\ 
  $\hat{f}_4$ & 0.1435515 & 0.1237342 & 0.73271435 \\ 
  $\hat{f}_5$ & 0.6463701 & 0.1790915 & 0.17453834 \\ 
  $\hat{f}_6$ & 0.4900437 & 0.2188221 & 0.29113423 \\ 
  $\hat{f}_7$ & 0.7037633 & 0.1544497 & 0.14178702 \\ 
  $\hat{f}_8$ & 0.5732760 & 0.2016811 & 0.22504289 \\ 
  $\hat{f}_9$ & 0.2762114 & 0.1941330 & 0.52965565 \\ 
   \hline
\end{tabular}
\caption{Estimated admixture components.}
\label{comp}
\end{table}



\begin{table}[ht]
\centering
\begin{tabular}{rrrr}
  \hline
 & \text{low usage} & \text{medium usage} & \text{high usage} \\ 
  \hline
$\hat{f}_1$ & 0.3957634 & 0.2164343 & 0.38780228 \\
$\hat{f}_2$ & 0.8192425 & 0.1050865 & 0.07567096 \\
$\hat{f}_3$ & 0.7581605 & 0.1315361 & 0.11030344 \\
$\hat{f}_4$ & 0.1435515 & 0.1237342 & 0.73271435 \\
$\hat{f}_5$ & 0.6463701 & 0.1790915 & 0.17453834 \\
$\hat{f}_6$ & 0.4900437 & 0.2188221 & 0.29113423 \\
$\hat{f}_7$ & 0.7037633 & 0.1544497 & 0.14178702 \\
$\hat{f}_8$ & 0.5732760 & 0.2016811 & 0.22504289 \\
$\hat{f}_9$ & 0.2762114 & 0.1941330 & 0.52965565 
   \hline
\end{tabular}
\caption{Estimated admixture components.}
\label{table_comp}
\end{table}


$$ \begin{matrix}
 & \text{low usage} & \text{medium usage} & \text{high usage} \\
\hat{f}_1 & 0.3957634 & 0.2164343 & 0.38780228 \\
\hat{f}_2 & 0.8192425 & 0.1050865 & 0.07567096 \\
\hat{f}_3 & 0.7581605 & 0.1315361 & 0.11030344 \\
\hat{f}_4 & 0.1435515 & 0.1237342 & 0.73271435 \\
\hat{f}_5 & 0.6463701 & 0.1790915 & 0.17453834 \\
\hat{f}_6 & 0.4900437 & 0.2188221 & 0.29113423 \\
\hat{f}_7 & 0.7037633 & 0.1544497 & 0.14178702 \\
\hat{f}_8 & 0.5732760 & 0.2016811 & 0.22504289 \\
\hat{f}_9 & 0.2762114 & 0.1941330 & 0.52965565 
\end{matrix}  $$


The estimated admixture components can be interpreted as group of words. That is, there are subsections of the $5,000$ words that we analyzed whose frequency vectors are close in the Euclidean metric to $\hat{f}_\ell$, for some $\ell \in \{1,\ldots,9\}$. We prefer the expression ``group of words'' to the word ``topic'', more widely used in the topic modeling literature. This because the words associated with some of the groups cohere into something that could be called a latent theme or topic, but there are other groups that are not as coherent. This phenomenon is rather common (as explained in \cite{bail}), an makes post-hoc interpretation of topic models rather dangerous. To avoid this, we interpret the estimated admixture components as group of words. The nine groups are provided in Appendix \ref{groups} in a word cloud format.








We now put Theorems \ref{polya} and \ref{thm1} to use. In \cite[Pages 1 and 2]{engel}, the authors show that an admixture model can be seen as an optimization problem. Finding the cheapest finite admixture model in the context of this paper too can be seen as such:
\begin{align}\label{opt_probl}
\begin{split}
\min_{M \in \mathbb{N}} \text{ } &\pi_i=\sum_{\ell=1}^M \varphi_{i,\ell} f_\ell\\
\text{s.t. } &\forall i, \sum_{\ell=1}^M \varphi_{i,\ell}=1 \text{; } \forall i,\forall \ell, \varphi_{i,\ell} \geq 0\\
&\forall \ell, \sum_{j=1}^J f_{\ell,j}=1 \text{; } \forall j,\forall \ell, f_{\ell,j} \geq 0.
\end{split}
\end{align}
In \cite{engel}, the authors use sparse factor analysis to analyze population structure. In particular, they use a version of the expectation maximization (EM) algorithm to perform parameter estimation of their model, which enables the application of sparse factor analysis to genome-wide data. We too propose a variation of the EM algorithm to solve \eqref{opt_probl}. Fix any initial number $L$ of components for our mixture model, then proceed as follows:
\begin{enumerate}
\item Compute E and M steps to estimate admixture components and admixture weights;
\item When the algorithm converges, compute the convex hull of the components, and keep only the extremal elements of the convex hull: these are the components of the cheapest model;
\item If the extremal elements of the convex hull correspond to the admixture components estimated by the EM algorithm, 
\begin{itemize}
\item the analysis ends.
\end{itemize}
If instead the extremal elements of the convex hull are a subset of the admixture components estimated by the EM steps
\begin{itemize}
\item use the Pólya approach outlined in Theorems \ref{polya} and \ref{thm1} to retrieve the weights of these fewer components.
\end{itemize}
\end{enumerate}
In an applied perspective, we should choose the highest $L$ that makes sense in the context of the model. If $L$ is too large, step (2) of the algorithm corrects for possible overspecifications.
\fi






%We state the growth rate of number of components for a finite admixture model both in expectation and in distribution. The tools we use in our analysis combine the concept of a Choquet measure with classic results in stochastic geometry on the number of extrema of random polytopes. We show that if our admixture
%weights are probability vectors from a unit $(J-1)$-simplex then the number of admixture components grows as $(\log n)^{J-1}$; in the standard mixture case we recover the $\log n$ rate. We also state a central limit
%theorem for the number of mixture components. In addition, we state the convergence of the sequence of the empirical measures generated by our model to the Choquet measure. Lastly, we relate our model to
%a classical non-parametric density estimator based on a P\'olya tree.


%We first established how an $m$-dimensional admixture model on $\mathbb{R}^J$, for some $J \in \mathbb{N}$, can be represented by an element of the convex hull of $n$ points drawn from a uniform on the $(J -1)$-dimensional simplex, $J \leq m \leq n$. We used this result to show that, if $m$ is unknown, the growth rate of $\mathbb{E}m$ is $(\log n)^{J-1}$.
%We then inspected the theoretical properties of our model. We saw how, using the Choquet Theorem, a unique measure on the extrema -- that we call the Choquet measure -- represents any point inside of the convex hull. 
%Then, we studied how to approximate the distribution of the sequence of points $\{X_1,\ldots,X_n \}$ we drew; we used Aldous' version of de Finetti's theorem, together with a result from Diaconis and Freedman in \cite{diaconis1980}.
%We then turned our attention to the behavior of the extrema of the convex hull. We showed that our convex hull, which is a convex polytope, converges to a smooth convex body as the number of its extreme points approaches infinity; we also saw the growth rate of the expected value of the number of extrema, and we provided a Central Limit Theorem for the number of extrema.
%After this, we showed that the sequence of empirical probability measures on the extreme points converges strongly to the Choquet measure.
%Finally, we proved that the posterior of an appropriately defined P\'olya tree is weakly consistent for the Choquet measure; the P\'olya tree posterior converges to the Choquet measure at rate $({\log n}/{n})^{{\alpha}/(2 \alpha + 1)}$, for a properly specified $\alpha$.

%To conclude our work, we pointed out how our results carry over to a non-parametric analysis and to finite mixture models. In particular, in finite mixture models the rate of growth for the expected number of components is $\log n$, a result that matches the rate of growth for the {actual} number of components found in \cite{miller}.
%To conclude our work, we provided an application of our model to the field of genomics: we generalized the hierarchical model in \cite{Pritchard945} by using admixtures to approximate unknown distributions; we then used our main result to establish the growth rate of the expected number of components of these admixtures. 
%In this work, we approached finite admixture models from an innovative perspective: we used many concepts developed in other literatures, making them relevant in the study of finite admixture models. In addition, we provided two novel results, namely that a convex hull generated by $n$ iid uniform random variables within a convex polytope tends to a smooth convex body as the number of  its extrema goes to infinity, and that the sequence of empirical probability measures on the extrema of our convex hull converges strongly to the Choquet measure. 






%\section{Groups of words resulting from our analysis}\label{groups}
%The bar graphs that describe the top terms for each of the 10 topics resulting from the analysis in section \ref{algorithm} are given in Figure \ref{fig3}.
%\begin{figure}[h!]
%\centering
%\includegraphics[width=\textwidth]{LDA_rich_cheap}
%\caption{Bar graphs that describe the top terms for each of the 10 topics.}
%\label{fig3}
%\centering
%\end{figure}




%\bibliographystyle{plain}


%\bibliographystyle{plain}
%\bibliography{Paper_fmm}  



\section{Conclusion}\label{concl}
In this paper there are two key ideas. The first one is that we can use techniques from stochastic convex geometry on the growth rate of the expected number of extremal elements of random polytopes to provide insights into the asymptotic growth rate of the expected number of identifiable admixture components. We prove that the expected number of identifiable components grows at rate $(\log n)^{J-1}$ where $J$ is the dimension of the Euclidean space we work with.
%, and $n$ is the amount of data points we collect.
%For the generic mixture model case, we recover a $\log n$ growth rate.  
We also show that the number of identifiable admixture components concentrates around its expected value, and we provide a central limit theorem for its distribution. 
%of the number of identifiable components. 
The other key concept is that we can retrieve the identifiable admixture weights using techniques from Choquet theory. In particular, we show that if the convex hull $\mathcal{K}_M$ generated by the identifiable admixture components is a simplex, we can apply Glivenko-Cantelli to recover the identifiable admixture weights. We also give an algorithm to estimate the richest cheap finite admixture model. 

%Choquet measure for $\pi_i$, for any $\pi_i \in \mathcal{K}_M$. In turn, this gives us the

An interesting open question is whether there are other instances in (Bayesian) inference where coupling results from stochastic (convex) geometry with results from Choquet theory allows to develop novel analyses, insights, models, or algorithms. For example, studying the properties of an apeirogon -- a polytope with infinitely many sides -- could give us insights on infinite mixture models. Another research direction for the future regards a formal procedure to initialize $L_0$ in Algorithm \ref{alg:2step}. A promising way to tackle this issue is to use the backward induction approach of \cite[Section 1.2]{peskir}, where the authors find an optimal stopping time for a given expected utility maximization problem. In our framework, such optimal stopping time could be interpreted as number $L_0\in\mathbb{N}$ that strikes a balance between being not too large, so to avoid the problems highlighted by \cite{dwivedi}, and not too small, so to satisfy the assumption of Theorem \ref{algo_cons}. We also plan to further generalize the results in section \ref{from-to}.
%, an expected utility maximization framework seems to be the correct method.



\section*{Acknowledgments} 
The authors would like to thank Andrea Aveni, Jordan Brian, Pierpaolo De Blasi, Federico Ferrari, Jürgen Jost, Jeremias Knoblauch, Rostislav Matveev, Vittorio Orlandi, Sonia Petrone, Mai Zhou, Alessandro Zito, and one anonymous reviewer for their helpful comments.  Michele Caprio would like to acknowledge funding from the following grants: NIH 1R01MH118927-01, ARO MURI W911NF2010080. Sayan Mukherjee would like to acknowledge funding from NSF DEB-1840223, NIH R01 DK116187-01, HFSP RGP0051/2017, NSF DMS 17-13012, and NSF CCF-1934964.


\appendix


\section{Further results}\label{app1}
\subsection{Joint distribution of admixture components}
\label{infiniteex}

In this section, we assume that the number $L$ of admixture components in \eqref{fam_first} is known, but the components $\{f_1,\ldots,f_L\}$ are not. We assume they are identically distributed random vectors, but we do not require independence. After realizing that collection $\{f_1,\ldots,f_L\}$ can be seen as a finite exchangeble sequence, we inspect how to approximate its joint distribution applying de Finetti's theorem and a result by Diaconis and Freedman \cite{diaconis1980}. 
%to the finite exchangeable sequence $(f_1,\ldots,f_L)$.

Following \cite{aldous_2010}, we can state de Finetti's result from a functional analytic viewpoint in our framework as follows. Let $\mathcal{S} \equiv \Delta^{J-1} \subset \mathbb{R}^{J}$, and recall that a sequence of random variables $X_i$'s is exchangeable if
$$(X_i)_{i \geq 1} \overset{\text{d}}{=} (X_{\text{perm}(i)})_{i \geq 1} \text{,}$$
for any finite permutation $\text{perm}$, where $\overset{\text{d}}{=}$ denotes equality in distribution. We can assume that the elements ${f}_1,\ldots,$ ${f}_{L}$ form a {finite} exchangeable sequence because the order in which they appear provides no additional information about the finite admixture model.

Let $\mathscr{P}(\mathcal{S})\equiv \mathscr{P}(\mathcal{S},\mathcal{B}(\mathcal{S}))$ be the set of probability measures on $(\mathcal{S},\mathcal{B}(\mathcal{S}))$, where $\mathcal{B}(\mathcal{S})$ is the Borel sigma-algebra for $\mathcal{S}$. Let then $\mathscr{P}( \mathscr{P}(\mathcal{S}) )$ be the set of probability measures on $\mathscr{P}(\mathcal{S})$. When we define an infinite exchangeable sequence of $\mathcal{S}$-valued random variables, we are actually defining an exchangeable measure $\Theta\in\mathscr{P}(\mathcal{S}^\infty)$, where $\Theta$ is the distribution of the sequence. Consider the set $\mathfrak{M}:=\{ \mu^\infty:=\mu \times \mu \times \cdots \text{ s.t. } \mu \in \mathscr{P}(\mathcal{S}) \} \subset \mathscr{P}(\mathcal{S}^\infty)$, that is the set of extremal elements of the convex set $\mathfrak{C}$ of exchangeable elements of $\mathscr{P}(\mathcal{S}^\infty)$. Then, we have
\begin{align}\label{defin_funct}
\Theta(A)=\int_{\mathscr{P}(\mathcal{S})} \mu^\infty(A) \text{ } \Lambda(\text{d}\mu), \quad \forall A \subset \mathcal{S}^\infty,
\end{align}
where $\mu^\infty \in \mathfrak{M}$ and we call $\Lambda$ the (unique) \textit{de Finetti measure}. Hence, there is a bijection between $\Lambda \in \mathscr{P} ( \mathscr{P}(\mathcal{S}) )$ and $\Theta \in \mathscr{P}(\mathcal{S}^\infty)$. 

Notice how (functional analytic) de Finetti's theorem \eqref{defin_funct} is similar to Theorem \ref{choq2}. There exists a unique measure ($\Lambda$ for de Finetti and $\nu_c$ for Choquet) supported on the extremal elements ($\mathfrak{M}$ for de Finetti and $E$ for Choquet) of a convex set ($\mathfrak{C}$ for de Finetti and $C$ for Choquet) that allows to represent any element ($\Theta$ for de Finetti and affine function $f(c)$, $c\in C$, for Choquet) within that set.







%Notice that a consequence of Proposition \ref{simplex_conseq} is that for all $x \in \Delta^{J-1}$, there exists a unique Choquet measure $\tilde{\nu}_x$, whose support are the extrema of $\Delta^{J-1}$, denoted as $\text{ex}(\Delta^{J-1})$, that allows to represent $x$. Then, there always exists a probability measure $\breve{\nu}_x$ supported on the whole simplex $\Delta^{J-1}$, whose restriction to $\text{ex}(\Delta^{J-1})$ is given by $\tilde{\nu}_x$. To this extent, $\breve{\nu}_x \times \breve{\nu}_x \times \cdots =:\breve{\nu}_x^\infty$ belongs to $\mathfrak{M}$.




%$\breve{\nu}:\mathcal{B}(\Delta^{J-1}) \rightarrow [0,1]$ such that $\breve{\nu}(B)>0$ if and only if $B=ex\Delta^{J-1}$; that is, $\breve{\nu}$ is an extension of $\tilde{\nu}$ to the Borel sets of $\Delta^{J-1}$, which gives nonzero measure to $ex\Delta^{J-1}$ only. 

As we pointed out before, we can assume $\mathbf{f}:=({f}_1,\ldots,$ ${f}_{L})$ to be a {finite} exchangeable sequence. We assumed that the admixture components are identically distributed but not necessarily independent, so let $f_1,\ldots,f_L \sim \mu$. Suppose, without loss of generality, that $\mathbf{f}$ is part of a much longer sequence of $m$ components
$$\left( {f}_1,\ldots,{f}_{L},\ldots,{f}_{m} \right).$$ 
Then, we can use \cite[Theorem 13]{diaconis1980} to compute an approximation of $\Theta_L$, the distribution of $\mathbf{f}$. Let us denote by $\Theta_m$ the distribution of $\left( {f}_1,\ldots,{f}_{L},\ldots,{f}_{m} \right)$; it is an exchangeable probability on $\mathcal{S}^m$. Then, $\Theta_L$, $L \leq m$, is the projection of $\Theta_m$ onto $\mathcal{S}^L$. Define the value $\beta(m,L)$ as
$$\beta(m,L):=1-\frac{m^{-L}m!}{(m-L)!},$$
and notice that $\beta(m,L) \leq \frac{1}{2} \frac{L(L-1)}{m}$. 




The theorem states that there exists $\tilde{\Lambda} \in \mathscr{P}(\mathscr{P}(\mathcal{S}))$ such that the probability $\Theta_{\mu L}$ defined on $\mathcal{S}^L$ as
$$\Theta_{\mu L}(A)=\int_{\mathscr{P}(\mathcal{S})} \mu^{L}(A) \text{ } \tilde{\Lambda}(\text{d}\mu) , \quad \forall A \subset \mathcal{S}^L$$
is such that
$d_{TV}(\Theta_L,\Theta_{\mu L}) \leq \beta(m,L)$, for all $L \leq m$. 
We denoted by $\mu^L$ the distribution of $L$ independent picks from $\mu$, that is, $\mu^L((y_1,\ldots,y_L))=\prod_{j=1}^L \mu(y_j)$, and by $d_{TV}$ the total variation distance $$d_{TV}(\Theta_L,\Theta_{\mu L}):=\sup\limits_{A \subset S^L} \left| \Theta_L(A)-\Theta_{\mu L}(A) \right| \text{.}$$
%We have that $\beta(m,L)$ is such that
%$$1-\beta(m,L)=\frac{m^{-L}m!}{(m-L)!} \quad \text{and} \quad \beta(m,L) \leq \frac{1}{2} \frac{L(L-1)}{m} \text{.}$$

Notice that $\tilde{\Lambda}$ depends on $m$ and $\Theta_m$, but not on $L$, and its analytical form is given in \cite[Proof of Theorem 13]{diaconis1980}.
\subsection{Number of extremal elements of the convex hull having the least amount of vertices}\label{app_2}

The following is an interesting result dealing with the number of extremal elements of a convex hull in $\Delta^{J-1}$ -- but not in any smaller-dimensional unit simplex -- having the least amount of vertices.

\begin{proposition}\label{rema2}
\iffalse
The number of random points we draw, $n$, is related to $f_0({K_n})$. In particular, for any $n \in \mathbb{N}$, $n=k f_0(K_n)$, where $k \geq 1$ is a positive constant.

Also,
\fi
Call $\mathscr{K}\subset \Delta^{J-1}$, $J\in\mathbb{N}$, a polytope such that
\begin{align}\label{ex_within}
\begin{split}
    \tilde{e}:=\#\text{ex}(\mathscr{K}) = \min_{n\in\mathbb{N}} \text{ } &n\\
    \text{subject to} \text{ } &\nexists q\in\{2,\ldots,J\} : \mathscr{K} \subset \Delta^{J-q}
\end{split}
\end{align}
%$\mathscr{K} \subset \Delta^{J-1}$ and $\mathscr{K} \not\subset \Delta^{J-q}$, for $J\in\mathbb{N}$ and $q\in\{2,\ldots,J\}$. 
Then, $\tilde{e}=J$.
%Let $\tilde{e}$ be the number of extrema of the convex hull (polytope) in our unit simplex with the least amount of vertices. Then, $\tilde{e}=J$. if we let $\tilde{e}:=\#\text{ex}(\mathscr{K})$, we have that
\end{proposition}

\iffalse
\subsection{Choquet simplex}\label{app_choquet}
The following is the general definition of a Choquet simplex.
\begin{definition}\label{simplex_def}
A nonempty convex set $C$ (not necessarily compact) of a locally convex space $V$ is a \textit{Choquet simplex} if it has the following property. Under the embedding of $V$ as the hyperplane $V \times \{1\}$ in the space $V \times \mathbb{R}$, the projecting cone 
$$\tilde{C}:=\left\lbrace{\alpha c \in V \times \mathbb{R} : c \in C \subset V \times \{1\} \text{, } \alpha \geq 0 }\right\rbrace$$
of $C$ transforms the space $V \times \mathbb{R}$ into a partially ordered space $P$ such that the space of differences $\tilde{C} - \tilde{C}$ generated by $P$ is a vector lattice in the order induced by $C$. That is, each pair $c_1,c_2 \in \tilde{C} - \tilde{C}$ has at least upper bound $c_1 \vee c_2 \in \tilde{C} - \tilde{C}$.
\end{definition} 
\fi 





 
\iffalse
Let us start with the former. Consider measurable space $(E_M,\mathcal{B}(E_M)=2^{E_M})$, call $\mathcal{M}_{E_M} \equiv \Delta(E_M,\mathcal{B}(E_M))$, and equip this latter with sigma-algebra $\mathcal{B}(\mathcal{M}_{E_M})$. A random measure $\mathcal{P}$ is a function on a generic probability space $(\Omega_1,\mathcal{F}_1,{P}_1)$,
$$\mathcal{P}:(\Omega_1,\mathcal{F}_1,{P}_1) \rightarrow (\mathcal{M}_{E_M},\mathcal{B}(\mathcal{M}_{E_M})),$$
so $\mathcal{P}(\omega)\in \mathcal{M}_{E_M}$, for all $\omega\in\Omega_1$. Let now $\mathcal{P}_k \sim DP \left(\alpha P_0 + \sum_{j=1}^k \delta_{e_j} \right)$ and $\mathcal{P}_\delta \sim \delta_{\nu_{\pi_i}}$. Then, equation \eqref{polya_eq} means that for all $\varepsilon>0$ and all continuous and bounded functionals $g$ on $E_M$,\footnote{The space of continuous and bounded functionals $g$ on $E_M$ is usually denoted by $C_b(E_M)$.}
\begin{align*}
    P_1\bigg( \bigg\{ &\omega\in\Omega_1 : \bigg| \sum_{e\in E_M} g(e) \mathcal{P}_k(\omega)(e) - \sum_{e\in E_M} g(e) \mathcal{P}_\delta(\omega)(e) \bigg| >\varepsilon \bigg\} \bigg)\\
    &=P_1\bigg( \bigg\{ \omega\in\Omega_1 : \bigg| \sum_{e\in E_M} g(e) \mathcal{P}_k(\omega)(e) - \sum_{e\in E_M} g(e) \nu_{\pi_i}(e) \bigg| >\varepsilon \bigg\} \bigg)\xrightarrow[k\rightarrow\infty]{}0.
\end{align*}

Let us then turn our attention to equation \eqref{gk_cvg}. In the alternative procedure explained in Remark \ref{rem_proc}, we have that 
\fi

\iffalse 
\subsection{Clarification of equation \eqref{gk_cvg}}\label{clar}
In this section, we elucidate the meaning of ``almost surely'' in equation \eqref{gk_cvg}. In Theorem \ref{thm1}, the labels are treated as random variables, so $\ell$ is regarded as a function on a generic probability space $(\Omega,\mathcal{F},\mathbb{P})$,
$$\ell:(\Omega,\mathcal{F},\mathbb{P}) \rightarrow (\{1,\ldots,M\},2^{\{1,\ldots,M\}}),$$
hence $\ell(\omega)\in\{1,\ldots,M\}$, for all $\omega\in\Omega$. Let now $\ell_1,\ldots,\ell_k\sim\zeta$ iid; then, equation \eqref{gk_cvg} coupled with \cite[Chapter 1, Remarks 1 and 2]{zhou} means that for all $\varepsilon>0$,
$$\mathbb{P}\left( \left\lbrace{\omega\in\Omega : \sup_{x\in\mathbb{R}} \left| \frac{1}{k}\sum_{j=1}^k \mathbb{I}_{[\ell_j(\omega),\infty)}(x)-F_\zeta(x) \right|>\varepsilon}\right\rbrace\right) \leq \frac{8k+4}{\exp(\frac{\varepsilon^2}{8}k)} \xrightarrow[k\rightarrow\infty]{}0.$$
\fi 





\section{Proofs}\label{proofs}
\iffalse
\begin{proof}[Proof of Theorem \ref{exp_pois}]
Let us denote by $\mathbb{K}^2_+(d)$ the set of compact convex sets in $\mathbb{R}^d$, $d \geq 2$, having nonempty interior, boundary of differentiability class $\mathcal{C}^2$, and positive Gaussian curvature. Pick any $K \in \mathbb{K}^2_+$, and let $K_n$ denote the convex hull of $n\in\mathbb{N}$ random points sampled independently according to the uniform distribution in $K$. 
%$\Pi_n$ denote the convex hull of the intersection of $K$ with a Poisson point process with intensity $n$.
Then, in \cite[page 485]{reitzner1}, the author shows that
$$\mathbb{E}[F_0(K_n) = \tau_d \Omega(K)n^{\frac{d-1}{d+1}}\left(1+o(1)\right)$$
as $n$ grows to infinity. Since $\Delta^{J-1} \in \mathbb{K}^2_+(J)$, and because we assumed $J\geq 2$, then \eqref{exp_pois_eq} follows.
\end{proof}

\begin{proof}[Proof of Corollary \ref{exp_pois_growth}]
Immediate from Theorem \ref{exp_pois}.
\end{proof}



\begin{proof}[Proof of Theorem \ref{conc_ineq_pp}]
    In \cite[Theorem 1.1]{grote}, the authors show that if we consider a stationary Poisson point process in $\mathbb{R}^d$ with intensity $\lambda>0$, let $\eta_\lambda$ denote its restriction to the unit ball $\mathbb{B}^d$, and $\Pi_\lambda$ be the convex hull of the points of $\eta_\lambda$, then if $y\geq 0$ and $\beta(d):=4+\frac{2}{d+1}$, we have that 
    \begin{align*}
    P\bigg( \left| F_0(\Pi_\lambda)-\mathbb{E}[F_0(\Pi_\lambda)] \right| \geq &y \sqrt{\mathbb{V}[F_0(\Pi_\lambda)]} \bigg)\\ &\leq 2 \exp \left( -\frac{1}{4} \min\left\lbrace{\frac{y^2}{2^{\beta(d)}},c_2(d) \lambda^{\frac{d-1}{2\beta(d)(d+1)}}y^{\frac{1}{\beta(d)}}}\right\rbrace \right),
\end{align*}
    where $c_1(d),c_2(d)\in(0,\infty)$ are constants depending (only) on $d$. This immediately gives \eqref{conc_ineq_pp_eq}. In \cite[Section 5]{grote}, the authors point out that constants $c_1(d),c_2(d)$ are finite and that in the proof may change from line to line, so there is no direct insight into how these constants depend on the dimension $d$ of the Eculidean space.
\end{proof}





\begin{proof}[Proof of Theorem \ref{clt_pois}]
    Consider again the set $\mathbb{K}^2_+(d)$ of compact convex sets in $\mathbb{R}^d$, $d \geq 2$, having nonempty interior, boundary of differentiability class $\mathcal{C}^2$, and positive Gaussian curvature. Pick any $K \in \mathbb{K}^2_+$, and let $\Pi_n$ denote the convex hull of the intersection of $K$ with a Poisson point process with intensity $n$. Then, in \cite[Proof of Theorem 2, pages 501--502]{reitzner1}, the author shows that
    $$\left| P\left( F_0(\Pi_n)\leq \mathbb{E}[F_0(\Pi_n)]) + x \sqrt{\mathbb{V}[F_0(\Pi_n)]}\right) - \Phi(x)\right| = \mathcal{O}\left( n^{-\frac{1}{2}+\frac{1}{J+1}}\left(\log n\right)^{2+\frac{2}{J+1}} \right).$$
    Since $\Delta^{J-1}\in\mathbb{K}^2_+(J)$, and because we assumed $J\geq 2$, then \eqref{clt_pois_eq} follows.
\end{proof}
\fi


Recall that a chain $\mathcal{F}_0(P)\subset \mathcal{F}_1(P) \subset \cdots \subset \mathcal{F}_{h}(P)$ of $i$-dimensional faces of a polytope $P$ is called a \textit{tower} (or a \textit{flag}) of $P$, and we denote it as $\tau(P)$.

\begin{proof}[Proof of Theorem \ref{growth}]
 In \cite[Theorem 6]{reitzner} and \cite[Theorem 5]{barany2}, the authors show that, given a convex polytope $P$ in $\mathbb{R}^d$, if we call $P_n$ the convex hull of $n$ points sampled iid from a uniform on $P$, then 
$$\mathbb{E} \left[F_0(P_n)\right] = \frac{\tau(P)}{(d+1)^{d-1}  (d-1)!}  (\log n)^{d-1} + \mathcal{O} \left( (\log n)^{(d-2)} \log\log n \right).$$
Then, since $\Delta^{J-1}$ is a convex polytope in $\mathbb{R}^J$, and given the way we defined $K_n$, equations \eqref{first_theorem} and \eqref{equation_imp_first} follow immediately after realizing that $\tau(\Delta^{J-1})=(J-1)! \cdot J$.
\end{proof}





\begin{proof}[Proof of Theorem \ref{var_thm}]


\iffalse
We first show the following claim.

\begin{claim}\label{claim1}
    Let $\text{Vol}$ denote the volume operator. Then, for all $J\in\mathbb{N}$, $\text{Vol}(\Delta^{J-1})=\frac{1}{J!}$.
\end{claim} 

\begin{proof}
    Notice that the unit simplex $\Delta^{J-1}$ is the set of elements $(x_1,\ldots,x_J)$ such that $x_j\geq 0$, for all $j\in\{1,\ldots,J\}$, and $\sum_{j=1}^J x_j \leq 1$. Foliating the simplex with ``horizontal'' hyperplanes $x_J=z$, $z\in[0,1]$, and applying Fubini's theorem we obtain
    $$\text{Vol}(\Delta^{J-1})=\int_0^1 (1-z)^{J-1} \text{Vol}(\Delta^{J-2}) \text{d}z=\frac{1}{n} \text{Vol}(\Delta^{J-2}),$$
    which gives us a recursion, concluding the proof.
\end{proof}
\fi
Let $\text{Vol}$ denote the volume operator. In \cite[Theorem 1.3]{reitzner_var}, the authors show that, given a convex polytope $P$ in $\mathbb{R}^d$ such that $\text{Vol}(P)=1$, if we call $P_n$ the convex hull of $n$ points sampled iid from a uniform on $P$, then there exists a constant $C>0$ such that
$$C \tau(P)(\log n)^{d-1} < \mathbb{V}[F_0(P_n)] < C \tau(P)^3(\log n)^{d-1}.$$
Recall that $\text{Vol}(\Delta^{J-1})=\sqrt{J}/[(J-1)!]$. Then, since $\text{Vol}(\Delta^{J-1})$ is in a fixed relation 
$$x \mapsto f(x)=\frac{\sqrt{J}}{(J-1)!} \cdot x$$ 
with $\text{Vol}(P)$, because $\Delta^{J-1}$ is a convex polytope in $\mathbb{R}^J$, and given the way we defined $K_n$, equation \eqref{var_eq} follows immediately.
\end{proof}




\begin{proof}[Proof of Theorem \ref{clt_2}]
In \cite[Corollary 1.2]{pardon}, the author shows that, given a convex polytope $P$ in $\mathbb{R}^2$ of unit area, if we call $P_n$ the convex hull of $n$ points sampled iid from a uniform on $P$, then
$$\lim_{n\rightarrow\infty} \sup_{x\in\mathbb{R}} \left| \mathbb{P}\left( \frac{F_0(P_n)-\mathbb{E}[F_0(P_n)]}{\sqrt{\mathbb{V}[F_0(P_n)}} \leq x \right) - \Phi(x) \right|=0.$$
Recall that the area of the unit simplex in $\mathbb{R}^3$ is $\sqrt{3}/2$. Then, since the area of $\Delta^2$ is in a fixed relation $x\mapsto f(x)=\frac{\sqrt{3}}{2}x$ with the area of $P$, because $\Delta^2$ is a convex polytope in $\mathbb{R}^2$, and given the way we defined $K_n$, equation \eqref{clt_2_eq} follows immediately.
\end{proof}



\begin{proof}[Proof of Theorem \ref{clt_gen}]
In \cite[Theorems 2, 6]{reitzner1} the author shows that, given a smooth compact convex set $P$ in $\mathbb{R}^d$, if we call $P_n$ the convex hull of $n$ points sampled iid from a uniform on $P$, then
$$\left| \mathbb{P} \left( F_0(P_n)\leq \mathbb{E}[F_0(P_n)]+x\sqrt{\mathbb{V}[F_0(P_n)]} \right) - \Phi(x) \right|= \mathcal{O}\left( n^{-\frac{1}{2(J+1)}}\left(\log n\right)^{2+\frac{2}{J+1}} \right).$$
Since $\hat{\Delta}^{J-1}_\epsilon$ is a smooth compact convex set in $\mathbb{R}^J$, and given the way we defined $\hat{K}_n$, equation \eqref{clt_gen_eq} follows immediately.
\end{proof}



\begin{proof}[Proof of Theorem \ref{conc_unif_thm}]
    In \cite[Theorem 2.11, Section 7]{vu}, the author shows that given a smooth compact convex set $P$ in $\mathbb{R}^d$, if we call $P_n$ the convex hull of $n$ points sampled iid from a uniform on $P$, then there exist positive constants $c,\Xi,\epsilon_0$ such that for any $\epsilon\in(0,\epsilon_0]$, $V\geq \Xi n^{\frac{d-1}{d+1}}$, $C\geq n\epsilon$, and $\lambda\in(0,\frac{V}{4C^2})$, the following holds
        $$\mathbb{P}\left( \left| F_0(P_n) -\mathbb{E}[F_0(P_n)] \right| \geq \sqrt{\lambda V} \right) \leq 2 \exp(-\lambda/4)+\exp(-c\epsilon n) + \exp\left(-cn^\frac{d-1}{3d+5}\right).$$
        Since $\hat{\Delta}^{J-1}_\epsilon$ is a smooth compact convex set in $\mathbb{R}^J$, and given the way we defined $\hat{K}_n$, equation \eqref{conc_unif_eq} follows immediately.
\end{proof}





















\iffalse
\begin{proof}[Proof of Theorem \ref{growthdist}]
Let us denote by $\mathbb{K}^2_+$ the set of compact convex sets in $\mathbb{R}^d$, $d \geq 2$, having nonempty interior, boundary of differentiability class $\mathcal{C}^2$, and positive Gaussian curvature. Pick any $K \in \mathbb{K}^2_+$, and sample $n$ points iid from the uniform on $K$. Call their convex hull $P_n$. Then, in \cite[Theorem 6]{reitzner1}, the author shows that  there are numbers $d_n$ bounded between two positive constants depending on $K$, and a constant $c(K)$, such that 
\begin{equation}\label{thm_2_mid}
\left| \mathbb{P} \left( \frac{F_i(P_n)-\mathbb{E}[F_i(P_n)]}{\sqrt{d_n n^{1-\frac{2}{d+1}}}} \leq {t} \right) - \Phi({t}) \right| \leq {c}(K) n^{-\frac{1}{2(d+1)}} (\log n)^{2+3i+\frac{2}{d+1}}.
\end{equation}
The denominator $\sqrt{d_n n^{1-\frac{2}{d+1}}}$ is of the same asymptotic order as the standard deviation of $F_i(P_n)$, so the inequality in \eqref{thm_2_mid} implies a central limit theorem for $F_i(P_n)$.

Notice then that $\Delta^{J-1} \in \mathbb{K}^2_+$ for any $\mathbb{R}^{J}$, $J \geq 2$. Hence, given the way we defined $K_n$, equation \eqref{second_theorem} follows immediately. As we can see, the rate of convergence of the distribution of $F_0(K_n)$ to $\Phi$ is given by $n^{-\frac{1}{2(J+1)}}(\log n)^{2+\frac{2}{J+1}}$.
\end{proof}
\fi

\iffalse
\begin{proof}[Proof of Theorem \ref{hullconv}]
%We prove the convergence for $K_n$ only, as the proof is the same for $\hat{K}_n$. 
Call $\hat{E}_n$ the extremal set of $\hat{K}_n$, that is, $\hat{E}_n= \text{ex}(\hat{K}_n)=F_0(\hat{K}_n)$.
%, where we assumed without loss of generality that the first $M(n)$ mixture components are the identifiable ones. 
Let $F_0(\hat{K}_n) \rightarrow \infty$, and call $\hat{E} \neq \emptyset$ the set that $\hat{E}_n$ tends to in the Hausdorff distance
\begin{align*}
d_H(\hat{E}_n,\hat{{E}})&:=\max\left\lbrace{ \sup_{{f} \in \hat{E}_n} \inf_{{g} \in \hat{{E}}}d_2 \left( {f},{g} \right) \text{ , } \sup_{{g} \in \hat{{E}}} \inf_{{f} \in \hat{E}_n} d_2 \left( {f},{g} \right) }\right\rbrace\\
&=\sup_{{s} \in \hat{E}_n \cup \hat{{E}}} \left| \inf_{{f} \in \hat{E}_n} d_2 \left( {s},{f} \right) - \inf_{{g} \in \hat{{E}}} d_2 \left( {s},{g} \right) \right| \text{,}
\end{align*}
as the cardinality of $\hat{E}_n$  approaches infinity. Here $d_2$ denotes the usual Euclidean distance, and the second equality is an equivalent way of writing the Hausdorff distance. Let $\hat{{K}}$ be the convex hull of $\hat{{E}}$.

\iffalse
Normally, the extrema of a convex set are defined to be the points that cannot be written as convex combinations of other points. Here we consider an equivalent definition: the extrema of $K_n$ are the points in  ${E}_n=\bigcap_k {E}_k$, where ${E}_k=\{ {f}_{e_1},\ldots,{f}_{e_k} \} \subset \{ {f}_1,\ldots,{f}_{M(n)} \}$ such that for all $k$, $\text{Conv}(\tilde{f}_{e_1},\ldots,\tilde{f}_{e_k})=\text{Conv}(\tilde{f}_1,\ldots,\tilde{f}_{M(n)})=K_n$.
\fi

{Step 1}: We first show that $\hat{{K}}$ is well defined. By construction, we know that $\hat{{E}} \neq \emptyset$; $\hat{{K}}$ is then the convex hull of the points in $\hat{{E}}$, which is well defined as we can always construct the convex hull of any given (sub)set of a vector space.

{Step 2}: Now, we show that $\hat{K}$ has infinitely many sides. Suppose for the sake of contradiction that $\hat{K}$ has finitely many sides. Then, it has a finite number of $\ell$-faces, for some $\ell$, which implies a finite number of vertices. But $\hat{K}$ is the convex hull of the elements in $\hat{{E}}$, that are infinite, a contradiction.

{Step 3}: $\hat{K}$ is convex: this is immediate from it being the convex hull of $\hat{{E}}$.

{Step 4}: We are left to show that $\hat{K}$ is the limit of $\hat{K}_n$. We have seen that $\hat{E}_n \rightarrow \hat{{E}}$ in the Hausdorff metric as $n$ goes to infinity; we also know that $\hat{K}_n=\text{Conv}(\hat{E}_n)$, for all $n$.  
%(there is a small abuse of notation here: $K_n$ is the convex hull of {the elements} of ${E}_n$; since no confusion arises and since we save some notation, we leave it as it is). 
But then
$$\hat{K}_n = \text{Conv}\left( \hat{E}_n \right) \xrightarrow[n \rightarrow \infty]{d_H} \text{Conv} (\hat{{E}} )=\hat{K} \text{,}$$
which concludes our proof.
\end{proof}
\fi

\iffalse
\begin{proof}[Proof of Theorem \ref{extrematopost}]
The proof consists of showing that if condition $\mathbb{E}[T(n)]=\gamma \cdot \mathbb{E}[M(n)]$ holds, then the growth rate of the extrema stated in Theorem \ref{growth} will hold for a  more general procedure. We already know from Theorem \ref{growth} that if ${S}_1,\ldots,{S}_{n} \sim \mbox{Uniform}(\Delta^{J-1})$ iid, then $\lim_{n \rightarrow \infty} (\log n)^{-(J-1)} \mathbb{E}[F_0(K_n)] = c(J)$, a value depending on the dimension of the Euclidean space $\mathbb{R}^J$ we work in. Recall that the number of extrema $F_0(K_n)$ of the convex body associated with our mixture model in the uniform case corresponds to $M(n)$. We now relax the assumption in equation \eqref{eq1_fmm}. 
%, and the support is a convex body.
%First assume that the support $\mathfrak{S}_n$ of the posterior distribution is a convex body, for all $n$ and all $x_1,\ldots,x_n$. 
Fix any $n \in \mathbb{N}$ 
%and consider observations $x_1,\ldots,x_n$
. Let then $\mathbb{E}[T(n)]$ be the expected number of extrema of $\breve{K}_n$.
%, where $\tilde{f}_1,\ldots,\tilde{f}_{T(n)} \sim \mathbb{P}(f_1,\ldots,f_{T(n)} \mid  x_1,\ldots,x_n)$. 
Assume that $\mathbb{E}[T(n)]=\gamma \cdot \mathbb{E}[M(n)]$, for some $\gamma \in \mathbb{R}$, $\gamma\neq 0$. Then, by Theorem \ref{growth}, we have that  $\lim_{n \rightarrow \infty} (\log n)^{-(J-1)} \mathbb{E}[T(n)] = \gamma \cdot c(J)$.
%, which implies that $$\lim_{n \rightarrow \infty} (\log n)^{-(J-1)} \mathbb{E}[L(n)] = \gamma \cdot b \cdot c(J).$$
Equation \eqref{gen_post} then follows by putting $c^{\prime}(J,\gamma)=\gamma \cdot c(J)$.
\end{proof}
\fi

\begin{proof}[Proof of Theorem \ref{extrema2}]
%First notice that sequence $(\gamma_n)$ exists because we can always write a natural number as a linear function of another natural number, using a rational coefficient. 
By hypothesis, we have that for all $n\geq N$, $\mathbb{E}[T(n)]=\gamma_n \mathbb{E}[M(n)]$.
%and $\lim_{n\rightarrow\infty}\gamma_n/r_{\gamma_n} =\zeta$. 
In addition, 
\iffalse
so 
\begin{equation}\label{algebr_man}
\frac{1}{\gamma_n}\mathbb{E}[T(n)]= \mathbb{E}[M(n)].
\end{equation}
Then, 
\fi
by Theorem \ref{growth} we have that 
$$\lim_{n \rightarrow \infty} \frac{\mathbb{E}[M(n)]}{(\log n)^{J-1}} =c(J).$$
Hence we obtain that
$$\lim_{n\rightarrow\infty} \frac{\mathbb{E}[T(n)]}{{\gamma_n}(\log n)^{J-1}}=\lim_{n\rightarrow\infty} \frac{\gamma_n\mathbb{E}[M(n)]}{{\gamma_n}(\log n)^{J-1}}= \lim_{n\rightarrow\infty} \frac{\mathbb{E}[M(n)]}{(\log n)^{J-1}}= c(J),$$
concluding the proof.
\iffalse
and by substituting according to \eqref{algebr_man}, we obtain 
$$\lim_{n \rightarrow \infty} \frac{1}{\gamma_n(\log n)^{J-1}} \mathbb{E}[T(n)]=c(J),$$
concluding the proof.

the left hand side in \eqref{gen_gr_rate} is equal to
$$\lim_{n \rightarrow \infty} \frac{1}{r_{\gamma_n} (\log n)^{J-1}} \gamma_n \mathbb{E}[M(n)],$$
which in turn is equal to 
$$\underbrace{\lim_{n \rightarrow \infty} \frac{1}{r_{\gamma_n}} \gamma_n}_{=1}\cdot \underbrace{\lim_{n \rightarrow \infty} \frac{1}{(\log n)^{J-1}}  \mathbb{E}[M(n)]}_{=c(J) \text{ by Theorem 2.1}}=c(J).$$
\fi
\end{proof}

\begin{proof}[Proof of Corollary \ref{extrema2_cor}]
Suppose that the assumptions of Theorem \ref{extrema2} hold and that $\gamma_n=\mathcal{O}(\varpi_n)$. This latter means that there exists $R\in\mathbb{R}$ and $N\in\mathbb{N}$ such that for all $n\geq N$,
$$\frac{\gamma_n}{\varpi_n} \leq R.$$
Then,
$$\lim_{n\rightarrow\infty} \frac{\mathbb{E}[T(n)]}{\varpi_n (\log n)^{J-1}}=\lim_{n\rightarrow\infty} \frac{\gamma_n \mathbb{E}[M(n)]}{\varpi_n (\log n)^{J-1}}=\lim_{n\rightarrow\infty} \frac{\gamma_n}{\varpi_n} \lim_{n\rightarrow\infty} \frac{\mathbb{E}[M(n)]}{(\log n)^{J-1}} \leq R c(J),$$
concluding the proof. 

%Immediate from Theorem \ref{extrema2} and the definition of big O of a function, that we denoted by $\mathcal{O}(\cdot)$.
\end{proof}


\begin{proof}[Proof of Proposition \ref{simplex_conseq}]
The proposition is an immediate consequence of 
%Definition \ref{simplex_def} and 
Theorem \ref{choq2}.
\end{proof}



%\begin{proof}[Proof of Theorem \ref{polya}]
%The result follows directly from \cite[Theorem 3.3.5]{ghosh}.
%\end{proof}


\begin{proof}[Proof of Theorem \ref{thm1}]
The uniform almost sure convergence statement is a consequence of the Glivenko-Cantelli's Theorem \cite{vdv_asy}, while the rate of convergence can be easily derived from \cite[Lemma 2.1]{zhou}, via a Borel-Cantelli argument.\footnote{We refer the reader that is further interested in the rates of the Glivenko-Cantelli convergence to the recent work \cite{dolera}.}
\end{proof}
%We can apply \cite[Theorem 2.1$^*$]{zhou} to obtain that, given any $0 < \eta < 0.5$, $\sup_{x \in\mathbb{R}} |F_K(x) - F_\zeta(x)|$ goes to zero faster than $\frac{1}{k^\eta}$. That is, $\sup_{x \in\mathbb{R}} |F_K(x) - F_\zeta(x)| = o(\frac{1}{k^\eta})$.
%, while the rate of convergence comes from  \cite[Chapter 1, Remarks 1 and 2]{zhou}; see also Appendix \ref{clar}.
%The need for $\mathcal{K}_M$ to be a (Choquet) simplex comes from Proposition \ref{simplex_conseq}. The weak almost sure convergence statement comes from \cite[Corollary 4.17]{ghosal_vdv}, while the rate of convergence comes from \cite[Example 8.5]{ghosal_vdv}. In this latter, the authors give it relative to the semimetric $d(P,Q)=|P(A)-Q(A)|$, for a fixed $A$. But $\nu_{\pi_i}$ is defined on $E_M$ that is finite and independent of $k$, so there is only a bounded number of choices for $A$. In turn this implies that the rate of convergence holds in terms of the total variation distance as well.
%because this holds for all measurable sets $A$, then the rate of convergence holds in terms of the total variation distance as well.
%


\begin{proof}[Proof of Theorem \ref{algo_cons}]
    First notice that $L_T$ depends on the amount $n$ of data available to perform the estimating procedure in Algorithm \ref{alg:2step}, so we can write $L_T\equiv L_T(n)$. Suppose now for the sake of contradiction that $\tilde{M}:=\lim_{n\rightarrow\infty} L_T(n) \neq M^\star$. Then, we have two cases, either $\tilde{M} > M^\star$, or $\tilde{M} < M^\star$.

    Case 1: If $\tilde{M} > M^\star$, then there exist a collection $\{\hat{f}_1,\ldots,\hat{f}_{\tilde{M}-M^\star}\}$ of admixture components that can be written as a convex combination of the remaining $M^\star$ components. Together with the assumption of the MLE in \eqref{mixtureml} being consistent, this contradicts Step 2 of Algorithm \ref{alg:2step}.

    Case 2: Suppose now $\tilde{M} < M^\star$. Then, since by Algorithm \ref{alg:2step} $L_T(n) \leq L_0$, for all $n$, we have that $\tilde{M}\leq L_0$. So, if $\tilde{M} < M^\star$, it follows that either $L_0 < M^\star$, or $L_0 \geq M^\star$. If $L_0 \geq M^\star$, Algorithm \ref{alg:2step} would have stopped at $\tilde{M} = M^\star$. But since we are assuming $\tilde{M} < M^\star$, then this means that $L_0 < M^\star$, which contradicts the assumption of our theorem. This concludes the proof.
\end{proof}

\begin{proof}[Proof of Proposition \ref{rema2}]
\iffalse
Notice that $n$ and $f_0(K_n)$ are related by construction, since this latter is the number of extreme points of the convex hull generated by the $n$ random points that we initially draw. Suppose then, for the sake of contradiction, that there exists $n^* \in \mathbb{N}$ such that $n^*=k^*  f_0(K_n)$, where $k^* < 1$. This implies that the number of random points drawn is smaller than the number of extrema of the convex hull they form, a contradiction.
\fi
Suppose for the sake of contradiction that $\tilde{e}\neq J$. This means that either $\tilde{e}>J$, or $\tilde{e}<J$. If the latter holds, then there exists $q^\prime\in\{2,\ldots,J\}$ such that $\mathscr{K} \subset \Delta^{J-q^\prime}$, which contradicts \eqref{ex_within}. If instead $\tilde{e}>J$, then we can find $\mathscr{K}^\prime \subsetneq \mathscr{K}$ such that $\# \text{ex}(\mathscr{K}^\prime)<\tilde{e}$, but $\mathscr{K}^\prime$ is still a proper subset of $\Delta^{J-1}$, thus again contradicting \eqref{ex_within}. This concludes the proof.
%To see that $\tilde{e}=J$ it is enough to notice that in $\Delta^1$, the convex hull (which is a polytope) with the least amount of vertices is a line segment (also called dion), in $\Delta^2$ it is a triangle, in $\Delta^3$ it is a tetrahedron, and in $\Delta^4$ it is a $5$-cell. Then, an induction argument proves the claim.
\end{proof}


\bibliographystyle{plain}
\bibliography{References}
\end{document}




